\documentclass[11pt,a4paper,oneside,openany]{book}

\usepackage[utf8]{inputenc}
\usepackage[T1]{fontenc}
\usepackage[english]{babel}
\usepackage{csquotes}
\usepackage[charter]{mathdesign}
\usepackage{amsmath}
\usepackage{lipsum}
\usepackage{amsthm}
\usepackage{mathtools}
\usepackage{physics}
\usepackage{graphicx}
\usepackage{enumitem}
\usepackage{tikz-cd}
\usepackage{geometry}

% Page header management
\usepackage{fancyhdr}
\pagestyle{fancy}
\fancyhf{} 
\fancyhead[L]{\slshape\nouppercase{\leftmark}} 
\fancyhead[R]{\thepage} 
\renewcommand{\headrulewidth}{0pt}

% Redefinition of the plain style for chapter starting pages
\fancypagestyle{plain}{
	\fancyhf{}
	\renewcommand{\headrulewidth}{0pt}
	\fancyfoot[R]{\thepage} 
}

\usepackage[backend=biber, sorting=none, style=alphabetic]{biblatex}
\addbibresource{bib.bib}
\usepackage[colorlinks=true,linkcolor=black,citecolor=black,urlcolor=black]{hyperref}

\usepackage{epigraph}
\setlength{\epigraphwidth}{0.5\textwidth}
\renewcommand{\epigraphflush}{flushright}
\renewcommand{\textflush}{flushright}

\newtheorem{theorem}{Theorem}[chapter]
\newtheorem{definition}{Definition}[section]
\newtheorem{lemma}[definition]{Lemma}
\newtheorem{proposition}[definition]{Proposition}
\newtheorem{corollary}[definition]{Corollary}
\theoremstyle{remark}
\newtheorem{remark}[definition]{Remark}
\newtheorem{example}[definition]{Example}

\let\mathcal\relax
\newcommand{\mathcal}[1]{\text{\usefont{OMS}{cmsy}{m}{n}#1}}
\DeclareMathOperator{\Spec}{Spec}
\DeclareMathOperator{\Ad}{Ad}
\DeclareMathOperator{\ad}{ad}
\newcommand{\poisson}[2]{\{#1,#2\}}
% ============================================================
% Comandi Matematici
% ============================================================

% --- Insiemi numerici ---
\newcommand{\R}{\mathbb{R}}
\newcommand{\C}{\mathbb{C}}
\newcommand{\N}{\mathbb{N}}
\newcommand{\K}{\mathbb{K}}
\newcommand{\A}{\mathfrak{A}}
\newcommand{\B}{\mathcal{B}}

% --- Spazio di Hilbert e spazio proiettivo ---
% ATTENZIONE: \H è già un comando di LaTeX (accento ungherese, es. \H{o}=ő).
% Il renewcommand lo sovrascrive: nella tesi non ti servirà mai quell'accento,
% quindi è sicuro, ma se preferisci non toccare comandi di sistema usa \Hil
% al posto di \H ovunque sotto.
\renewcommand{\H}{\mathcal{H}}
\newcommand{\PH}{\mathbb{P}\H}                    % P H, spazio proiettivo

% --- Algebre C* ---
\newcommand{\Cstar}{C^*}                          % da usare in math mode: $\Cstar$
\newcommand{\Csalg}{$C^*$-algebra}                % da usare nel testo: "è una \Csalg\ commutativa"
\newcommand{\Csalgs}{$C^*$-algebre}               % plurale
\newcommand{\BH}{\mathcal{B}(\H)}                           % operatori limitati su H
\newcommand{\Cz}[1]{C_0(#1)}                      % C_0(M)
\newcommand{\Cinf}[1]{C^\infty(#1)}               % C^infty(M)
\newcommand{\Cinfc}[1]{C^\infty_c(#1)}            % C^infty_c(M)
\newcommand{\Mat}[2]{M_{#1}(#2)}                  % M_n(C), es. \Mat{2}{\C}

% --- Algebre di Lie (fraktur generico) ---
\newcommand{\mf}[1]{\mathfrak{#1}}                % \mf{g}, \mf{su}(2), \mf{u}(n), \mf{X}(M)

\newcommand{\qcomm}[2]{\dfrac{1}{i\hbar}[#1,#2]}  % (1/ih)[A,B]  -- condizione di Dirac, Qubit IV

% --- Notazione bra-ket ---
\newcommand{\melem}[3]{\langle#1|#2|#3\rangle}    % elemento di matrice <psi|A|psi>
\newcommand{\expect}[2]{\langle#1\rangle_{#2}}    % <A>_psi

% --- Campo vettoriale hamiltoniano ---
\newcommand{\Ham}[1]{X_{#1}}                      % \Ham{f}, \Ham{f_H}, \Ham{f_A}

% --- Simboli e operatori vari ---
\newcommand{\Id}{\mathbb{1}}
\renewcommand{\Re}{\mathrm{Re}}                   % la tesi usa Re/Im in tondo, non ℜ/ℑ
\renewcommand{\Im}{\mathrm{Im}}
\newcommand{\LC}{\varepsilon_{ijk}}               % simbolo di Levi-Civita (indici fissi ijk)
\DeclareMathOperator{\Ran}{Ran}
\DeclareMathOperator{\Ker}{Ker}
\DeclareMathOperator{\Span}{span}
\DeclareMathOperator{\sing}{sing}
\newgeometry{margin=2cm}
\setlength{\headwidth}{\textwidth}

\begin{document}
\setcounter{chapter}{4}
% =====================================================================
% Chapter 4: Recovery of Analytical Expression of Hamiltonians
% DRAFT v1. Every \cite{TODO-...} and every \ref to Chapters 1-3 must be
% matched to the real bibliography keys and labels. The map is at the end.
% Environments assumed: definition, theorem, proposition, lemma, corollary,
% remark, proof (amsthm). No custom macros are used.
% =====================================================================
% =====================================================================
% Chapter 4: Recovery of Analytical Expression of Hamiltonians
% DRAFT v3. Every \cite{TODO-...} and every \ref to Chapters 1-3 must be
% matched to the real bibliography keys and labels. The map is at the end.
% Environments assumed: definition, theorem, proposition, lemma, corollary,
% remark, proof (amsthm). No custom macros are used.
% =====================================================================


\chapter{Recovery of Analytical Expression of Hamiltonians}
\label{ch:analytical-hamiltonians}

The geometric formulation of Chapter~\ref{ch:projective} identifies the dynamics of a quantum system with the Hamiltonian flow of the mean value function $\langle H\rangle$ on the projective Hilbert space. It does not, however, say anything about the operator $H$ itself. Proposition~\ref{prop:kahler-isomorphisms} shows that every continuous one-parameter group of Kähler isomorphisms is generated by a selfadjoint operator, unique up to an additive constant, and nothing in the Kähler structure of $\mathbb{P}\mathcal{H}$ distinguishes the generator of a free particle from the generator of an interacting one. The analytical expression of a Hamiltonian is therefore not a consequence of the geometry alone and requires an additional physical input.

In this chapter the additional input is a symmetry principle. We require that a group of spacetime transformations acts on $\mathbb{P}\mathcal{H}$ by Kähler isomorphisms, with the time translations producing the dynamics. The multiplication law of the group turns into algebraic relations between the generators, and for the Galilei group these relations determine the Hamiltonian of an elementary system up to an additive constant. This reverses the logic of canonical quantization, where the kinetic term is postulated by analogy with the classical energy. Here the form $|\mathbf{p}|^2/2m$ is a consequence of the covariance of the state space under boosts, while the mass $m$ enters as a cohomological invariant of the projective action, so that it is not assumed at the outset.

Section~\ref{sec:galilei} introduces the Galilei group together with the central extension that its projective representations require. Section~\ref{sec:invariance} formulates Galilean invariance of $\mathbb{P}\mathcal{H}$ and lifts it to a unitary representation. Section~\ref{sec:galilean-functions} translates the generators into Hamiltonian functions on $\mathbb{P}\mathcal{H}$, extending the formulas of Chapter~\ref{ch:projective} to unbounded operators. Section~\ref{sec:free-hamiltonian} proves the main result and draws its geometric consequences. The relativistic case is left to a later chapter, where the same strategy will be applied to the Poincaré group.

\section{Cohomology and Cocycles}\label{sec:cohomology}

Several statements of Chapters 1 to 3 hold locally and fail globally, and the failure always has the same structure. In Remark~1.5.11 a momentum map cannot be defined when the $1$-form $\iota_{\xi_P}\omega$ is closed but not exact. The symplectic form of the sphere in Section~1.7 is closed, and it cannot be the differential of a global form, as recalled in Remark~3.6.3 for the projective Hilbert space. In Section~2.4 the Weyl operators compose according to the group law of $\mathbb{R}^2$ only up to the phases of Equation~(2.44), and no choice of phases turns them into a genuine representation. Similarly, Proposition~3.3.6 shows that a symmetry of $\mathbb{P}\H$ is implemented by a unitary operator determined only up to a phase, while Proposition~3.6.1 assumes from the outset that the group acts through a genuine unitary representation. In every case an identity satisfied locally is a necessary condition for a global object to exist, and the question is whether this condition is also sufficient. Cohomology is the algebraic framework which answers it. The necessary identities define the cocycles, the objects which satisfy them for trivial reasons define the coboundaries, and the quotient measures the obstruction.

This section develops cohomology in the generality required by the rest of the chapter. We start from de Rham cohomology, where the terms closed and exact were already used in Chapter 1, and we extract from it the abstract notion of a cochain complex. We then specialize this notion to groups, where it classifies projective unitary representations and central extensions, and to Lie algebras, where the cocycle identity becomes linear algebra. The two specializations are linked to each other and to differential forms by Proposition~\ref{prop:invariant-forms}, which identifies the cocycles of a Lie algebra with the closed left-invariant forms on the corresponding group. After showing how a cocycle arises in classical mechanics, we state Bargmann's theorem, which transfers the classification of projective representations from a simply connected group to its Lie algebra, and we illustrate the possible behaviours with three examples.

%--------------------------------------------------------------------
\subsection{Closed and Exact Forms}\label{subsec:de-rham}

The exterior derivative of Definition~1.1.24 satisfies $d\circ d=0$, so that in every degree the image of $d$ is contained in the kernel of $d$ taken in the next degree. These two subspaces of differential forms carry names which are used throughout the chapter.

\begin{definition}[Closed and exact forms]\label{def:closed-exact}
	Let $M$ be a smooth manifold as per Definition~1.1.4 and let $k\geq0$. With the convention $\Omega^{-1}(M):=\{0\}$, the space of \emph{closed} $k$-forms and the space of \emph{exact} $k$-forms on $M$ are
	\begin{equation*}
		Z^k(M):=\{\alpha\in\Omega^k(M)\mid d\alpha=0\}, \qquad B^k(M):=\{d\beta\mid\beta\in\Omega^{k-1}(M)\},
	\end{equation*}
	where $\Omega^k(M)$ is as per Definition~1.1.17.
\end{definition}

Every exact form is closed, since $d(d\beta)=0$ by item iii of Definition~1.1.24, so that $B^k(M)\subseteq Z^k(M)$. This inclusion is what allows us to quotient one space by the other, and the quotient counts the closed forms which are not exact.

\begin{definition}[de Rham cohomology]\label{def:de-rham}
	Let $M$ be a smooth manifold as per Definition~1.1.4 and let $k\geq0$. The \emph{$k$-th de Rham cohomology group} of $M$ is the quotient vector space
	\begin{equation*}
		H^k_{\mathrm{dR}}(M):=Z^k(M)/B^k(M),
	\end{equation*}
	with $Z^k(M)$ and $B^k(M)$ as per Definition~\ref{def:closed-exact}. The class of $\alpha\in Z^k(M)$ is denoted by $[\alpha]$.
\end{definition}

The equation $d\beta=\alpha$ for the unknown form $\beta$ has the closedness of $\alpha$ as a necessary condition, and $H^k_{\mathrm{dR}}(M)$ records whether this condition is also sufficient. In degree zero the answer can be computed at once, and it recovers a fact already met in Remark~1.4.14.

\begin{proposition}[Degree zero]\label{prop:H0}
	Let $M$ be a smooth manifold as per Definition~1.1.4. Then $H^0_{\mathrm{dR}}(M)$ is the space of locally constant functions on $M$, \textit{i.e.}, the space of real functions on the set of connected components of $M$. In particular $H^0_{\mathrm{dR}}(M)\cong\mathbb{R}$ if $M$ is connected.
\end{proposition}

\begin{proof}
	Since $B^0(M)=\{0\}$ we have $H^0_{\mathrm{dR}}(M)=Z^0(M)$. Let $f\in C^\infty(M)=\Omega^0(M)$ as per Remark~1.1.18, and let $(U,\varphi)$ be a chart with $\varphi(U)$ connected. By Remark~1.1.13, $df=\partial_i(f\circ\varphi^{-1})\,dx^i$ on $U$, hence $df=0$ on $U$ if and only if all partial derivatives of $f\circ\varphi^{-1}$ vanish, that is, if and only if $f$ is constant on $U$. Such charts cover $M$, so $df=0$ if and only if $f$ is locally constant. A locally constant function is constant on every connected component, and connected components of a manifold are open because $M$ is locally Euclidean.
\end{proof}

In degrees $k\geq1$ the condition $d\alpha=0$ is a differential equation and can be tested locally, whereas the existence of $\beta$ with $d\beta=\alpha$ is a global statement. The Poincar\'e lemma shows that the difference between the two lies entirely in the global shape of $M$.

\begin{theorem}[Poincar\'e lemma]\label{thm:poincare}
	Let $U\subseteq\mathbb{R}^n$ be an open set which is star-shaped with respect to one of its points. Then every closed $k$-form on $U$ with $k\geq1$ is exact.
\end{theorem}

\noindent See \cite[TODO-location]{TODO-Lee13}. Hence every closed form is exact on a coordinate ball, and a nonzero class in $H^k_{\mathrm{dR}}(M)$ is an obstruction which cannot be seen on small open sets.

The simplest manifold on which such an obstruction exists is the circle, which is the configuration space of the planar pendulum of Section~1.1. The form below is the model of a closed form which is not exact.

\begin{example}[The circle]\label{ex:circle-form}
	Let $j:S^1\hookrightarrow\mathbb{R}^2$ be the inclusion of the unit circle and let $\eta:=j^*(x\,dy-y\,dx)\in\Omega^1(S^1)$, where the pullback is as per Definition~1.1.26. The form $\eta$ is closed because $\Omega^2(S^1)=\{0\}$ by Remark~1.1.18. Suppose that $\eta=df$ for some $f\in C^\infty(S^1)$, and consider the smooth map $\gamma:\mathbb{R}\to S^1$, $\gamma(t)=(\cos t,\sin t)$. Since $d$ commutes with pullbacks, as recalled before Definition~1.1.26,
	\begin{equation*}
		\gamma^*\eta=(\cos^2t+\sin^2t)\,dt=dt, \qquad \gamma^*df=d(f\circ\gamma)=(f\circ\gamma)'\,dt.
	\end{equation*}
	Integrating over $[0,2\pi]$ yields $2\pi=f(\gamma(2\pi))-f(\gamma(0))=0$, which is a contradiction. Hence $[\eta]\neq0$ in $H^1_{\mathrm{dR}}(S^1)$.
\end{example}

The same phenomenon occurs in degree two for the symplectic forms of Chapter 1. Its proof requires the integration of top-degree forms, whose properties we recall.

\begin{theorem}[Stokes]\label{thm:stokes}
	Let $M$ be a nonempty compact oriented smooth manifold without boundary of dimension $m$. Then $\int_M\nu>0$ for every $m$-form $\nu$ which is positive with respect to the orientation, and
	\begin{equation*}
		\int_Md\beta=0, \qquad \beta\in\Omega^{m-1}(M).
	\end{equation*}
\end{theorem}

\noindent See \cite[TODO-location]{TODO-Lee13}. A nowhere vanishing $m$-form singles out an orientation for which it is positive, and this is the way Theorem~\ref{thm:stokes} is applied below.

\begin{proposition}[Compact symplectic forms are not exact]\label{prop:compact-not-exact}
	Let $(M,\omega)$ be a symplectic manifold as per Definition~1.3.8 which is nonempty, compact and without boundary. Then $\omega$ is closed and it is not exact, so that $[\omega]\neq0$ in $H^2_{\mathrm{dR}}(M)$.
\end{proposition}

\begin{proof}
	Closedness is part of Definition~1.3.8. Let $2n=\dim M$ and suppose that $\omega=d\alpha$ with $\alpha\in\Omega^1(M)$. Since $d\omega=0$, the antiderivation property of item i of Definition~1.1.24 yields $d(\omega^{n-1})=0$, and therefore
	\begin{equation*}
		d(\alpha\wedge\omega^{n-1})=d\alpha\wedge\omega^{n-1}-\alpha\wedge d(\omega^{n-1})=\omega^n.
	\end{equation*}
	By Corollary~1.3.4 the form $\omega^n$ is nowhere vanishing, so it fixes an orientation of $M$ for which it is positive. Theorem~\ref{thm:stokes} then gives both $\int_M\omega^n>0$ and $\int_M\omega^n=\int_Md(\alpha\wedge\omega^{n-1})=0$, which is a contradiction.
\end{proof}

\noindent This proves what Remark~3.6.3 asserts for $\mathbb{P}\H$ when $\dim\H<\infty$, and it applies to the sphere $S^2$ of Section~1.7.

The momentum maps of Chapter 1 involve closed $1$-forms instead. For these, the global obstruction is controlled by the fundamental group of $M$, in accordance with the hypothesis of simple connectivity appearing in Remark~3.6.3.

\begin{proposition}[Simply connected manifolds]\label{prop:simply-connected}
	Let $M$ be a connected and simply connected smooth manifold. Then every closed $1$-form on $M$ is exact, \textit{i.e.}, $H^1_{\mathrm{dR}}(M)=0$.
\end{proposition}

\noindent See \cite[TODO-location]{TODO-Lee13}. The idea is that the line integral of a closed $1$-form along a loop vanishes when the loop is contractible, so that $f(x):=\int_{x_0}^x\alpha$ is well defined and $df=\alpha$.

\begin{remark}\label{rem:momentum-existence}
	Let $(P,\omega)$ be a symplectic manifold as per Definition~1.3.8 and let $X\in\mathfrak{X}(P)$ satisfy $\mathcal{L}_X\omega=0$. By Remark~1.4.3 the $1$-form $\iota_X\omega$ is closed. If $P$ is simply connected, Proposition~\ref{prop:simply-connected} shows that $\iota_X\omega=dH$ for some $H\in C^\infty(P)$. Applied to $X=\xi_P$ this entails that the function $\hat J(\xi)$ of Definition~1.5.7 exists for every $\xi$ in the Lie algebra, which is the topological hypothesis used in Remark~3.6.3.
\end{remark}

%--------------------------------------------------------------------
\subsection{Cochain Complexes}\label{subsec:cochain}

The definition of $H^k_{\mathrm{dR}}(M)$ uses only two properties of $(\Omega^k(M),d)$, namely that $d$ raises the degree by one and that $d\circ d=0$. Abstracting them yields a notion which applies to any sequence of spaces connected by such operators, and in particular to the objects of Sections~\ref{subsec:group-cohomology} and~\ref{subsec:lie-cohomology}.

\begin{definition}[Cochain complex]\label{def:cochain-complex}
	A \emph{cochain complex} is a sequence $(C^k, d_k)_{k\geq0}$, where each $C^k$ is an abelian group, written additively, and each $d_k:C^k\to C^{k+1}$ is a group homomorphism such that $d_{k+1}\circ d_k=0$ for all $k\geq0$. The complex is denoted by $C^\bullet$.
\end{definition}

\begin{definition}[Cocycles, coboundaries and cohomology]\label{def:abstract-cohomology}
	Let $C^\bullet$ be a cochain complex as per Definition~\ref{def:cochain-complex}. The groups of \emph{cocycles} and of \emph{coboundaries} in degree $k$ are
	\begin{equation*}
		Z^k(C^\bullet):=\ker{d}_k, \qquad B^k(C^\bullet):=\operatorname{im} d_{k-1},
	\end{equation*}
	with $B^0(C^\bullet):=\{0\}$. The \emph{$k$-th cohomology group} of $C^\bullet$ is $H^k(C^\bullet):=Z^k(C^\bullet)/B^k(C^\bullet)$.
\end{definition}

Since $d_k\circ d_{k-1}=0$, one has $B^k(C^\bullet)\subseteq Z^k(C^\bullet)$, and the quotient is well defined. The de Rham complex $(\Omega^k(M),d)$ is the first example, with $Z^k$, $B^k$ and $H^k$ as in Definitions~\ref{def:closed-exact} and~\ref{def:de-rham}.

\begin{remark}[Cohomology as an obstruction]\label{rem:obstruction}
	In every cochain complex the equation $d_{k-1}x=z$ for the unknown $x\in C^{k-1}$ has $d_kz=0$ as a necessary condition, because $d_k\circ d_{k-1}=0$. The group $H^k(C^\bullet)$ measures the failure of this condition to be sufficient: it vanishes if and only if every cocycle is a coboundary. The three problems recalled at the beginning of the section have this form. In Remark~1.5.11 the equation is $d\hat J(\xi)=\iota_{\xi_P}\omega$, for the projective representations of Section~\ref{subsec:group-cohomology} it is the rephasing equation of Proposition~\ref{prop:rephasing}, and for Lie algebra extensions it is the equation of Proposition~\ref{prop:lie-extension}.
\end{remark}

Maps between complexes which respect the differentials descend to cohomology. This is what makes cohomology functorial, and it is the reason why classes can be compared across different manifolds, groups or algebras.

\begin{proposition}[Cochain maps]\label{prop:cochain-map}
	Let $C^\bullet$ and $D^\bullet$ be cochain complexes with differentials $d$ and $d'$, and let $\varphi^k:C^k\to D^k$ be group homomorphisms with $d'_k\circ\varphi^k=\varphi^{k+1}\circ d_k$ for all $k\geq0$. Then $\varphi^k$ maps $Z^k(C^\bullet)$ into $Z^k(D^\bullet)$ and $B^k(C^\bullet)$ into $B^k(D^\bullet)$, and it induces a homomorphism $H^k(\varphi):H^k(C^\bullet)\to H^k(D^\bullet)$, $[z]\mapsto[\varphi^k(z)]$.
\end{proposition}

\begin{proof}
	If $d_kz=0$, then $d'_k\varphi^k(z)=\varphi^{k+1}(d_kz)=0$. If $z=d_{k-1}x$, then $\varphi^k(z)=d'_{k-1}\varphi^{k-1}(x)$. Hence $\varphi^k$ descends to the quotients.
\end{proof}

\begin{example}[Pullback]\label{ex:pullback}
	Let $F:M\to N$ be a smooth map between smooth manifolds as per Definition~1.1.5. Since $d$ commutes with pullbacks, $F^*:\Omega^\bullet(N)\to\Omega^\bullet(M)$ is a cochain map, and Proposition~\ref{prop:cochain-map} yields $F^*:H^k_{\mathrm{dR}}(N)\to H^k_{\mathrm{dR}}(M)$. If $F$ is a diffeomorphism, then $(F^{-1})^*F^*=\mathrm{id}$ and $F^*(F^{-1})^*=\mathrm{id}$, so that $H^k_{\mathrm{dR}}(M)\cong H^k_{\mathrm{dR}}(N)$.
\end{example}

%--------------------------------------------------------------------

%=====================================================================
% PART A of the rewritten subsection "Group Cohomology, Multipliers and
% Central Extensions".
% Replaces everything from \subsection{Group Cohomology, ...} up to
% (and excluding) Definition~\ref{def:projective-representation}.
% Part B (projective representations, classification of extensions,
% commutator pairing, Klein example) follows in the next message.
%
% Environments assumed: definition, proposition, remark, example, proof.
% No custom macros. The only \cite is the pre-existing TODO placeholder of rem:topological-groups.
%
% Labels removed from this part (to be re-created in Part B):
%   rem:associator, prop:cocycle-extension, prop:commutator-pairing,
%   rem:class-of-multiplier, ex:klein-pauli.
% Labels moved here from later in the subsection:
%   def:central-extension, def:equiv-extensions.
% New labels: ex:heisenberg-group, eq:heisenberg-law, def:section,
%   def:defect, eq:defect, eq:product-defect, rem:holonomy, def:simplex,
%   rem:labels, ex:triangle-tetrahedron, rem:boundary-stokes,
%   prop:defect-properties, eq:defect-associativity, eq:change-section,
%   rem:picture-reading, ex:area-defect, fig:triangle, fig:tetrahedron,
%   fig:quadrilateral.
%=====================================================================

\subsection{Group Cohomology, Multipliers and Central Extensions}\label{subsec:group-cohomology}

By Proposition~3.3.6, a symmetry group acting on $\mathbb{P}\mathcal{H}$ by Kähler isomorphisms is implemented on $\mathcal{H}$ by unitary operators which are determined by the action on rays only up to a phase. Choosing one representative for each group element produces a family of unitary operators which composes according to the group law only up to phases, and the problem is whether a different choice of representatives can remove them. This is a problem of the same type as the ones recalled at the beginning of the section. A local identity, here the compatibility of the phases with the associativity of the product in the group, is necessary for the existence of a global object, here a genuine representation, and one asks whether it is also sufficient.

There is however a difference with the previous subsections. A group $G$, regarded as a set equipped with a multiplication, is not a manifold, so that it carries neither differential forms nor an exterior derivative. The structure of a cochain complex nevertheless survives, and its origin is geometric. Rather than stating the definitions at once, which would hide the reason for their form, the subsection builds them from a picture. The starting point is the extension problem, which consists in enlarging $G$ by a group of phases so that the phases become an additional coordinate and no longer an obstruction. Its solution requires a function of two group elements, called the defect, and two questions about it, namely which functions can occur and how they depend on a choice. We then represent this function as a number attached to a triangle, much as an area is, and we show that the associativity of the product is the statement that the four triangles bounding a tetrahedron are compatible with one another. Only at this point the cochains, the differentials, the cocycles and the coboundaries are defined, as the algebraic transcription of the picture and as the group counterpart of the closed and exact forms of Section~\ref{subsec:de-rham}. Projective unitary representations are treated afterwards, and their multiplier is recognized as the defect of an extension. The triangle and the tetrahedron remain the guiding intuition for the rest of the chapter, and the computations on the Galilei group of Section~\ref{sec:galilei} are instances of them.

%--------------------------------------------------------------------
\subsubsection*{The Extension Problem}

A symmetry group acts on rays, whereas we need it to act by unitary operators. Passing from the first to the second amounts to adjoining to the group a copy of the phases, and the mathematical notion describing this operation is the central extension. We state it for a general group $G$ and a general abelian group $A$, written additively, since the same notion will be used both with $A=\mathbb{R}$ and with $A=U(1)$. For $A=U(1)$ the group operation is the multiplication of phases, the zero element is $1$ and the opposite of $a$ is $a^{-1}$.

\begin{definition}[Central extension]\label{def:central-extension}
	Let $G$ be a group and let $A$ be an abelian group. A \emph{central extension} of $G$ by $A$ is a group $\widehat{G}$ together with an injective homomorphism $\iota:A\to\widehat{G}$ whose image is contained in the center of $\widehat{G}$, and a surjective homomorphism $\pi:\widehat{G}\to G$ with $\ker\pi=\iota(A)$.
\end{definition}

In a central extension the preimage $\pi^{-1}(\kappa)$ of an element $\kappa\in G$ is a coset of $\iota(A)$, \textit{i.e.}, a copy of $A$, and the elements of $\iota(A)$ commute with every element of $\widehat{G}$. The extra coordinate is therefore invisible in the quotient $G$, and its only trace is the way in which two elements multiply. To compare two extensions of the same group by the same group we need a notion of isomorphism which respects this structure.

\begin{definition}[Equivalent extensions]\label{def:equiv-extensions}
	Two central extensions $(\widehat{G},\iota,\pi)$ and $(\widehat{G}',\iota',\pi')$ of $G$ by $A$ are \emph{equivalent} if there is an isomorphism $\Theta:\widehat{G}\to\widehat{G}'$ with $\Theta\circ\iota=\iota'$ and $\pi'\circ\Theta=\pi$.
\end{definition}

The simplest extension is the direct product $A\times G$, with $\iota(a)=(a,e)$ and $\pi(a,\kappa)=\kappa$, which we call the trivial extension. The extension problem asks for all central extensions of $G$ by $A$ up to equivalence, and in particular for the ones which are not equivalent to the trivial one. The following example shows that such extensions exist, and it is the guiding example of the whole subsection.

\begin{example}[A nontrivial extension of $\mathbb{R}^2$]\label{ex:heisenberg-group}
	Let $G=\mathbb{R}^2$ with the addition and let $A=\mathbb{R}$. For $\mathbf{u}=(u_1,u_2)$ and $\mathbf{u}'=(u_1',u_2')$ we set $\mathbf{u}\wedge\mathbf{u}':=u_1u_2'-u_2u_1'$. On the set $\mathbb{R}\times\mathbb{R}^2$ we define
	\begin{equation}\label{eq:heisenberg-law}
		(a,\mathbf{u})(a',\mathbf{u}')=\Big(a+a'+\tfrac12\,\mathbf{u}\wedge\mathbf{u}',\ \mathbf{u}+\mathbf{u}'\Big).
	\end{equation}
	The product is associative, because $\mathbf{u}\wedge\mathbf{u}'$ is bilinear, so that both ways of bracketing three factors yield
	\begin{equation*}
		\Big(a+a'+a''+\tfrac12\big(\mathbf{u}\wedge\mathbf{u}'+\mathbf{u}\wedge\mathbf{u}''+\mathbf{u}'\wedge\mathbf{u}''\big),\ \mathbf{u}+\mathbf{u}'+\mathbf{u}''\Big).
	\end{equation*}
	The identity is $(0,\mathbf{0})$, the inverse of $(a,\mathbf{u})$ is $(-a,-\mathbf{u})$ because $\mathbf{u}\wedge\mathbf{u}=0$, and the elements $(a,\mathbf{0})$ are central. With $\iota(a)=(a,\mathbf{0})$ and $\pi(a,\mathbf{u})=\mathbf{u}$ this is a central extension of $\mathbb{R}^2$ by $\mathbb{R}$. Since $\mathbf{u}\wedge\mathbf{u}'=-\mathbf{u}'\wedge\mathbf{u}$, Equation~\eqref{eq:heisenberg-law} entails
	\begin{equation*}
		(a,\mathbf{u})(a',\mathbf{u}')=\iota\big(\mathbf{u}\wedge\mathbf{u}'\big)\,(a',\mathbf{u}')(a,\mathbf{u}),
	\end{equation*}
	so that two elements above $\mathbf{u}$ and $\mathbf{u}'$, which commute in $\mathbb{R}^2$, commute only up to the central element $\iota(\mathbf{u}\wedge\mathbf{u}')$. The trivial extension $\mathbb{R}\times\mathbb{R}^2$ is an abelian group, and an equivalence is an isomorphism, hence this extension is not equivalent to the trivial one. It is the Heisenberg Lie group of Definition~2.4.14 up to the sign of the central coordinate, as shown in Example~\ref{ex:heisenberg-extension}.
\end{example}

The example exhibits a phenomenon which is invisible at the level of $G$ alone, namely that the lifts of two commuting elements may fail to commute. To control it in general, we need a description of the product of $\widehat{G}$ in terms of data living on $G$ and $A$ only. The next paragraph provides it.

%--------------------------------------------------------------------
\subsubsection*{Sections and the Defect}

Let $(\widehat{G},\iota,\pi)$ be a central extension of $G$ by $A$. Since $\pi$ is surjective, we can choose for every $\kappa\in G$ one element of $\widehat{G}$ lying above it. Such a choice is the only extra datum at our disposal, and it is not required to respect the group law.

\begin{definition}[Section]\label{def:section}
	Let $(\widehat{G},\iota,\pi)$ be a central extension of $G$ by $A$ as per Definition~\ref{def:central-extension}. A \emph{section} is a map $\varsigma:G\to\widehat{G}$ with $\pi\circ\varsigma=\mathrm{id}_G$ and $\varsigma(e)=e$.
\end{definition}

A section exists because $\pi$ is surjective, and it is in general not a homomorphism, which is the point of the construction. Every element of $\widehat{G}$ is determined by its image in $G$ together with a displacement along $\iota(A)$. Indeed, for $x\in\widehat{G}$ the element $\kappa=\pi(x)$ is forced, and $x\,\varsigma(\kappa)^{-1}$ lies in $\ker\pi=\iota(A)$, so that it equals $\iota(a)$ for a unique $a\in A$ since $\iota$ is injective. Hence $x=\iota(a)\varsigma(\kappa)$ in a unique way, and $\widehat{G}$ is in bijection with $A\times G$. The multiplication of $\widehat{G}$ is then determined once we know how to multiply two sections. The elements $\varsigma(\kappa)\varsigma(\kappa')$ and $\varsigma(\kappa\kappa')$ lie above the same element $\kappa\kappa'$ of $G$, hence they differ by an element of $\iota(A)$.

\begin{definition}[Defect of a section]\label{def:defect}
	Let $\varsigma$ be a section as per Definition~\ref{def:section}. The \emph{defect} of $\varsigma$ is the map $\zeta_\varsigma:G\times G\to A$ uniquely determined by
	\begin{equation}\label{eq:defect}
		\varsigma(\kappa)\,\varsigma(\kappa')=\iota\big(\zeta_\varsigma(\kappa,\kappa')\big)\,\varsigma(\kappa\kappa'), \qquad \kappa,\kappa'\in G.
	\end{equation}
\end{definition}

The defect measures how far $\varsigma$ is from being a homomorphism, since $\zeta_\varsigma$ vanishes identically if and only if $\varsigma$ is one. Because $\iota(A)$ is central, the product of two arbitrary elements of $\widehat{G}$ is expressed through the defect alone,
\begin{equation}\label{eq:product-defect}
	\big(\iota(a)\varsigma(\kappa)\big)\big(\iota(a')\varsigma(\kappa')\big)=\iota\big(a+a'+\zeta_\varsigma(\kappa,\kappa')\big)\,\varsigma(\kappa\kappa').
\end{equation}
It descends that, in the coordinates $(a,\kappa)$, the group $\widehat{G}$ is the set $A\times G$ with the product $(a,\kappa)(a',\kappa')=(a+a'+\zeta_\varsigma(\kappa,\kappa'),\kappa\kappa')$. This is the form of Equation~\eqref{eq:heisenberg-law}, with $\zeta_\varsigma(\mathbf{u},\mathbf{u}')=\tfrac12\mathbf{u}\wedge\mathbf{u}'$. The defect therefore encodes the whole extension, and the extension problem splits into two questions. The first asks which functions $\zeta:G\times G\to A$ occur as defects, since only these can be used to build extensions. The second asks how the defect depends on the choice of the section, since two defects which differ by a change of section describe the same extension. Both questions have a short algebraic answer, which we state in Proposition~\ref{prop:defect-properties} below. Its content, however, is best understood through the picture of the next paragraph.

%--------------------------------------------------------------------
\subsubsection*{A Picture: Simplices with Vertices in a Group}

The identities which answer the two questions involve products of group elements, and they are difficult to read in the abstract. We therefore look for a picture in which the defect is a number attached to a geometric object, and in which the identities become statements on how such objects fit together. The Heisenberg group suggests the picture, and the following remark shows that the object involved is a triangle in general.

\begin{remark}[The defect as a displacement]\label{rem:holonomy}
	Equation~\eqref{eq:defect} can be rewritten as $\varsigma(\kappa)\varsigma(\kappa')\varsigma(\kappa\kappa')^{-1}=\iota\big(\zeta_\varsigma(\kappa,\kappa')\big)$. The left-hand side is the lifted product along a closed path of $G$ with three steps, from $e$ to $\kappa$, then from $\kappa$ to $\kappa\kappa'$ by multiplication with $\kappa'$ on the right, and last from $\kappa\kappa'$ back to $e$. The lift of this path does not close in $\widehat{G}$. It ends at $\iota(\zeta_\varsigma(\kappa,\kappa'))$, that is, it is displaced along the fibre $A$ by the defect. For the section $\varsigma(\mathbf{u})=(0,\mathbf{u})$ of Example~\ref{ex:heisenberg-group}, Equation~\eqref{eq:heisenberg-law} yields
	\begin{equation*}
		\varsigma(\mathbf{u})\varsigma(\mathbf{u}')\varsigma(\mathbf{u}+\mathbf{u}')^{-1}=\big(\tfrac12\,\mathbf{u}\wedge\mathbf{u}',\ \mathbf{0}\big),
	\end{equation*}
	and $\tfrac12\mathbf{u}\wedge\mathbf{u}'$ is the signed area of the triangle of the plane with vertices $\mathbf{0}$, $\mathbf{u}$ and $\mathbf{u}+\mathbf{u}'$. The displacement accumulated when going around the boundary of a triangle is its area, and in a general extension the defect plays the same role.
\end{remark}

To make this precise for an arbitrary group, we need a notion of triangle, and more generally of simplex, whose vertices are elements of $G$ itself. A $k$-simplex is the convex hull of $k+1$ points, \textit{i.e.}, a point for $k=0$, a segment for $k=1$, a triangle for $k=2$ and a tetrahedron for $k=3$. Only its combinatorial structure matters here, namely the ordered list of vertices and the faces obtained by omitting one of them. The order is needed to orient the edges, from the vertex of lower index to the vertex of higher index.

\begin{definition}[Simplex of a group]\label{def:simplex}
	Let $G$ be a group and let $k\geq0$. A \emph{$k$-simplex of $G$} is an ordered $(k+1)$-tuple $\sigma=(g_0,\dots,g_k)$ with $g_0,\dots,g_k\in G$, whose entries are called the vertices of $\sigma$.
\end{definition}

A simplex of $G$ has no geometry of its own, so that the quantity attached to an edge must be read off from the group. The relevant quantity is the group element which carries one vertex to the other. In the plane, with the addition as group law, it is the vector joining the two vertices, and this is the reading to keep in mind.

\begin{remark}[Edge labels and left invariance]\label{rem:labels}
	We label the edge from $g_i$ to $g_j$, with $i<j$, by the element $g_i^{-1}g_j\in G$. It is the element which carries $g_i$ to $g_j$ by left multiplication, since $g_i\,(g_i^{-1}g_j)=g_j$. The \emph{consecutive} edges carry the labels $\kappa_i:=g_{i-1}^{-1}g_i$ for $i=1,\dots,k$, and every other label is a product of consecutive ones, $g_i^{-1}g_j=\kappa_{i+1}\kappa_{i+2}\cdots\kappa_j$.
	
	Two simplices which differ by the simultaneous left multiplication of all vertices by a fixed $h\in G$ have the same labels, because $(hg_i)^{-1}(hg_j)=g_i^{-1}g_j$. Conversely, a simplex with consecutive labels $\kappa_1,\dots,\kappa_k$ is the left translate by $g_0$ of the simplex $(e,\kappa_1,\kappa_1\kappa_2,\dots,\kappa_1\cdots\kappa_k)$. Hence a function of $k$ group elements $(\kappa_1,\dots,\kappa_k)$ is the same thing as a function of $k$-simplices which takes the same value on simplices related by a left translation. This invariance is the reason why the defect $\zeta_\varsigma(\kappa,\kappa')$ is a function of two group elements and not of three vertices, since it describes the product of $\kappa$ and $\kappa'$ and does not depend on the vertex at which one decides to start. For $k=0$ the left translations act transitively on the vertices, so that an invariant function on $0$-simplices is a constant.
	
	If $\kappa_i=e$ for some $i$, then $g_{i-1}=g_i$ and the simplex is degenerate, \textit{i.e.}, collapsed along one edge. A function which vanishes on degenerate simplices is called normalized. The defect is normalized, since $\varsigma(e)=e$ entails $\zeta_\varsigma(\kappa,e)=\zeta_\varsigma(e,\kappa)=0$, as proved in Proposition~\ref{prop:defect-properties}, and a collapsed triangle has no area.
\end{remark}

We now name the edges and the faces of the two simplices which matter for the defect. The faces of a simplex are obtained by omitting one vertex, and their labels are computed from those of the simplex by the rule of Remark~\ref{rem:labels}.

\begin{example}[The triangle and the tetrahedron]\label{ex:triangle-tetrahedron}
	Let $(g_0,g_1,g_2)$ be a $2$-simplex with consecutive labels $\kappa_1=g_0^{-1}g_1$ and $\kappa_2=g_1^{-1}g_2$. We call $\kappa_1$ and $\kappa_2$ the short edges and $\kappa_1\kappa_2=g_0^{-1}g_2$ the long edge, as in Figure~\ref{fig:triangle}. Omitting $g_0$, $g_1$ and $g_2$ yields the three edges $(g_1,g_2)$, $(g_0,g_2)$ and $(g_0,g_1)$, with labels $\kappa_2$, $\kappa_1\kappa_2$ and $\kappa_1$.
	
	Let $(g_0,g_1,g_2,g_3)$ be a $3$-simplex with consecutive labels $\kappa_1,\kappa_2,\kappa_3$. It has six edges, as in Figure~\ref{fig:tetrahedron}. Three are the consecutive ones, with labels $\kappa_1$, $\kappa_2$ and $\kappa_3$. The edge $(g_0,g_2)$ has label $\kappa_1\kappa_2$, the edge $(g_1,g_3)$ has label $\kappa_2\kappa_3$ and the edge $(g_0,g_3)$, which we call the long edge, has label $\kappa_1\kappa_2\kappa_3$. It has four faces, which are triangles:
	\begin{center}
		\begin{tabular}{lll}
			\hline
			omitted vertex & face & labels of the short edges \\ \hline
			$g_0$ & $(g_1,g_2,g_3)$ & $(\kappa_2,\kappa_3)$ \\
			$g_1$ & $(g_0,g_2,g_3)$ & $(\kappa_1\kappa_2,\kappa_3)$ \\
			$g_2$ & $(g_0,g_1,g_3)$ & $(\kappa_1,\kappa_2\kappa_3)$ \\
			$g_3$ & $(g_0,g_1,g_2)$ & $(\kappa_1,\kappa_2)$ \\ \hline
		\end{tabular}
	\end{center}
\end{example}

\begin{figure}[ht]
	\centering
	% ---------------------------------------------------------------
	% PROMPT FOR A TIKZ GENERATOR (paste everything between the rules)
	% ---------------------------------------------------------------
	% Write TikZ code for a figure, to be compiled inside a LaTeX
	% document that already loads tikz. Use only the libraries
	% arrows.meta, calc and positioning. Black and white only, line
	% width 0.8pt, final width about 0.45\textwidth, labels in math mode
	% with the standard LaTeX math font.
	% Content: a triangle with three ordered vertices, g_0 bottom left,
	% g_1 bottom right, g_2 on top and roughly above the midpoint of g_0
	% and g_1. Write the vertex names $g_0$, $g_1$, $g_2$ outside the
	% triangle, next to the vertices, and mark each vertex with a small
	% filled dot. Draw the three edges as arrows with a closed arrow tip
	% (Stealth), directed from the lower to the higher index: g_0 to
	% g_1, g_1 to g_2, and g_0 to g_2. Label the edge g_0 to g_1 with
	% $\kappa_1=g_0^{-1}g_1$, the edge g_1 to g_2 with
	% $\kappa_2=g_1^{-1}g_2$, and the edge g_0 to g_2 with
	% $\kappa_1\kappa_2=g_0^{-1}g_2$. Place each label on the outer side
	% of its edge. Next to the labels of the two edges kappa_1 and
	% kappa_2 add the small gray italic text "short edge", and next to
	% the label of the edge kappa_1 kappa_2 add "long edge". Inside the
	% triangle draw a small counterclockwise circular arrow, indicating
	% the orientation g_0, g_1, g_2, and write next to it the symbol
	% $\zeta(\kappa_1,\kappa_2)$.
	% ---------------------------------------------------------------
	\fbox{\parbox{0.45\textwidth}{\centering\vspace{3cm}TODO: triangle figure\vspace{3cm}}}
	\caption{A $2$-simplex of $G$ with its consecutive labels $\kappa_1$ and $\kappa_2$ and its long edge $\kappa_1\kappa_2$. The defect $\zeta_\varsigma(\kappa_1,\kappa_2)$ is the number attached to the triangle.}
	\label{fig:triangle}
\end{figure}

\begin{figure}[ht]
	\centering
	% ---------------------------------------------------------------
	% PROMPT FOR A TIKZ GENERATOR (paste everything between the rules)
	% ---------------------------------------------------------------
	% Write TikZ code for a figure, to be compiled inside a LaTeX
	% document that already loads tikz. Use only the libraries
	% arrows.meta, calc and positioning. Black and white only, line
	% width 0.8pt, final width about 0.55\textwidth, labels in math mode
	% with the standard LaTeX math font.
	% Content: a tetrahedron drawn in oblique projection with four
	% ordered vertices and these coordinates (in cm): g_0 at (0,0),
	% g_1 at (2.2,-0.9), g_2 at (4.4,0), g_3 at (2.2,2.6). Mark every
	% vertex with a small filled dot and write the names $g_0$, $g_1$,
	% $g_2$, $g_3$ outside the solid, next to the vertices. Draw six
	% edges as arrows with a closed arrow tip (Stealth), each directed
	% from the lower to the higher index. The edge from g_0 to g_2 lies
	% behind the solid and must be dashed, all others are solid. Label
	% the edges, placing each label near the middle of its edge, as
	% follows: g_0 to g_1: $\kappa_1$, g_1 to g_2: $\kappa_2$,
	% g_2 to g_3: $\kappa_3$, g_0 to g_2: $\kappa_1\kappa_2$,
	% g_1 to g_3: $\kappa_2\kappa_3$, g_0 to g_3: $\kappa_1\kappa_2\kappa_3$.
	% Use a slightly smaller font for the three labels
	% $\kappa_1\kappa_2$, $\kappa_2\kappa_3$, $\kappa_1\kappa_2\kappa_3$
	% and avoid overlaps between labels and edges.
	% ---------------------------------------------------------------
	\fbox{\parbox{0.55\textwidth}{\centering\vspace{3.5cm}TODO: tetrahedron figure\vspace{3.5cm}}}
	\caption{A $3$-simplex of $G$ with its six edge labels. The edge from $g_0$ to $g_2$ lies behind the solid.}
	\label{fig:tetrahedron}
\end{figure}

The faces of a simplex come with signs, and these signs are the last ingredient of the picture. They are the same signs which appear in the boundary of a simplex in the cohomology of manifolds, where they have a precise meaning.

\begin{remark}[The boundary of a simplex and Stokes' theorem]\label{rem:boundary-stokes}
	The \emph{oriented boundary} of a simplex $\sigma=(g_0,\dots,g_k)$ is the formal sum $\partial\sigma:=\sum_{j=0}^k(-1)^j\sigma_j$, where $\sigma_j$ is the face obtained by omitting $g_j$. The alternating signs orient the faces coherently with the order of the vertices of $\sigma$. For a triangle one finds $\partial(g_0,g_1,g_2)=(g_1,g_2)-(g_0,g_2)+(g_0,g_1)$, which is the path $g_0\to g_1\to g_2$ followed by the path $g_2\to g_0$ along the long edge run backwards. For a tetrahedron
	\begin{equation*}
		\partial(g_0,g_1,g_2,g_3)=(g_1,g_2,g_3)-(g_0,g_2,g_3)+(g_0,g_1,g_3)-(g_0,g_1,g_2).
	\end{equation*}
	A function on simplices extends linearly to formal sums, so that it can be evaluated on $\partial\sigma$. This is the mechanism of the cohomology of manifolds. By Stokes' theorem, the exterior derivative of Definition~1.1.24 is characterized, in terms of integration on simplices, by $\int_\sigma d\alpha=\int_{\partial\sigma}\alpha$, so that $d\alpha$ is the function on $k$-simplices obtained by evaluating $\alpha$ on the boundary. Closed forms are the ones with $d\alpha=0$, exact forms are the ones of the shape $d\beta$, and the identity $d\circ d=0$ expresses that the boundary of a boundary is empty. We keep this heuristic reading and we apply the same recipe to functions of simplices of $G$, which are invariant under left translation by Remark~\ref{rem:labels}.
\end{remark}

%--------------------------------------------------------------------
\subsubsection*{The Two Questions in the Picture}

With the picture in hand we can answer the two questions on the defect. The answers are two short computations, which we state first and then read geometrically.

\begin{proposition}[Properties of the defect]\label{prop:defect-properties}
	Let $(\widehat{G},\iota,\pi)$ be a central extension of $G$ by $A$ as per Definition~\ref{def:central-extension}, let $\varsigma$ be a section and let $\zeta:=\zeta_\varsigma$ be its defect as per Definition~\ref{def:defect}.
	\begin{enumerate}
		\item $\zeta(\kappa,e)=\zeta(e,\kappa)=0$ for all $\kappa\in G$, and for all $\kappa_1,\kappa_2,\kappa_3\in G$
		\begin{equation}\label{eq:defect-associativity}
			\zeta(\kappa_1,\kappa_2)+\zeta(\kappa_1\kappa_2,\kappa_3)=\zeta(\kappa_2,\kappa_3)+\zeta(\kappa_1,\kappa_2\kappa_3).
		\end{equation}
		\item If $\varsigma'$ is another section, there is a unique map $\mu:G\to A$ with $\mu(e)=0$ and $\varsigma'(\kappa)=\iota(\mu(\kappa))\varsigma(\kappa)$, and
		\begin{equation}\label{eq:change-section}
			\zeta_{\varsigma'}(\kappa,\kappa')=\zeta_\varsigma(\kappa,\kappa')+\mu(\kappa)+\mu(\kappa')-\mu(\kappa\kappa').
		\end{equation}
	\end{enumerate}
\end{proposition}

\begin{proof}
	Item 1. Since $\varsigma(e)=e$, Equation~\eqref{eq:defect} with $\kappa'=e$ reads $\varsigma(\kappa)=\iota(\zeta(\kappa,e))\varsigma(\kappa)$, hence $\iota(\zeta(\kappa,e))=e$ and $\zeta(\kappa,e)=0$ because $\iota$ is injective. The same argument with $\kappa=e$ yields $\zeta(e,\kappa)=0$. For the identity we bracket $\varsigma(\kappa_1)\varsigma(\kappa_2)\varsigma(\kappa_3)$ in the two possible ways, using Equation~\eqref{eq:defect} twice and the centrality of $\iota(A)$ to collect the factors in $\iota(A)$ on the left,
	\begin{align*}
		\big(\varsigma(\kappa_1)\varsigma(\kappa_2)\big)\varsigma(\kappa_3)&=\iota\big(\zeta(\kappa_1,\kappa_2)\big)\varsigma(\kappa_1\kappa_2)\varsigma(\kappa_3)\\
		&=\iota\big(\zeta(\kappa_1,\kappa_2)+\zeta(\kappa_1\kappa_2,\kappa_3)\big)\varsigma(\kappa_1\kappa_2\kappa_3),\\
		\varsigma(\kappa_1)\big(\varsigma(\kappa_2)\varsigma(\kappa_3)\big)&=\iota\big(\zeta(\kappa_2,\kappa_3)\big)\varsigma(\kappa_1)\varsigma(\kappa_2\kappa_3)\\
		&=\iota\big(\zeta(\kappa_2,\kappa_3)+\zeta(\kappa_1,\kappa_2\kappa_3)\big)\varsigma(\kappa_1\kappa_2\kappa_3).
	\end{align*}
	The two expressions coincide by associativity of the product of $\widehat{G}$, and $\iota$ is injective, which yields Equation~\eqref{eq:defect-associativity}.
	
	Item 2. The elements $\varsigma'(\kappa)$ and $\varsigma(\kappa)$ lie above the same element of $G$, so that $\varsigma'(\kappa)\varsigma(\kappa)^{-1}\in\ker\pi=\iota(A)$ defines a unique $\mu(\kappa)\in A$, and $\mu(e)=0$ because $\varsigma'(e)=\varsigma(e)=e$. Using the centrality of $\iota(A)$ and $\varsigma(\kappa\kappa')=\iota(-\mu(\kappa\kappa'))\varsigma'(\kappa\kappa')$,
	\begin{align*}
		\varsigma'(\kappa)\varsigma'(\kappa')&=\iota\big(\mu(\kappa)+\mu(\kappa')+\zeta_\varsigma(\kappa,\kappa')\big)\varsigma(\kappa\kappa')\\
		&=\iota\big(\mu(\kappa)+\mu(\kappa')+\zeta_\varsigma(\kappa,\kappa')-\mu(\kappa\kappa')\big)\varsigma'(\kappa\kappa'),
	\end{align*}
	which is Equation~\eqref{eq:change-section} by Definition~\ref{def:defect}.
\end{proof}

\noindent The proposition answers the second question completely, while for the first one it provides a necessary condition, whose sufficiency is proved in Proposition~\ref{prop:cocycle-extension}.

The two equations of the proposition acquire a meaning once they are read on the simplices of Example~\ref{ex:triangle-tetrahedron}. The reading shows that Equation~\eqref{eq:defect-associativity} is a closure condition of the tetrahedron, and that it is the exact counterpart of the associativity of the product of $\widehat{G}$.

\begin{remark}[Reading Proposition~\ref{prop:defect-properties} on simplices]\label{rem:picture-reading}
	\emph{Associativity is the closure of the tetrahedron.} Consider four vertices $g_0,\dots,g_3$ and the quadrilateral formed by the three consecutive edges $\kappa_1,\kappa_2,\kappa_3$ together with the long edge from $g_0$ to $g_3$, as in Figure~\ref{fig:quadrilateral}. It can be cut into two triangles in two ways. The cut along the diagonal from $g_0$ to $g_2$ yields the triangles $(g_0,g_1,g_2)$ and $(g_0,g_2,g_3)$, to which the defect assigns $\zeta(\kappa_1,\kappa_2)$ and $\zeta(\kappa_1\kappa_2,\kappa_3)$. The cut along the diagonal from $g_1$ to $g_3$ yields $(g_1,g_2,g_3)$ and $(g_0,g_1,g_3)$, with values $\zeta(\kappa_2,\kappa_3)$ and $\zeta(\kappa_1,\kappa_2\kappa_3)$. Equation~\eqref{eq:defect-associativity} states that the total value on the quadrilateral does not depend on the diagonal, just as the area of a quadrilateral does not. The four triangles are the four faces of the tetrahedron $(g_0,g_1,g_2,g_3)$, and with the signs of Remark~\ref{rem:boundary-stokes} the equation reads $\zeta(\text{face }0)-\zeta(\text{face }1)+\zeta(\text{face }2)-\zeta(\text{face }3)=0$. The defect evaluated on the boundary of a tetrahedron vanishes, \textit{i.e.}, the tetrahedron closes.
	
	The two diagonals correspond to the two bracketings of a product of three factors. The first merges $\varsigma(\kappa_1)$ and $\varsigma(\kappa_2)$ into a section above $\kappa_1\kappa_2$, which is the cut along the diagonal from $g_0$ to $g_2$, and the second merges $\varsigma(\kappa_2)$ and $\varsigma(\kappa_3)$, which is the cut along the diagonal from $g_1$ to $g_3$. This explains the proof of item 1: associativity of $\widehat{G}$ is the independence of the total defect from the way the quadrilateral is cut.
	
	\emph{A change of section is a boundary.} A section is a choice of a number $\mu(\kappa)$ for every edge. Equation~\eqref{eq:change-section} states that changing the section modifies the value on a triangle by $\mu(\kappa_1)+\mu(\kappa_2)-\mu(\kappa_1\kappa_2)$, the oriented sum of $\mu$ on the boundary of the triangle. This modification does not spoil the closure of the tetrahedron. Each of the six edges of a tetrahedron belongs to exactly two of its four faces, and it enters their boundaries with opposite signs, so that the corresponding contributions cancel. This is the statement that the boundary of a boundary is empty, and it is the content of Proposition~\ref{prop:group-complex} below. It is also the reason why a change of section turns a defect into another defect.
\end{remark}

\begin{figure}[ht]
	\centering
	% ---------------------------------------------------------------
	% PROMPT FOR A TIKZ GENERATOR (paste everything between the rules)
	% ---------------------------------------------------------------
	% Write TikZ code for a figure made of two quadrilaterals side by
	% side, separated by the symbol "=" drawn in the middle, to be
	% compiled inside a LaTeX document that already loads tikz. Use only
	% the libraries arrows.meta, calc and positioning. Black and white
	% (light gray fills allowed), line width 0.8pt, total width about
	% 0.95\textwidth, labels in math mode with the standard LaTeX math
	% font.
	% Each quadrilateral is a square of side 2cm with vertices g_0 at
	% (0,0), g_1 at (2,0), g_2 at (2,2), g_3 at (0,2). Write the names
	% $g_0$, $g_1$, $g_2$, $g_3$ outside the square next to the vertices
	% and mark the vertices with small filled dots. Draw the boundary as
	% arrows with a closed arrow tip (Stealth): g_0 to g_1 labelled
	% $\kappa_1$, g_1 to g_2 labelled $\kappa_2$, g_2 to g_3 labelled
	% $\kappa_3$, and the long edge g_0 to g_3 labelled
	% $\kappa_1\kappa_2\kappa_3$ on the left side. Place the labels
	% outside the square.
	% Left square: draw the diagonal from g_0 to g_2 as an arrow. Fill
	% the two triangles (g_0,g_1,g_2) and (g_0,g_2,g_3) with two
	% different light gray tones and write inside them
	% $\zeta(\kappa_1,\kappa_2)$ and $\zeta(\kappa_1\kappa_2,\kappa_3)$
	% respectively. Label the diagonal $\kappa_1\kappa_2$.
	% Right square: draw the diagonal from g_1 to g_3 as an arrow. Fill
	% the two triangles (g_1,g_2,g_3) and (g_0,g_1,g_3) with two
	% different light gray tones and write inside them
	% $\zeta(\kappa_2,\kappa_3)$ and $\zeta(\kappa_1,\kappa_2\kappa_3)$
	% respectively. Label the diagonal $\kappa_2\kappa_3$.
	% Under the left square write the caption text
	% "bracketing $(\varsigma(\kappa_1)\varsigma(\kappa_2))\varsigma(\kappa_3)$"
	% and under the right square
	% "bracketing $\varsigma(\kappa_1)(\varsigma(\kappa_2)\varsigma(\kappa_3))$".
	% ---------------------------------------------------------------
	\fbox{\parbox{0.9\textwidth}{\centering\vspace{3.5cm}TODO: quadrilateral figure\vspace{3.5cm}}}
	\caption{The two ways of cutting the quadrilateral $(g_0,g_1,g_2,g_3)$ into triangles. Equation~\eqref{eq:defect-associativity} states that the total value of the defect is the same for both.}
	\label{fig:quadrilateral}
\end{figure}

The guiding example makes the closure condition literal, since the defect of the Heisenberg group is an area and areas are additive.

\begin{example}[The defect of the Heisenberg group]\label{ex:area-defect}
	Consider the section $\varsigma(\mathbf{u})=(0,\mathbf{u})$ of Example~\ref{ex:heisenberg-group}. Equation~\eqref{eq:heisenberg-law} yields $\varsigma(\mathbf{u})\varsigma(\mathbf{u}')=\iota\big(\tfrac12\mathbf{u}\wedge\mathbf{u}'\big)\varsigma(\mathbf{u}+\mathbf{u}')$, hence $\zeta_\varsigma(\mathbf{u},\mathbf{u}')=\tfrac12\mathbf{u}\wedge\mathbf{u}'$. For a triangle $(g_0,g_1,g_2)$ of $\mathbb{R}^2$ the consecutive labels are $\mathbf{u}=g_1-g_0$ and $\mathbf{u}'=g_2-g_1$, and the signed area of the triangle is
	\begin{equation*}
		\tfrac12\,(g_1-g_0)\wedge(g_2-g_0)=\tfrac12\,\mathbf{u}\wedge(\mathbf{u}+\mathbf{u}')=\tfrac12\,\mathbf{u}\wedge\mathbf{u}',
	\end{equation*}
	which is also $\tfrac12(g_0\wedge g_1+g_1\wedge g_2+g_2\wedge g_0)$. Hence $\zeta_\varsigma$ is the signed area, and Equation~\eqref{eq:defect-associativity} is the additivity of the area. In the alternating sum of the four faces of a tetrahedron $(g_0,g_1,g_2,g_3)$ each of the six terms $g_i\wedge g_j$ with $i<j$ occurs in exactly two faces with opposite signs, as one checks by direct expansion, so that the sum vanishes.
	
	The defect is not the differential of a $1$-cochain. Since $\mathbb{R}^2$ is abelian, every expression $\mu(\mathbf{u})+\mu(\mathbf{u}')-\mu(\mathbf{u}+\mathbf{u}')$ is symmetric in $(\mathbf{u},\mathbf{u}')$, whereas $\zeta_\varsigma(\mathbf{e}_1,\mathbf{e}_2)=\tfrac12=-\zeta_\varsigma(\mathbf{e}_2,\mathbf{e}_1)$ for the canonical basis. By Equation~\eqref{eq:change-section}, no section of this extension has vanishing defect, in agreement with the fact that the extension is not abelian and therefore not trivial.
\end{example}

%--------------------------------------------------------------------
\subsubsection*{The Algebraic Transcription: Cochains and Differentials}

We now transcribe the picture into definitions. A $k$-cochain is a normalized function of $k$ group elements, \textit{i.e.}, a normalized function on $k$-simplices which is invariant under left translation by Remark~\ref{rem:labels}. The differential is the evaluation on the oriented boundary of Remark~\ref{rem:boundary-stokes}, written in terms of the labels of the faces computed in Example~\ref{ex:triangle-tetrahedron}. We keep the cochains of degree up to three, which suffice for the classification of central extensions.

\begin{definition}[Cochain complex of a group]\label{def:group-cochains}
	Let $G$ be a group with identity $e$ and let $A$ be an abelian group, written additively, on which $G$ acts trivially. The \emph{normalized cochains} of $G$ with values in $A$ are
	\begin{align*}
		C^0(G,A)&:=A,\\
		C^1(G,A)&:=\{\mu:G\to A\mid\mu(e)=0\},\\
		C^2(G,A)&:=\big\{\zeta:G\times G\to A\mid\zeta(\kappa,e)=\zeta(e,\kappa)=0\ \ \forall\kappa\in G\big\},\\
		C^3(G,A)&:=\{\text{maps }G\times G\times G\to A\},
	\end{align*}
	with $C^k(G,A):=\{0\}$ for $k\geq4$. The \emph{differentials} are $d_0:=0$,
	\begin{align*}
		(d_1\mu)(\kappa,\kappa')&:=\mu(\kappa)+\mu(\kappa')-\mu(\kappa\kappa'),\\
		(d_2\zeta)(\kappa,\kappa',\kappa'')&:=\zeta(\kappa',\kappa'')-\zeta(\kappa\kappa',\kappa'')+\zeta(\kappa,\kappa'\kappa'')-\zeta(\kappa,\kappa'),
	\end{align*}
	for all $\kappa,\kappa',\kappa''\in G$, and $d_k:=0$ for $k\geq3$.
\end{definition}

The definition is the dictionary of the picture. Elements of $C^1(G,A)$ attach a value to each edge, elements of $C^2(G,A)$ to each triangle and elements of $C^3(G,A)$ to each tetrahedron, and the trivial action of $G$ on $A$ reflects the centrality of $\iota(A)$. The normalization requires a cochain to vanish on degenerate simplices, which is the property recorded in item 1 of Proposition~\ref{prop:defect-properties} and in the condition $\mu(e)=0$. The differential $d_1\mu$ is the sum of $\mu$ on the oriented boundary of the triangle $(g_0,g_1,g_2)$, the two short edges with the sign $+$ and the long edge with the sign $-$. The differential $d_2\zeta$ is the sum of $\zeta$ on the oriented boundary of the tetrahedron $(g_0,g_1,g_2,g_3)$, with the faces omitting $g_0,g_1,g_2,g_3$ carrying the signs $+,-,+,-$ and the values read in the table of Example~\ref{ex:triangle-tetrahedron}. For $A=U(1)$ the group operation is the multiplication, the zero element is $1$ and the opposite of $a$ is $a^{-1}$, so that the differentials read
\begin{equation*}
	(d_1\mu)(\kappa,\kappa')=\frac{\mu(\kappa)\,\mu(\kappa')}{\mu(\kappa\kappa')}, \qquad (d_2\omega)(\kappa,\kappa',\kappa'')=\frac{\omega(\kappa',\kappa'')\,\omega(\kappa,\kappa'\kappa'')}{\omega(\kappa\kappa',\kappa'')\,\omega(\kappa,\kappa')}.
\end{equation*}
In this notation Equation~\eqref{eq:change-section} reads $\zeta_{\varsigma'}=\zeta_\varsigma+d_1\mu$, and Equation~\eqref{eq:defect-associativity} reads $d_2\zeta_\varsigma=0$.

\begin{remark}[The differential in every degree]\label{rem:group-general}
	The recipe of Remark~\ref{rem:boundary-stokes} produces a differential in every degree. For a $(k+1)$-simplex with consecutive labels $\kappa_1,\dots,\kappa_{k+1}$, the face omitting $g_0$ has consecutive labels $(\kappa_2,\dots,\kappa_{k+1})$, the face omitting $g_{k+1}$ has $(\kappa_1,\dots,\kappa_k)$, and the face omitting $g_j$ with $1\leq j\leq k$ has the labels in which $\kappa_j$ and $\kappa_{j+1}$ are replaced by their product $\kappa_j\kappa_{j+1}$. Evaluating a $k$-cochain $\zeta$ on the oriented boundary yields
	\begin{multline*}
		(d_k\zeta)(\kappa_1,\dots,\kappa_{k+1})=\zeta(\kappa_2,\dots,\kappa_{k+1})\\
		+\sum_{j=1}^k(-1)^j\zeta(\kappa_1,\dots,\kappa_j\kappa_{j+1},\dots,\kappa_{k+1})+(-1)^{k+1}\zeta(\kappa_1,\dots,\kappa_k),
	\end{multline*}
	which reduces to Definition~\ref{def:group-cochains} for $k=1,2$. The identity $d\circ d=0$ is the statement that the boundary of a boundary is empty. We only need it in degree two, where it is verified in Proposition~\ref{prop:group-complex}.
\end{remark}

\begin{proposition}\label{prop:group-complex}
	In Definition~\ref{def:group-cochains}, $d_1$ takes values in $C^2(G,A)$ and $d_2\circ d_1=0$. Hence $C^\bullet(G,A)$ is a cochain complex as per Definition~\ref{def:cochain-complex}.
\end{proposition}

\begin{proof}
	For $\mu\in C^1(G,A)$ we have $(d_1\mu)(\kappa,e)=\mu(\kappa)+\mu(e)-\mu(\kappa)=0$, and likewise $(d_1\mu)(e,\kappa)=0$. Substituting the definition of $d_1$ in that of $d_2$ yields
	\begin{align*}
		(d_2d_1\mu)(\kappa,\kappa',\kappa'')={}&[\mu(\kappa')+\mu(\kappa'')-\mu(\kappa'\kappa'')]-[\mu(\kappa\kappa')+\mu(\kappa'')-\mu(\kappa\kappa'\kappa'')]\\
		&+[\mu(\kappa)+\mu(\kappa'\kappa'')-\mu(\kappa\kappa'\kappa'')]-[\mu(\kappa)+\mu(\kappa')-\mu(\kappa\kappa')],
	\end{align*}
	and the twelve terms cancel in pairs. The remaining requirements of Definition~\ref{def:cochain-complex} hold because $d_0=0$ and $d_k=0$ for $k\geq3$.
\end{proof}

\noindent In the picture, the cancellation is the one of Remark~\ref{rem:picture-reading}, since each edge of the tetrahedron appears in two faces with opposite signs.

Since $C^\bullet(G,A)$ is a cochain complex, Definition~\ref{def:abstract-cohomology} applies and we can name the objects which answer the two questions on the defect. Cocycles are the cochains which are closed, in the sense of Remark~\ref{rem:boundary-stokes}, and coboundaries are the ones which are exact.

\begin{definition}[Cocycle]\label{def:cocycle}
	Let $G$ be a group with identity $e$ and let $A$ be an abelian group, written additively. A \emph{$2$-cocycle} on $G$ with values in $A$ is a map $\zeta\in C^2(G,A)$ with $d_2\zeta=0$, \textit{i.e.}, such that
	\begin{equation*}
		\zeta(\kappa,\kappa')+\zeta(\kappa\kappa',\kappa'')=\zeta(\kappa',\kappa'')+\zeta(\kappa,\kappa'\kappa''), \qquad \kappa,\kappa',\kappa''\in G.
	\end{equation*}
	A cocycle with values in $A=\mathbb{R}$ is called a real cocycle.
\end{definition}

By Proposition~\ref{prop:defect-properties}, the defect of every section of a central extension is a $2$-cocycle. Changing the section modifies it by the differential of a $1$-cochain, which is the structure singled out by the next definition.

\begin{definition}[Coboundary and second cohomology group]\label{def:cohomology-group}
	Let $G$ and $A$ be as in Definition~\ref{def:cocycle}. A $2$-cocycle $\zeta$ is a \emph{coboundary} if $\zeta=d_1\mu$ for some $\mu\in C^1(G,A)$, \textit{i.e.}, if $\zeta(\kappa,\kappa')=\mu(\kappa)+\mu(\kappa')-\mu(\kappa\kappa')$. Two cocycles are \emph{cohomologous} if their difference is a coboundary. The \emph{second cohomology group} $H^2(G,A)$ is the cohomology group of degree two of $C^\bullet(G,A)$ as per Definition~\ref{def:abstract-cohomology}.
\end{definition}

In degrees zero and one the same construction yields $H^0(G,A)=A$ and $H^1(G,A)=\{\mu:G\to A\mid\mu(\kappa\kappa')=\mu(\kappa)+\mu(\kappa')\}$, the group of homomorphisms from $G$ to $A$, because $d_0=0$. The group $H^2(G,A)$ is the object of interest, and it plays for the extension problem the role that $H^k_{\mathrm{dR}}(M)$ plays for the existence of a primitive of a closed form, in the sense of Remark~\ref{rem:obstruction}. The defects of the sections of one extension form a single class in $H^2(G,A)$ by Equation~\eqref{eq:change-section}, and the class vanishes if and only if some section has vanishing defect, \textit{i.e.}, if and only if the extension is trivial. The converse, that every class arises from an extension and that distinct classes yield inequivalent extensions, is the content of Proposition~\ref{prop:cocycle-extension}.

\begin{remark}[Topological groups]\label{rem:topological-groups}
	The cochains of Definition~\ref{def:group-cochains} are arbitrary maps, which is the appropriate choice for abstract groups. For a Lie group one works with Borel cochains \cite[TODO-location]{TODO-Varadarajan1985}, and all statements below on multipliers of Lie groups are to be read in this setting. Without a regularity condition the cohomology groups become much larger. For $G=\mathbb{R}^2$ every bi-additive antisymmetric map $B:\mathbb{R}^2\times\mathbb{R}^2\to\mathbb{R}$ produces the cocycle $\frac12B$ by the same computation as in Example~\ref{ex:area-defect}, and bi-additive maps need not be bilinear because additive maps $\mathbb{R}\to\mathbb{R}$ need not be continuous.
\end{remark}

The picture also explains why this machinery is needed for the Galilei group. In Example~\ref{ex:heisenberg-group} two commuting elements of $G$ have lifts which commute only up to a central element, and this element is the area of the parallelogram spanned by them. The same mechanism occurs for the boosts and the translations of the Galilei group, which commute in $\mathcal{G}$ while their lifts differ by a phase, as Lemma~\ref{lem:galilei-relations} shows. The size of this phase, which no choice of section can remove, is the mass, and it is read in Section~\ref{sec:galilei} as the value of an invariant cocycle in the same way as the area is read here. To transfer these statements to unitary operators, we now need the notion of projective representation, which is the object of the next paragraphs.

\subsubsection*{Projective Unitary Representations}

The subsection so far has been abstract, with a generic abelian group $A$. We now return to the problem which motivated it, and we take $A=U(1)$. A symmetry acts on rays, and its implementation on $\mathcal{H}$ is a family of unitary operators which is fixed only up to phases. We first define this family precisely. We then compute the two facts about its phases which the cochains of Definition~\ref{def:group-cochains} were designed to record, namely the identity which the phases satisfy and their transformation under a change of representatives. These are the computations of Proposition~\ref{prop:defect-properties}, carried out in the setting in which the central group is $U(1)$.

\begin{definition}[Projective unitary representation]\label{def:projective-representation}
	Let $G$ be a group with identity $e$ and let $\mathcal{H}$ be a Hilbert space as per Definition~2.2.6. A \emph{projective unitary representation} of $G$ on $\mathcal{H}$ is a map $V:G\to U(\mathcal{H})$, $\kappa\mapsto V_\kappa$, taking values in the unitary operators of Definition~2.2.22 with $V_e=\mathbb{1}$, such that for all $\kappa,\kappa'\in G$ there is $\omega(\kappa,\kappa')\in U(1)$ with
	\begin{equation*}
		V_\kappa V_{\kappa'}=\omega(\kappa,\kappa')\,V_{\kappa\kappa'}.
	\end{equation*}
	The function $\omega:G\times G\to U(1)$ is called the \emph{multiplier} of $V$.
\end{definition}

A projective unitary representation is a bookkeeping device for a symmetry of rays, and its relation with the extensions of the previous paragraphs is the content of the following remark.

\begin{remark}\label{rem:projective-quotient}
	A projective unitary representation is the same thing as a homomorphism of $G$ into the quotient group $U(\mathcal{H})/U(1)$, \textit{i.e.}, into the group of unitary transformations of the rays of $\mathcal{H}$, together with a choice of one representative for each group element. The multiplier records the discrepancy between this choice and the group law. The map $V$ is a genuine representation as per Definition~1.2.18 if and only if $\omega\equiv1$, and it is a strongly continuous one as per Definition~1.2.28 if moreover $\kappa\mapsto V_\kappa$ is strongly continuous. If $G$ is a Lie group, we call $V$ continuous if $\kappa\mapsto|(x,V_\kappa y)|$ is continuous for all $x,y\in\mathcal{H}$.
\end{remark}

The choice of one representative for each element is the analogue of the choice of a section in Definition~\ref{def:section}, and the multiplier is the analogue of the defect of Definition~\ref{def:defect}. The analogy is made exact in Proposition~\ref{prop:cocycle-extension}, and for the moment it explains why the multiplier must satisfy the identity of Proposition~\ref{prop:defect-properties}. We verify it directly, since it is the first appearance of a cocycle in representation theory.

\begin{proposition}[The multiplier is a cocycle]\label{prop:multiplier-cocycle}
	Let $V$ be a projective unitary representation of $G$ with multiplier $\omega$. Then $\omega(\kappa,e)=\omega(e,\kappa)=1$ and
	\begin{equation}\label{eq:multiplier-identity}
		\omega(\kappa,\kappa')\,\omega(\kappa\kappa',\kappa'')=\omega(\kappa',\kappa'')\,\omega(\kappa,\kappa'\kappa''), \qquad \kappa,\kappa',\kappa''\in G.
	\end{equation}
\end{proposition}

\begin{proof}
	Since $V_e=\mathbb{1}$, the relations $V_\kappa=V_\kappa V_e=\omega(\kappa,e)V_\kappa$ and $V_\kappa=V_eV_\kappa=\omega(e,\kappa)V_\kappa$ entail $\omega(\kappa,e)=\omega(e,\kappa)=1$, because $V_\kappa$ is invertible. Bracketing the product $V_\kappa V_{\kappa'}V_{\kappa''}$ in the two possible ways yields
	\begin{align*}
		(V_\kappa V_{\kappa'})V_{\kappa''}&=\omega(\kappa,\kappa')\,\omega(\kappa\kappa',\kappa'')\,V_{\kappa\kappa'\kappa''},\\
		V_\kappa(V_{\kappa'}V_{\kappa''})&=\omega(\kappa',\kappa'')\,\omega(\kappa,\kappa'\kappa'')\,V_{\kappa\kappa'\kappa''}.
	\end{align*}
	The two expressions coincide by associativity of the product of operators, and $V_{\kappa\kappa'\kappa''}$ is invertible, so that Equation~\eqref{eq:multiplier-identity} holds.
\end{proof}

\noindent Equation~\eqref{eq:multiplier-identity} is the associativity of the product of operators, rewritten in terms of the phases, \textit{i.e.}, the closure of the tetrahedron of Remark~\ref{rem:picture-reading} with the values of $\omega$ on its faces. It is the only constraint which we will find on $\omega$, and Definition~\ref{def:cocycle} names it.

The operators $V_\kappa$ are determined by the action on rays only up to a phase. A change of representatives modifies the multiplier in a way which can be computed explicitly, and it is natural to ask whether this freedom can be used to make the multiplier trivial.

\begin{proposition}[Rephasing]\label{prop:rephasing}
	Let $V$ be a projective unitary representation of $G$ with multiplier $\omega$, and let $\mu:G\to U(1)$ be a map with $\mu(e)=1$. Then $V'_\kappa:=\mu(\kappa)V_\kappa$ is a projective unitary representation with multiplier
	\begin{equation}\label{eq:rephasing}
		\omega'(\kappa,\kappa')=\omega(\kappa,\kappa')\,\frac{\mu(\kappa)\,\mu(\kappa')}{\mu(\kappa\kappa')}.
	\end{equation}
	In particular, $V$ can be rephased into a genuine representation if and only if there is such a map $\mu$ with $\omega(\kappa,\kappa')=\mu(\kappa\kappa')\big/\big(\mu(\kappa)\mu(\kappa')\big)$ for all $\kappa,\kappa'\in G$.
\end{proposition}

\begin{proof}
	We have $V'_e=\mathbb{1}$ and
	\begin{equation*}
		V'_\kappa V'_{\kappa'}=\mu(\kappa)\mu(\kappa')\,\omega(\kappa,\kappa')\,V_{\kappa\kappa'}=\frac{\mu(\kappa)\mu(\kappa')}{\mu(\kappa\kappa')}\,\omega(\kappa,\kappa')\,V'_{\kappa\kappa'},
	\end{equation*}
	which is Equation~\eqref{eq:rephasing}. The new multiplier is trivial if and only if $\omega(\kappa,\kappa')=\mu(\kappa\kappa')\big/\big(\mu(\kappa)\mu(\kappa')\big)$.
\end{proof}

\noindent The freedom in the choice of representatives acts on the multiplier through the factor $\mu(\kappa)\mu(\kappa')/\mu(\kappa\kappa')$, which is the differential $d_1\mu$ of Definition~\ref{def:group-cochains} written multiplicatively, and which equals $1$ exactly when $\mu$ is a homomorphism.

In the notation of Definition~\ref{def:group-cochains} with $A=U(1)$, the two propositions state that the multiplier $\omega$ belongs to $C^2(G,U(1))$ and satisfies $d_2\omega=1$, and that a change of representatives replaces $\omega$ by $\omega\cdot d_1\mu$. This is the exact counterpart of Proposition~\ref{prop:defect-properties}, with the multiplier in the role of the defect and the rephasing in the role of the change of section. The consequence is that only the class of the multiplier can carry information.

\begin{remark}[The class of the multiplier]\label{rem:class-of-multiplier}
	Proposition~\ref{prop:multiplier-cocycle} states that the multiplier of a projective unitary representation is a $2$-cocycle on $G$ with values in $U(1)$, and Equation~\eqref{eq:rephasing} states that a change of representatives replaces it by $\omega\cdot d_1\mu$, which is cohomologous to $\omega$. It descends that the class $[\omega]\in H^2(G,U(1))$ does not depend on the operators chosen within their rays. By Proposition~\ref{prop:rephasing}, $V$ can be rephased into a genuine representation if and only if $[\omega]=0$, so that this class is the obstruction to removing the phases.
\end{remark}

%--------------------------------------------------------------------
\subsubsection*{Realizing a Cocycle: the Classification of Central Extensions}

When the class $[\omega]$ does not vanish, the phases cannot be removed, and the previous picture suggests the alternative. Rather than eliminating them, one promotes them to an additional coordinate of the group, which is the construction of a central extension. We have so far proved only one direction of the correspondence between cocycles and extensions, namely that the defect of a section is a cocycle. The converse requires building a group out of a cocycle. The construction is the associativity computation of Proposition~\ref{prop:defect-properties} run backwards, and we display it first, since it explains the shape of $d_2$.

\begin{remark}[The meaning of $d_2$]\label{rem:associator}
	Let $\zeta\in C^2(G,A)$ be an arbitrary normalized map and consider on the set $A\times G$ the operation $(a,\kappa)(a',\kappa'):=\big(a+a'+\zeta(\kappa,\kappa'),\,\kappa\kappa'\big)$, which is Equation~\eqref{eq:product-defect} read as a definition. Bracketing a triple product in the two possible ways yields
	\begin{align*}
		\big((a,\kappa)(a',\kappa')\big)(a'',\kappa'')&=\big(a+a'+a''+\zeta(\kappa,\kappa')+\zeta(\kappa\kappa',\kappa''),\,\kappa\kappa'\kappa''\big),\\
		(a,\kappa)\big((a',\kappa')(a'',\kappa'')\big)&=\big(a+a'+a''+\zeta(\kappa',\kappa'')+\zeta(\kappa,\kappa'\kappa''),\,\kappa\kappa'\kappa''\big).
	\end{align*}
	The difference between the first components of the two expressions is $-(d_2\zeta)(\kappa,\kappa',\kappa'')$, so that the operation is associative if and only if $d_2\zeta=0$. In the picture of Remark~\ref{rem:picture-reading}, the first line corresponds to the cut of the quadrilateral along the diagonal from $g_0$ to $g_2$ and the second to the cut along the diagonal from $g_1$ to $g_3$. Each of the four terms of $d_2\zeta$ comes from one of the four triangles which appear in the two cuts.
\end{remark}

The statement below collects the converse, the comparison between an extension and its defect, and the classification. Its last item is the dictionary with unitary representations, which closes the loop with the beginning of the subsection.

\begin{proposition}[Cocycles, extensions and projective representations]\label{prop:cocycle-extension}
	Let $G$ be a group, $A$ an abelian group and $\zeta$ a $2$-cocycle on $G$ with values in $A$.
	\begin{enumerate}
		\item The set $\widehat{G}_\zeta:=A\times G$ with the product
		\begin{equation*}
			(a,\kappa)(a',\kappa'):=\big(a+a'+\zeta(\kappa,\kappa'),\,\kappa\kappa'\big)
		\end{equation*}
		is a group with identity $(0,e)$. The maps $\iota(a)=(a,e)$ and $\pi(a,\kappa)=\kappa$ make it a central extension of $G$ by $A$, as per Definition~\ref{def:central-extension}.
		\item Every central extension $(\widehat{G},\iota,\pi)$ of $G$ by $A$ admits a section $\varsigma$ as per Definition~\ref{def:section}. The defect $\zeta_\varsigma$ is a cocycle, and $\Theta(a,\kappa):=\iota(a)\varsigma(\kappa)$ is an equivalence of $\widehat{G}_{\zeta_\varsigma}$ with $\widehat{G}$.
		\item $\widehat{G}_\zeta$ and $\widehat{G}_{\zeta'}$ are equivalent if and only if $\zeta$ and $\zeta'$ are cohomologous. Hence $H^2(G,A)$ classifies the central extensions of $G$ by $A$ up to equivalence.
		\item Let $\chi:A\to U(1)$ be a homomorphism and let $\Pi$ be a unitary representation of $\widehat{G}_\zeta$ with $\Pi(\iota(a))=\chi(a)\mathbb{1}$ for all $a\in A$. Then $V_\kappa:=\Pi(0,\kappa)$ is a projective unitary representation of $G$ with multiplier $\chi\circ\zeta$. Conversely, if $V$ is a projective unitary representation with multiplier $\chi\circ\zeta$, then $\Pi(a,\kappa):=\chi(a)V_\kappa$ is such a representation of $\widehat{G}_\zeta$.
	\end{enumerate}
\end{proposition}

\begin{proof}
	Item 1. By Remark~\ref{rem:associator} the product is associative, because $d_2\zeta=0$. The element $(0,e)$ is a two-sided identity because $\zeta(e,\kappa)=\zeta(\kappa,e)=0$. Moreover $(a,\kappa)\big(-a-\zeta(\kappa,\kappa^{-1}),\kappa^{-1}\big)=(0,e)$, so that every element $x$ has a right inverse $y$. Then $y$ has a right inverse $z$, and $x=x(yz)=(xy)z=z$, so that $yx=yz=(0,e)$ and $y$ is a two-sided inverse. Hence $\widehat{G}_\zeta$ is a group. The same normalization shows that $(a,e)(a',\kappa')=(a+a',\kappa')=(a',\kappa')(a,e)$, so that $\iota(A)$ is central, and $\ker\pi=\iota(A)$ is immediate.
	
	Item 2. A section exists because $\pi$ is surjective, and $\zeta_\varsigma$ is a cocycle by item 1 of Proposition~\ref{prop:defect-properties}. By Equation~\eqref{eq:product-defect},
	\begin{equation*}
		\Theta(a,\kappa)\Theta(a',\kappa')=\iota\big(a+a'+\zeta_\varsigma(\kappa,\kappa')\big)\varsigma(\kappa\kappa')=\Theta\big((a,\kappa)(a',\kappa')\big),
	\end{equation*}
	so that $\Theta$ is a homomorphism. It commutes with $\iota$ and $\pi$. It is injective, because $\iota(a)\varsigma(\kappa)=e$ implies $\kappa=e$ after applying $\pi$ and then $a=0$, and it is surjective, because $x\varsigma(\pi(x))^{-1}\in\ker\pi=\iota(A)$ for every $x\in\widehat{G}$.
	
	Item 3. An equivalence $\Theta:\widehat{G}_\zeta\to\widehat{G}_{\zeta'}$ satisfies $\Theta(a,e)=(a,e)$ and $\Theta(0,\kappa)=(\mu(\kappa),\kappa)$ for a map $\mu:G\to A$ with $\mu(e)=0$. Since $(a,\kappa)=(a,e)(0,\kappa)$, one has $\Theta(a,\kappa)=(a+\mu(\kappa),\kappa)$. Comparing $\Theta(a,\kappa)\Theta(a',\kappa')$ with $\Theta\big((a,\kappa)(a',\kappa')\big)$ shows that $\Theta$ is a homomorphism if and only if $\zeta'(\kappa,\kappa')=\zeta(\kappa,\kappa')-\mu(\kappa)-\mu(\kappa')+\mu(\kappa\kappa')$, that is, if and only if $\zeta'-\zeta=-d_1\mu$. Conversely, for every such $\mu$ the map $\Theta$ just written is an equivalence, with inverse built from $-\mu$. It follows that the assignment $[\zeta]\mapsto[\widehat{G}_\zeta]$ is well defined and injective on $H^2(G,A)$, and it is surjective by item 2.
	
	Item 4. We have $V_e=\Pi(0,e)=\mathbb{1}$ and
	\begin{equation*}
		\Pi(0,\kappa)\Pi(0,\kappa')=\Pi\big(\zeta(\kappa,\kappa'),\kappa\kappa'\big)=\Pi\big(\iota(\zeta(\kappa,\kappa'))\,(0,\kappa\kappa')\big)=\chi\big(\zeta(\kappa,\kappa')\big)\,V_{\kappa\kappa'}.
	\end{equation*}
	Conversely, $\Pi(\iota(a))=\Pi(a,e)=\chi(a)\mathbb{1}$ and
	\begin{equation*}
		\Pi(a,\kappa)\Pi(a',\kappa')=\chi(a+a')\,V_\kappa V_{\kappa'}=\chi\big(a+a'+\zeta(\kappa,\kappa')\big)\,V_{\kappa\kappa'}=\Pi\big((a,\kappa)(a',\kappa')\big).
	\end{equation*}
\end{proof}

\noindent Items 1 to 3 answer the first previous question, since the functions which occur as defects are exactly the cocycles, and the second, since the class of the defect does not depend on the section.

Item 4 is the dictionary used in the sequel. A projective representation with multiplier $\chi\circ\zeta$ is the same as a genuine representation of the extension in which the central subgroup acts by the character $\chi$, so that the phases are absorbed by the extra coordinate. Applied to a multiplier $\omega$ with $A=U(1)$, $\chi=\mathrm{id}$ and $\zeta=\omega$, which is a cocycle by Proposition~\ref{prop:multiplier-cocycle}, it yields a genuine unitary representation $\Pi(z,\kappa)=zV_\kappa$ of the extension $\widehat{G}_\omega$ of $G$ by $U(1)$. In $\widehat{G}_\omega$ the section $\kappa\mapsto(1,\kappa)$ has defect $\omega$, which completes the correspondence between the multiplier and the defect announced after Remark~\ref{rem:projective-quotient}. The sign in item 3 differs from the one in Equation~\eqref{eq:rephasing}, but this is immaterial, since the two formulas are exchanged by $\mu\mapsto-\mu$ and $\mu$ ranges over a group.

%--------------------------------------------------------------------
\subsubsection*{A Test for Nontrivial Classes: the Commutator Pairing}

A cocycle depends on the choice of the section $\varsigma$ and of the phases, hence it is difficult to inspect directly. Showing that it is not a coboundary by definition would require excluding every possible map $\mu$. There is however a quantity built from the cocycle which does not depend on any of these choices, and which provides a practical test. Its picture is the one already met in Example~\ref{ex:heisenberg-group}. If two elements $\kappa$ and $\kappa'$ commute, the two broken paths from $e$ to $\kappa\kappa'=\kappa'\kappa$, one through $\kappa$ and one through $\kappa'$, bound a parallelogram, as in Figure~\ref{fig:parallelogram}. The displacement along the fibre accumulated around its boundary is the difference between the values of the cocycle on the two triangles into which the diagonal cuts it.

\begin{figure}[ht]
	\centering
	% ---------------------------------------------------------------
	% PROMPT FOR A TIKZ GENERATOR (paste everything between the rules)
	% ---------------------------------------------------------------
	% Write TikZ code for a figure, to be compiled inside a LaTeX
	% document that already loads tikz. Use only the libraries
	% arrows.meta, calc and positioning. Black and white (light gray
	% fills allowed), line width 0.8pt, final width about 0.5\textwidth,
	% labels in math mode with the standard LaTeX math font.
	% Content: a parallelogram with vertices $e$ at (0,0), $\kappa$ at
	% (3,-0.6), $\kappa'$ at (0.8,2) and $\kappa\kappa'=\kappa'\kappa$
	% at (3.8,1.4). Mark the vertices with small filled dots and write
	% their names outside the figure, next to the vertices. Draw the
	% four sides as arrows with a closed arrow tip (Stealth): the lower
	% path from $e$ to $\kappa$ and from $\kappa$ to $\kappa\kappa'$,
	% labelled $\kappa$ and $\kappa'$, and the upper path from $e$ to
	% $\kappa'$ and from $\kappa'$ to $\kappa\kappa'$, labelled $\kappa'$
	% and $\kappa$. Place the labels on the outer side of the sides.
	% Draw the diagonal from $e$ to $\kappa\kappa'$ as a dashed arrow.
	% Fill the lower triangle (e, kappa, kappa kappa') with a light gray
	% and write inside it $\zeta(\kappa,\kappa')$, and fill the upper
	% triangle (e, kappa', kappa kappa') with a lighter gray and write
	% inside it $\zeta(\kappa',\kappa)$. Under the figure write
	% "$\Omega_\zeta(\kappa,\kappa')=\zeta(\kappa,\kappa')-\zeta(\kappa',\kappa)$".
	% ---------------------------------------------------------------
	\fbox{\parbox{0.5\textwidth}{\centering\vspace{3cm}TODO: parallelogram figure\vspace{3cm}}}
	\caption{Two commuting elements $\kappa$ and $\kappa'$ and the two broken paths from $e$ to $\kappa\kappa'$. The pairing $\Omega_\zeta(\kappa,\kappa')$ is the difference between the values of the cocycle on the two triangles, and in $\mathbb{R}^2$ it is twice the area of the parallelogram.}
	\label{fig:parallelogram}
\end{figure}

\begin{proposition}[Commutator pairing]\label{prop:commutator-pairing}
	Let $\zeta$ be a $2$-cocycle on $G$ with values in $A$, and let $\kappa,\kappa'\in G$ satisfy $\kappa\kappa'=\kappa'\kappa$. Set $\Omega_\zeta(\kappa,\kappa'):=\zeta(\kappa,\kappa')-\zeta(\kappa',\kappa)\in A$. Then
	\begin{enumerate}
		\item in $\widehat{G}_\zeta$, any lifts $\hat\kappa,\hat\kappa'$ of $\kappa,\kappa'$ satisfy $\hat\kappa\hat\kappa'=\iota\big(\Omega_\zeta(\kappa,\kappa')\big)\hat\kappa'\hat\kappa$,
		\item if $\zeta'$ is cohomologous to $\zeta$, then $\Omega_{\zeta'}(\kappa,\kappa')=\Omega_\zeta(\kappa,\kappa')$ for every commuting pair. In particular $\Omega_\zeta$ vanishes on all commuting pairs if $\zeta$ is a coboundary,
		\item if $V$ is a projective unitary representation with multiplier $\omega$, then $V_\kappa V_{\kappa'}=\Omega_\omega(\kappa,\kappa')\,V_{\kappa'}V_\kappa$, where $\Omega_\omega(\kappa,\kappa'):=\omega(\kappa,\kappa')\,\omega(\kappa',\kappa)^{-1}$ is the pairing of item 1 written multiplicatively for $A=U(1)$.
	\end{enumerate}
\end{proposition}

\begin{proof}
	Item 1. For the lifts $(0,\kappa)$ and $(0,\kappa')$ one has $(0,\kappa)(0,\kappa')=(\zeta(\kappa,\kappa'),\kappa\kappa')$ and $(0,\kappa')(0,\kappa)=(\zeta(\kappa',\kappa),\kappa\kappa')$, which yields the claim. Replacing a lift by another one multiplies it by a central element, which does not affect the relation. Item 2. If $\zeta'=\zeta+d_1\mu$, then $\Omega_{\zeta'}(\kappa,\kappa')-\Omega_\zeta(\kappa,\kappa')=-\mu(\kappa\kappa')+\mu(\kappa'\kappa)=0$, because $\kappa\kappa'=\kappa'\kappa$. Item 3. We have $V_\kappa V_{\kappa'}=\omega(\kappa,\kappa')V_{\kappa\kappa'}$ and $V_{\kappa'}V_\kappa=\omega(\kappa',\kappa)V_{\kappa'\kappa}=\omega(\kappa',\kappa)V_{\kappa\kappa'}$, and eliminating $V_{\kappa\kappa'}$ yields the claim.
\end{proof}

\noindent The pairing $\Omega_\zeta$ is the precise sense in which a cocycle measures how wrongly the lifted operators commute. If it is nonzero for one pair of commuting elements, no choice of phases can make the representation genuine, and $\zeta$ is not a coboundary.

In the guiding example, Example~\ref{ex:area-defect} yields $\Omega_\zeta(\mathbf{u},\mathbf{u}')=\mathbf{u}\wedge\mathbf{u}'$, which is twice the area of the parallelogram spanned by $\mathbf{u}$ and $\mathbf{u}'$ and is nonzero as soon as the two vectors are not parallel. Proposition~\ref{prop:commutator-pairing} therefore reproduces by a short test the conclusion which Example~\ref{ex:area-defect} reached by a direct argument, and it is the test that we apply to the Galilei group in Proposition~\ref{prop:bargmann-nontrivial}. The following finite example shows all the objects of the subsection at work, without any differentiable structure and without an area to compute. It also anticipates, in the simplest possible setting, the mechanism by which boosts and translations will produce the mass.

\begin{example}[The Klein group and the Pauli matrices]\label{ex:klein-pauli}
	Let $G=\{e,a,b,ab\}$ be the Klein four-group, \textit{i.e.}, the abelian group generated by two elements $a$ and $b$ of order two, and let $\mathcal{H}=\mathbb{C}^2$. Define $V_e:=\mathbb{1}$, $V_a:=\sigma_x$, $V_b:=\sigma_z$ and $V_{ab}:=\sigma_x\sigma_z$, where $\sigma_x$ and $\sigma_z$ are the Pauli matrices. Since $\sigma_x^2=\sigma_z^2=\mathbb{1}$ and $\sigma_z\sigma_x=-\sigma_x\sigma_z$, any product $V_\kappa V_{\kappa'}$ is a phase times $V_{\kappa\kappa'}$, so that $V$ is a projective unitary representation. Its multiplier satisfies $\omega(a,b)=1$ and $\omega(b,a)=-1$, hence $\Omega_\omega(a,b)=-1$. Since $a$ and $b$ commute, item 2 of Proposition~\ref{prop:commutator-pairing} entails that $\omega$ is not a coboundary. This is also visible directly, because in a genuine representation of an abelian group the operators $V_a$ and $V_b$ would commute, whereas a rephasing multiplies each of them by a phase and cannot change the relation $V_aV_b=-V_bV_a$.
	
	Equation~\eqref{eq:rephasing} can be tested on the same data. Choosing $\mu(a)=i$ and $\mu(b)=\mu(ab)=1$ yields $V'_a=i\sigma_x$ and $V'_b=V_b$, with multiplier $\omega'(a,a)=\omega(a,a)\,\mu(a)^2=-1$, $\omega'(a,b)=i$ and $\omega'(b,a)=-i$. The multiplier has changed at each of these pairs, as it must since $d_1\mu\neq1$, but $\Omega_{\omega'}(a,b)=i/(-i)=-1$ coincides with $\Omega_\omega(a,b)$, in accordance with Proposition~\ref{prop:commutator-pairing}. Example~\ref{ex:heisenberg-extension} below exhibits the same mechanism for the continuous group $\mathbb{R}^2$.
\end{example}
\subsection{Lie Algebra Cohomology and Invariant Forms}\label{subsec:lie-cohomology}

For Lie groups the classification of multipliers can be transferred to the Lie algebra, where the cocycle identity becomes a linear condition. In this subsection $\mathfrak{g}$ is the finite-dimensional real Lie algebra of a Lie group $G$ as per Definition~1.2.10. Its cochains are the alternating multilinear forms on $\mathfrak{g}$, and the differential is built from the bracket alone.

\begin{definition}[Chevalley-Eilenberg complex]\label{def:lie-cohomology-group}
	Let $\mathfrak{g}$ be a real Lie algebra as per Definition~1.2.1 and let $C^k(\mathfrak{g}):=\mathrm{Alt}^k(\mathfrak{g},\mathbb{R})$ be the space of alternating $k$-linear forms on $\mathfrak{g}$, with $C^0(\mathfrak{g}):=\mathbb{R}$. The \emph{Chevalley-Eilenberg differential} $d_k:C^k(\mathfrak{g})\to C^{k+1}(\mathfrak{g})$ is
	\begin{equation*}
		(d_k\psi)(X_0,\dots,X_k):=\sum_{0\leq i<j\leq k}(-1)^{i+j}\,\psi\big([X_i,X_j],X_0,\dots,\widehat{X_i},\dots,\widehat{X_j},\dots,X_k\big),
	\end{equation*}
	where a hat denotes an omitted argument. The cohomology groups of this complex are denoted by $H^k(\mathfrak{g},\mathbb{R})$.
\end{definition}

\noindent That $d_{k+1}\circ d_k=0$, which is required by Definition~\ref{def:cochain-complex}, is proved in Corollary~\ref{cor:lie-complex} below. In degrees one and two the formula can be evaluated at once, and it recovers the cocycle identity that is familiar from the group case.

\begin{proposition}[Degrees one and two]\label{prop:lie-degree-two}
	For $\lambda\in C^1(\mathfrak{g})$ and $\psi\in C^2(\mathfrak{g})$ one has
	\begin{align*}
		(d_1\lambda)(X,Y)&=-\lambda([X,Y]),\\
		(d_2\psi)(X,Y,Z)&=-\big(\psi([X,Y],Z)+\psi([Y,Z],X)+\psi([Z,X],Y)\big).
	\end{align*}
	Hence a $2$-cocycle of $\mathfrak{g}$ is a bilinear antisymmetric map $\psi:\mathfrak{g}\times\mathfrak{g}\to\mathbb{R}$ with $\psi([X,Y],Z)+\psi([Y,Z],X)+\psi([Z,X],Y)=0$, the coboundaries are the maps $\psi_\lambda(X,Y):=\lambda([X,Y])$ with $\lambda\in\mathfrak{g}^*$, and $H^2(\mathfrak{g},\mathbb{R})$ is the quotient of the former by the latter.
\end{proposition}

\begin{proof}
	The first formula is the case $k=1$ of Definition~\ref{def:lie-cohomology-group}. For $k=2$ the pairs $(i,j)=(0,1),(0,2),(1,2)$ contribute $-\psi([X,Y],Z)$, $+\psi([X,Z],Y)$ and $-\psi([Y,Z],X)$ respectively. The antisymmetry of the bracket and of $\psi$ gives $\psi([X,Z],Y)=-\psi([Z,X],Y)$. The sign in $d_1\lambda=-\psi_\lambda$ is immaterial because $\lambda$ ranges over a vector space. Finally $d_2d_1\lambda=0$ is the Jacobi identity applied to $\lambda$.
\end{proof}

The expression of $d_k$ in Definition~\ref{def:lie-cohomology-group} coincides with the invariant formula for the exterior derivative, in which the terms with derivatives of functions are missing. This is because the forms entering are constant along a frame of left-invariant vector fields, and it connects Lie algebra cohomology to the de Rham complex of the group.

\begin{definition}[Left-invariant form]\label{def:left-invariant-form}
	Let $G$ be a Lie group as per Definition~1.2.8. A form $\alpha\in\Omega^k(G)$ is \emph{left-invariant} if $L_\kappa^*\alpha=\alpha$ for all $\kappa\in G$, where $L_\kappa$ is the left multiplication and the pullback is as per Definition~1.1.26. The space of left-invariant $k$-forms is denoted by $\Omega^k(G)^G$.
\end{definition}

\begin{proposition}[Left-invariant forms and the Chevalley-Eilenberg complex]\label{prop:invariant-forms}
	Let $G$ be a Lie group with Lie algebra $\mathfrak{g}$. The map $\alpha\mapsto\alpha_e$, which evaluates a left-invariant form at the identity, is a linear isomorphism $\Omega^k(G)^G\to C^k(\mathfrak{g})$. Moreover $d$ maps $\Omega^k(G)^G$ into $\Omega^{k+1}(G)^G$ and $(d\alpha)_e=d_k(\alpha_e)$.
\end{proposition}

\begin{proof}
	Left invariance means $\alpha_\kappa(dL_\kappa v_1,\dots,dL_\kappa v_k)=\alpha_e(v_1,\dots,v_k)$, so $\alpha$ is determined by $\alpha_e$, and the map is injective. Conversely, for $\psi\in C^k(\mathfrak{g})$ the formula $\alpha_\kappa(w_1,\dots,w_k):=\psi(dL_{\kappa^{-1}}w_1,\dots,dL_{\kappa^{-1}}w_k)$ defines a left-invariant form, because $L_{(\kappa'\kappa)^{-1}}\circ L_{\kappa'}=L_{\kappa^{-1}}$. It is smooth because its components with respect to the global frame of left-invariant vector fields $X^{e_1},\dots,X^{e_n}$ of Definition~1.2.10 are constant. Hence the map is surjective.
	
	Since $d$ commutes with pullbacks, $L_\kappa^*d\alpha=dL_\kappa^*\alpha=d\alpha$, so $d\alpha$ is left-invariant. Let $v_0,\dots,v_k\in\mathfrak{g}$. The invariant formula for the exterior derivative \cite[TODO-location]{TODO-Lee13} applied to $X^{v_0},\dots,X^{v_k}$ reads
	\begin{equation*}
		(d\alpha)(X^{v_0},\dots,X^{v_k})=\sum_i(-1)^iX^{v_i}\big(\alpha(\dots,\widehat{X^{v_i}},\dots)\big)+\sum_{i<j}(-1)^{i+j}\alpha\big([X^{v_i},X^{v_j}],\dots,\widehat{X^{v_i}},\dots,\widehat{X^{v_j}},\dots\big).
	\end{equation*}
	In the first sum the functions $\alpha(X^{v_0},\dots)$ are constant by left invariance, so it vanishes. In the second sum $[X^{v_i},X^{v_j}]=X^{[v_i,v_j]}$ by Definition~1.2.10. Evaluating at $e$ yields $(d\alpha)_e(v_0,\dots,v_k)=(d_k\alpha_e)(v_0,\dots,v_k)$.
\end{proof}

\begin{corollary}\label{cor:lie-complex}
	The sequence $(C^k(\mathfrak{g}),d_k)$ of Definition~\ref{def:lie-cohomology-group} satisfies $d_{k+1}\circ d_k=0$, hence it is a cochain complex as per Definition~\ref{def:cochain-complex}. Moreover $H^k(\mathfrak{g},\mathbb{R})$ is isomorphic to the cohomology of the complex of left-invariant forms $(\Omega^\bullet(G)^G,d)$, and in degree two the $2$-cocycles of $\mathfrak{g}$ are the values at $e$ of the closed left-invariant $2$-forms, while the coboundaries are the values at $e$ of the forms $d\lambda$ with $\lambda\in\Omega^1(G)^G$.
\end{corollary}

\begin{proof}
	By Proposition~\ref{prop:invariant-forms}, $d_{k+1}d_k\alpha_e=(dd\alpha)_e=0$ by item iii of Definition~1.1.24. The remaining statements are the translation of the same proposition through the isomorphism $\alpha\mapsto\alpha_e$.
\end{proof}

\noindent In this sense the cocycle condition of Proposition~\ref{prop:lie-degree-two} is the closedness of a left-invariant $2$-form, exactly as the closedness of the symplectic form in Definition~1.3.8.

\begin{remark}[Invariant and de Rham cohomology]\label{rem:chevalley-eilenberg}
	The complex of left-invariant forms is a subcomplex of the de Rham complex, and the induced map $H^k(\mathfrak{g},\mathbb{R})\to H^k_{\mathrm{dR}}(G)$ is an isomorphism if $G$ is compact and connected \cite[TODO-location]{TODO-ChevalleyEilenberg1948}. It is not so in general. For $G=\mathbb{R}^2$, the form $dx\wedge dy$ is closed and left-invariant, and it is exact in the sense of de Rham since $dx\wedge dy=d(x\,dy)$. It is not the differential of any left-invariant form, because such forms have constant coefficients, hence they are closed. Thus $H^2(\mathbb{R}^2,\mathbb{R})\neq0$ while $H^2_{\mathrm{dR}}(\mathbb{R}^2)=0$, in agreement with Example~\ref{ex:heisenberg-extension} below.
\end{remark}

A cocycle allows us to build a larger Lie algebra, in analogy with Proposition~\ref{prop:cocycle-extension}. The following statement is the infinitesimal counterpart of the correspondence between $H^2$ and central extensions.

\begin{proposition}[Central extension of a Lie algebra]\label{prop:lie-extension}
	Let $\psi$ be a $2$-cocycle on $\mathfrak{g}$. The vector space $\mathfrak{g}_\psi:=\mathbb{R}\oplus\mathfrak{g}$ with the bracket $[(a,X),(b,Y)]:=\big(\psi(X,Y),[X,Y]\big)$ is a Lie algebra in which $\mathbb{R}\oplus\{0\}$ is central, and $\mathfrak{g}_\psi$ and $\mathfrak{g}_{\psi'}$ are equivalent extensions of $\mathfrak{g}$ if and only if $[\psi]=[\psi']$ in $H^2(\mathfrak{g},\mathbb{R})$.
\end{proposition}

\begin{proof}
	The bracket is bilinear and antisymmetric, and $[(a,0),(b,Y)]=(\psi(0,Y),0)=0$. For the Jacobi identity,
	\begin{equation*}
		\big[[(a,X),(b,Y)],(c,Z)\big]=\big(\psi([X,Y],Z),[[X,Y],Z]\big),
	\end{equation*}
	and the sum over cyclic permutations of $(X,Y,Z)$ vanishes by the Jacobi identity of $\mathfrak{g}$ and by the cocycle condition. An equivalence commuting with the inclusion of $\mathbb{R}$ and with the projection has the form $(a,X)\mapsto(a+\lambda(X),X)$ for a linear functional $\lambda$, and it is a Lie algebra homomorphism if and only if $\psi'(X,Y)=\psi(X,Y)+\lambda([X,Y])$, that is, if and only if $\psi'-\psi$ is a coboundary.
\end{proof}

%--------------------------------------------------------------------
\subsection{A Cocycle in Classical Mechanics}\label{subsec:classical-cocycle}

Cocycles are not an exclusively quantum phenomenon, and they already appear in the Hamiltonian formalism of Chapter 1. The momentum map of Definition~1.5.7 is required to satisfy $d\hat J(\xi)=\iota_{\xi_P}\omega$ only, and this determines each function $\hat J(\xi)$ up to an additive constant. The Lie algebra structure of the group is not required to be preserved by $J$, and the failure of this property is measured by a cocycle.

\begin{proposition}[Defect of the momentum map]\label{prop:momentum-defect}
	Let $(P,\omega,\Phi,J)$ be a Hamiltonian $G$-space as per Definition~1.5.7 with $P$ connected, and let $\mathfrak{g}$ be the Lie algebra of $G$. For $\xi,\eta\in\mathfrak{g}$ the function
	\begin{equation*}
		\{\hat J(\xi),\hat J(\eta)\}-\hat J([\xi,\eta])
	\end{equation*}
	is constant on $P$, where $\{\cdot,\cdot\}$ is the Poisson bracket of Definition~1.4.7. The constant, denoted by $\Sigma(\xi,\eta)$, defines a $2$-cocycle of $\mathfrak{g}$. If $\Sigma=\psi_\lambda$ for some $\lambda\in\mathfrak{g}^*$, then $J':=J+\lambda$ is a momentum map for the same action with $\Sigma'=0$.
\end{proposition}

\begin{proof}
	By the proof of Proposition~1.4.12, $X_{\{f,h\}}=-[X_f,X_h]$ for $f,h\in C^\infty(P)$. Using Remark~1.5.8 and Equation~(1.7),
	\begin{equation*}
		X_{\{\hat J(\xi),\hat J(\eta)\}}=-[\xi_P,\eta_P]=[\xi,\eta]_P=X_{\hat J([\xi,\eta])}.
	\end{equation*}
	Since $\iota_{X_H}\omega=dH$ and $\omega$ is non-degenerate, the function $\{\hat J(\xi),\hat J(\eta)\}-\hat J([\xi,\eta])$ has vanishing differential, and it is constant by Proposition~\ref{prop:H0}. The map $\Sigma$ is bilinear and antisymmetric. Writing $\hat J([\xi,\eta])=\{\hat J(\xi),\hat J(\eta)\}-\Sigma(\xi,\eta)$ and noting that the Poisson bracket between a constant and any function vanishes,
	\begin{equation*}
		\Sigma([\xi,\eta],\zeta)=\{\{\hat J(\xi),\hat J(\eta)\},\hat J(\zeta)\}-\hat J([[\xi,\eta],\zeta]).
	\end{equation*}
	The sum over cyclic permutations of the right-hand side vanishes by the Jacobi identity of Proposition~1.4.12 and of $\mathfrak{g}$, so $\Sigma$ satisfies the cocycle condition of Proposition~\ref{prop:lie-degree-two}. Finally, adding the constant $\lambda(\xi)$ to $\hat J(\xi)$ does not change $d\hat J(\xi)$ nor the Poisson brackets, and it replaces $\Sigma$ by $\Sigma-\psi_\lambda$.
\end{proof}

\noindent The class $[\Sigma]\in H^2(\mathfrak{g},\mathbb{R})$ is thus the obstruction to finding a momentum map which preserves brackets, and an explicit instance is computed in Example~\ref{ex:heisenberg-extension}.

%--------------------------------------------------------------------
\subsection{From the Group to the Lie Algebra: Bargmann's Theorem}\label{subsec:bargmann}

We have now two cohomologies attached to a Lie group $G$, the one of the group with values in $U(1)$ or $\mathbb{R}$, which involves arbitrary functions, and the one of its Lie algebra, which is finite-dimensional linear algebra. The bridge between the two is obtained by differentiating a cocycle, which is the content of the following definition. It reproduces at the infinitesimal level the commutator pairing of Proposition~\ref{prop:commutator-pairing}.

\begin{definition}[Infinitesimal cocycle]\label{def:infinitesimal-cocycle}
	Let $G$ be a Lie group as per Definition~1.2.8, with Lie algebra $\mathfrak{g}$ and exponential map as per Definition~1.2.13, and let $\zeta$ be a real cocycle, defined and smooth on a neighbourhood of $(e,e)$. The \emph{infinitesimal cocycle} of $\zeta$ is the bilinear map
	\begin{equation*}
		\psi_\zeta(X,Y):=\frac{\partial^2}{\partial s\,\partial t}\Big[\zeta(\exp sX,\exp tY)-\zeta(\exp tY,\exp sX)\Big]_{s=t=0}, \qquad X,Y\in\mathfrak{g}.
	\end{equation*}
\end{definition}

\begin{remark}
	If $[X,Y]=0$, then $\exp sX$ and $\exp tY$ commute for all $s,t$, and $$\psi_\zeta(X,Y)=\partial_s\partial_t\,\Omega_\zeta(\exp sX,\exp tY)|_{s=t=0}.$$ Thus $\psi_\zeta$ detects, at the infinitesimal level, the same failure of commutation as $\Omega_\zeta$.
\end{remark}

We can now state the result which reduces the classification of projective representations of a simply connected Lie group to a computation in its Lie algebra. The statement is the form of Bargmann's theorem used in the thesis, and it is recalled without proof. As explained in Remark~\ref{rem:topological-groups}, multipliers are continuous and rephasing functions are Borel maps.

\begin{proposition}[Bargmann]\label{prop:bargmann-correspondence}
	Let $G$ be a connected and simply connected Lie group with Lie algebra $\mathfrak{g}$, and let $\H$ be a separable Hilbert space.
	\begin{enumerate}
		\item Every continuous multiplier $\omega$ of $G$ has a local exponent, \textit{i.e.}, a real cocycle $\zeta$, smooth in a neighbourhood of $(e,e)$, with $\omega=e^{i\zeta}$ there. The infinitesimal cocycle $\psi_\zeta$ of Definition~\ref{def:infinitesimal-cocycle} is a $2$-cocycle of $\mathfrak{g}$, and its class in $H^2(\mathfrak{g},\mathbb{R})$ depends only on the equivalence class of $\omega$, where $\omega$ and $\omega'$ are equivalent multipliers if $\omega'\omega^{-1}$ is the coboundary of a Borel map $G\to U(1)$.
		\item The assignment in item 1 is a bijection from the equivalence classes of continuous multipliers of $G$ onto $H^2(\mathfrak{g},\mathbb{R})$.
		\item Let $V$ be a continuous projective unitary representation of $G$ on $\H$ with multiplier $\omega$, and let $\zeta$ be a continuous real cocycle on all of $G$ such that $[\psi_\zeta]$ is the class of $\omega$. Then the phases of $V$ can be chosen so that $V$ is strongly continuous with multiplier exactly $e^{i\zeta}$. By Proposition~\ref{prop:cocycle-extension}, $\Pi(a,\kappa):=e^{ia}V_\kappa$ is then a strongly continuous unitary representation of the central extension $\widehat{G}_\zeta$.
	\end{enumerate}
\end{proposition}

\noindent See \cite[TODO-location]{TODO-Bargmann1954} and \cite[TODO-location]{TODO-Varadarajan1985}. The point of the proposition is that the cohomology of the group is replaced by the cohomology of the Lie algebra, which is finite-dimensional linear algebra.

The proposition packages four ideas which are worth separating. We read it step by step, since the same pattern is followed whenever the theorem is applied.

\begin{remark}[Reading Proposition~\ref{prop:bargmann-correspondence} in four steps]\label{rem:bargmann-reading}
	\emph{Step 1, the multiplier.} A symmetry of rays is lifted to unitary operators $V_\kappa$, defined up to phase, and $V_\kappa V_{\kappa'}$ and $V_{\kappa\kappa'}$ differ by a phase $\omega(\kappa,\kappa')$, which is a cocycle by Proposition~\ref{prop:multiplier-cocycle}.
	
	\emph{Step 2, the class.} Changing the phases of the lifts changes $\omega$ by a coboundary, by Proposition~\ref{prop:rephasing}. Only the class of $\omega$ is meaningful. If the class is trivial the lifts can be chosen to form a genuine representation, and otherwise the phase is an obstruction that no choice removes.
	
	\emph{Step 3, the Lie algebra.} For a connected and simply connected group the class is determined by its infinitesimal version. Differentiating a projective representation yields operators $X_a$ with $[X_a,X_b]=if_{ab}^{\ \ c}X_c+i\psi(e_a,e_b)\mathbb{1}$, where $f_{ab}^{\ \ c}$ are the structure constants and $\psi$ is a bilinear antisymmetric form obeying the cocycle condition of Proposition~\ref{prop:lie-degree-two}. Redefining $X_a\to X_a+\lambda_a\mathbb{1}$ changes $\psi$ by the coboundary $\psi_\lambda$, so that only $[\psi]\in H^2(\mathfrak{g},\mathbb{R})$ matters.
	
	\emph{Step 4, the lift.} Once the class is known, Proposition~\ref{prop:cocycle-extension} converts the projective representation into a genuine representation of a central extension, in which the phases have become values of a character of the central subgroup. The problem is thereby reduced to the linear algebra of $H^2(\mathfrak{g},\mathbb{R})$. The assumption that $G$ is simply connected is essential, as Example~\ref{ex:so3-su2} shows.
\end{remark}

%--------------------------------------------------------------------
\subsection{Examples}\label{subsec:cohomology-examples}

We conclude with three examples showing the three possible behaviours of the second cohomology group. They are the abelian algebra with a nontrivial class, which is also the origin of the canonical commutation relations, a semisimple algebra with a trivial class, and a group where the obstruction is not algebraic but topological.

\begin{example}[Heisenberg extension of $\mathbb{R}^2$]\label{ex:heisenberg-extension}
	Let $\mathfrak{g}=\mathbb{R}^2$ with the zero bracket. The cocycle condition is empty, so every bilinear antisymmetric form is a cocycle, while every coboundary vanishes by Proposition~\ref{prop:lie-degree-two}. Hence $H^2(\mathbb{R}^2,\mathbb{R})\cong\mathbb{R}$, generated by $\psi(\mathbf{u},\mathbf{u}')=u_1u_2'-u_2u_1'$, which is the constant symplectic form $dq\wedge dp$ evaluated on $\mathbf{u}$ and $\mathbf{u}'$. The extension $\mathfrak{g}_\psi$ has a basis $q,p,z$ with $[q,p]=z$ and $z$ central, the Heisenberg algebra of Definition~2.4.1 up to the factor $i$ in Equation~(2.37).
	
	\emph{Group level.} Let $G=\mathbb{R}^2$ and $\zeta(\mathbf{u},\mathbf{u}')=\tfrac{1}{2}\psi(\mathbf{u},\mathbf{u}')$. This is a real cocycle because it is bilinear, and $\Omega_\zeta(\mathbf{u},\mathbf{u}')=\psi(\mathbf{u},\mathbf{u}')$. By Proposition~\ref{prop:commutator-pairing}, $c\zeta$ is not a coboundary for $c\neq0$, and a projective representation with multiplier $e^{ic\zeta}$ consists of unitary groups $U(s)=V_{(s,0)}$ and $W(t)=V_{(0,t)}$ with $U(s)W(t)=e^{icst}\,W(t)U(s)$, a relation which cannot be removed by rephasing. For $c=-1$ this is Equation~(2.44). More precisely, if $U(\alpha),V(\beta)$ satisfy Equations~(2.44) and~(2.45), then $V_{(\alpha,\beta)}:=e^{i\alpha\beta/2}U(\alpha)V(\beta)$ is a projective unitary representation of $\mathbb{R}^2$ with multiplier $e^{-i\zeta}$, as one checks by reordering with $V(\beta)U(\alpha')=U(\alpha')V(\beta)e^{i\alpha'\beta}$. The group law of the central extension $\widehat{G}_{-\zeta}$ of Proposition~\ref{prop:cocycle-extension} is that of Equation~(2.48), so $\widehat{G}_{-\zeta}$ is the Heisenberg Lie group of Definition~2.4.14. By item 4 of Proposition~\ref{prop:cocycle-extension}, with $\chi(a)=e^{ia}$, the representations of Remark~2.4.15 are the projective representations of $\mathbb{R}^2$ with multiplier $e^{-i\zeta}$.
	
	\emph{Classical origin.} Let $P=\mathbb{R}^2$ with $\omega=dq\wedge dp$ and let $G=\mathbb{R}^2$ act by translations in both coordinates. For $\xi=(\xi^1,\xi^2)$ one has $\xi_P=\xi^1\partial_q+\xi^2\partial_p$ and $\iota_{\xi_P}\omega=\xi^1dp-\xi^2dq=d(\xi^1p-\xi^2q)$, so that $\hat J(\xi)=\xi^1p-\xi^2q$. Corollary~1.4.9 yields $\{\hat J(\xi),\hat J(\eta)\}=\xi^1\eta^2-\xi^2\eta^1=\psi(\xi,\eta)$. Since $\mathfrak{g}$ is abelian, Proposition~\ref{prop:momentum-defect} gives $\Sigma=\psi$, which is a nonzero class. Thus no momentum map for this action preserves brackets, and the functions $\hat J(e_1)=p$, $\hat J(e_2)=-q$ and $1$ span the Heisenberg algebra under the Poisson bracket, since $\{q,p\}=1$ by Remark~1.4.10. In contrast, for the translations in the $q$-coordinates alone of Example~1.5.10 one has $\{p_i,p_j\}=0$, hence $\Sigma=0$.
\end{example}

\begin{example}[The algebra $\mathfrak{su}(2)$]\label{ex:su2-cohomology}
	For the semisimple algebra $\mathfrak{su}(2)$ of Example~1.2.5, Whitehead's second lemma yields $H^2(\mathfrak{su}(2),\mathbb{R})=0$ \cite[TODO-location]{TODO-Hum10}. Since $SU(2)$ is connected and simply connected, Proposition~\ref{prop:bargmann-correspondence} shows that every continuous projective unitary representation of $SU(2)$ can be rephased into a genuine strongly continuous representation. This is the reason why the representation $\pi_{1/2}$ of Section~2.7 is a genuine representation of the group $SU(2)$.
\end{example}

\begin{example}[$SO(3)$ and $SU(2)$]\label{ex:so3-su2}
	Let $d:SU(2)\to U(\mathbb{C}^2)$ be the defining representation and let $u\mapsto R_u$ be the covering homomorphism $SU(2)\to SO(3)$, whose kernel is $\{\pm1\}$. Choose a map $\varrho:SO(3)\to SU(2)$ with $R_{\varrho(R)}=R$ and $\varrho(\mathbb{1})=1$, and set $V_R:=d(\varrho(R))$. Since $\varrho(R)\varrho(R')$ and $\varrho(RR')$ lie over the same rotation, they differ at most by the sign $-1$, so that $V$ is a projective unitary representation of $SO(3)$ with multiplier $\omega(R,R')\in\{\pm1\}$. This multiplier cannot be removed by a rephasing which preserves continuity, since a loop of rotations by an angle going from $0$ to $2\pi$ is represented by a path from $\mathbb{1}$ to $-\mathbb{1}$ \cite[TODO-location]{TODO-Hal03}. The obstruction has topological origin, because $\pi_1(SO(3))=\mathbb{Z}_2$ while $SU(2)$ is simply connected, and it disappears when $SO(3)$ is replaced by its universal covering $SU(2)$, in accordance with Example~\ref{ex:su2-cohomology}.
\end{example}

% =====================================================================
\section{The Galilei Group and its Central Extension}
\label{sec:galilei}

Galilean invariance expresses the equivalence of inertial observers whose frames differ by a time translation, a spatial translation, a rotation or a boost, that is, a change to a frame moving with constant velocity. These transformations form a Lie group, and its action on the states of a quantum system is the structure from which the Hamiltonian will be recovered. Two features distinguish the Galilei group from the symmetry groups met so far. The first is that boosts do not commute with time translations, and this noncommutativity is what ties the Hamiltonian to the momentum. The second is that the projective representations of the group cannot be reduced to genuine ones, since a central extension is unavoidable. We therefore introduce the group together with its central extension, and we record the relations that hold in the extended group, because the derivation of the Hamiltonian uses them in the form of relations between unitary operators.

\begin{definition}[Universal covering of the Galilei group]\label{def:galilei-group}
	Let $R_u\in SO(3)$ denote the rotation associated with $u\in SU(2)$ through the covering homomorphism $SU(2)\to SO(3)$. The \emph{Galilei group} $\mathcal{G}$ is the Lie group, as per Definition~\ref{def:lie-group}, with underlying manifold $\mathbb{R}\times\mathbb{R}^3\times\mathbb{R}^3\times SU(2)$ and product
	\begin{equation}
		(s,\mathbf{a},\mathbf{v},u)(s',\mathbf{a}',\mathbf{v}',u') := \big(s+s',\ \mathbf{a}+R_u\mathbf{a}'+s'\mathbf{v},\ \mathbf{v}+R_u\mathbf{v}',\ uu'\big).
	\end{equation}
\end{definition}

\begin{remark}
	The product of Definition~\ref{def:galilei-group} is the composition of the transformations $(t,\mathbf{x})\mapsto(t+s,\,R_u\mathbf{x}+\mathbf{v}t+\mathbf{a})$ of the events of a Newtonian spacetime, in which $s$ is a time translation, $\mathbf{a}$ a spatial translation, $\mathbf{v}$ the velocity of a boost and $u$ a rotation. The manifold of $\mathcal{G}$ is $\mathbb{R}^9\times S^3$, hence $\mathcal{G}$ is a connected and simply connected Lie group of dimension $10$. Quotienting by the subgroup $\{\pm 1\}\subset SU(2)$ yields the usual Galilei group.
\end{remark}

The four families of transformations just mentioned are the building blocks of every element of $\mathcal{G}$, and they will be handled separately throughout the chapter.

\begin{definition}[Elementary transformations]\label{def:elementary-transformations}
	Let $\mathcal{G}$ be the Galilei group as per Definition~\ref{def:galilei-group}. For $s\in\mathbb{R}$, $\mathbf{a},\mathbf{v}\in\mathbb{R}^3$ and $u\in SU(2)$ we set
	\begin{equation}
		\tau_s := (s,\mathbf{0},\mathbf{0},1), \qquad \alpha_{\mathbf{a}} := (0,\mathbf{a},\mathbf{0},1), \qquad \beta_{\mathbf{v}} := (0,\mathbf{0},\mathbf{v},1), \qquad \rho_u := (0,\mathbf{0},\mathbf{0},u),
	\end{equation}
	and we call them the time translation, the spatial translation, the boost and the rotation.
\end{definition}

The product of two elements of $\mathcal{G}$ is a group law, whereas a unitary representation of the projective action will satisfy the same law only up to phases. Those phases are encoded by a real function on $\mathcal{G}\times\mathcal{G}$, which we define now and then show to be a cocycle.

\begin{definition}[Bargmann cocycle]\label{def:bargmann-cocycle}
	The \emph{Bargmann cocycle} is the map $\zeta:\mathcal{G}\times\mathcal{G}\to\mathbb{R}$ defined, for $\kappa=(s,\mathbf{a},\mathbf{v},u)$ and $\kappa'=(s',\mathbf{a}',\mathbf{v}',u')$, by
	\begin{equation}
		\zeta(\kappa,\kappa') := \mathbf{v}\cdot R_u\mathbf{a}' + \tfrac{1}{2}|\mathbf{v}|^2 s'.
	\end{equation}
\end{definition}

\begin{lemma}\label{lem:cocycle}
	Let $\zeta$ be as per Definition~\ref{def:bargmann-cocycle}. Then $\zeta$ is a real $2$-cocycle on $\mathcal{G}$ as per Definition~\ref{def:cocycle}, \textit{i.e.}, $\zeta(\kappa,e)=\zeta(e,\kappa)=0$ for the identity $e$ of $\mathcal{G}$, and for all $\kappa,\kappa',\kappa''\in\mathcal{G}$
	\begin{equation}
		\zeta(\kappa,\kappa')+\zeta(\kappa\kappa',\kappa'') = \zeta(\kappa',\kappa'')+\zeta(\kappa,\kappa'\kappa'').
		\label{eq:cocycle-identity}
	\end{equation}
\end{lemma}

\begin{proof}
	The identity element is $e=(0,\mathbf{0},\mathbf{0},1)$, so both $\zeta(\kappa,e)$ and $\zeta(e,\kappa)$ vanish by inspection. Write $\kappa''=(s'',\mathbf{a}'',\mathbf{v}'',u'')$ and $R=R_u$, $R'=R_{u'}$. By Definition~\ref{def:galilei-group} the product $\kappa\kappa'$ has boost velocity $\mathbf{v}+R\mathbf{v}'$ and rotation $RR'$, while $\kappa'\kappa''$ has translation $\mathbf{a}'+R'\mathbf{a}''+s''\mathbf{v}'$. Expanding the left-hand side of Equation~\eqref{eq:cocycle-identity} and using $R\mathbf{v}'\cdot RR'\mathbf{a}''=\mathbf{v}'\cdot R'\mathbf{a}''$, one obtains
	\begin{equation*}
		\mathbf{v}\cdot R\mathbf{a}'+\mathbf{v}\cdot RR'\mathbf{a}''+\mathbf{v}'\cdot R'\mathbf{a}''+s''\,\mathbf{v}\cdot R\mathbf{v}'+\tfrac{1}{2}|\mathbf{v}|^2(s'+s'')+\tfrac{1}{2}|\mathbf{v}'|^2s''.
	\end{equation*}
	The same expansion of the right-hand side yields exactly this expression, which proves Equation~\eqref{eq:cocycle-identity}.
\end{proof}

\noindent Equation~\eqref{eq:cocycle-identity} is the associativity condition for the product defined below, and the normalization on the identity ensures that the neutral element of the extended group is the expected one.

\begin{definition}[Bargmann group]\label{def:bargmann-group}
	The \emph{Bargmann group} is the Lie group $\widehat{\mathcal{G}}:=\mathbb{R}\times\mathcal{G}$ with product
	\begin{equation}
		(c,\kappa)(c',\kappa') := \big(c+c'+\zeta(\kappa,\kappa'),\ \kappa\kappa'\big),
	\end{equation}
	where $\zeta$ is as per Definition~\ref{def:bargmann-cocycle}. For $\kappa\in\mathcal{G}$ and $c\in\mathbb{R}$ we write $\hat\kappa:=(0,\kappa)$ and $z(c):=(c,e)$.
\end{definition}

\begin{remark}\label{rem:central-extension}
	By Lemma~\ref{lem:cocycle} and Proposition~\ref{prop:cocycle-extension}, $\widehat{\mathcal{G}}$ is the central extension of $\mathcal{G}$ by $\mathbb{R}$ associated with the cocycle $\zeta$, in the sense of Definition~\ref{def:central-extension}, with $\iota=z$ and $\pi(c,\kappa)=\kappa$. Concretely, every element of $\mathcal{G}$ appears in $\widehat{\mathcal{G}}$ carrying an additional real coordinate $c$. The elements $z(c)$ are central, \textit{i.e.}, they commute with every element of $\widehat{\mathcal{G}}$, and the coordinate $c$ disappears in the quotient $\mathcal{G}$. Its only trace is the extra term $\zeta(\kappa,\kappa')$ added to the coordinates of two factors in the product of Definition~\ref{def:bargmann-group}. Moreover $(c,\kappa)=z(c)\,\hat\kappa$. The group $\widehat{\mathcal{G}}$ is the one used by L\'evy-Leblond in the geometric treatment of Galilean quantum mechanics \cite[TODO-location]{TODO-LevyLeblond1971}.
\end{remark}

The multiplication law of $\widehat{\mathcal{G}}$ has a consequence that is used repeatedly below. It shows precisely how the cocycle $\zeta$ and the central elements $z(c)$ are related, namely that the product of the lifts of two elements is the lift of their product up to a central factor whose coordinate is $\zeta$. The same identity explains how a unitary representation turns the real number $\zeta$ into a phase.

\begin{lemma}[Product of lifts]\label{lem:lift-product}
	For all $\kappa,\kappa'\in\mathcal{G}$,
	\begin{equation*}
		\hat\kappa\,\hat\kappa'=\big(\zeta(\kappa,\kappa'),\,\kappa\kappa'\big)=z\big(\zeta(\kappa,\kappa')\big)\,\widehat{\kappa\kappa'}.
	\end{equation*}
	Consequently, if $\Pi$ is a representation of $\widehat{\mathcal{G}}$ such that $\Pi(z(c))=e^{imc/\hbar}\mathbb{1}$ for some $m\in\mathbb{R}$ and all $c\in\mathbb{R}$, the operators $V_\kappa:=\Pi(\hat\kappa)$ satisfy
	\begin{equation*}
		V_\kappa V_{\kappa'}=e^{im\zeta(\kappa,\kappa')/\hbar}\,V_{\kappa\kappa'}.
	\end{equation*}
\end{lemma}

\begin{proof}
	By Definition~\ref{def:bargmann-group}, $(0,\kappa)(0,\kappa')=\big(0+0+\zeta(\kappa,\kappa'),\,\kappa\kappa'\big)$. Since $(c,\kappa'')=z(c)\,(0,\kappa'')$ for every $c\in\mathbb{R}$ and $\kappa''\in\mathcal{G}$, this equals $z\big(\zeta(\kappa,\kappa')\big)\,\widehat{\kappa\kappa'}$. Applying $\Pi$ and using that $\Pi(z(c))$ is the scalar $e^{imc/\hbar}$ yields the second statement.
\end{proof}

\noindent In words, the map $\kappa\mapsto\hat\kappa$ is a homomorphism only modulo the center, and the failure is recorded by $\zeta$ in the central coordinate. A representation then converts this central element into the phase $e^{im\zeta/\hbar}$, which is the multiplier of the projective representation $\kappa\mapsto V_\kappa$.

\begin{remark}[Meaning of the cocycle]\label{rem:cocycle-meaning}
	Let $V_\kappa$ be as in Lemma~\ref{lem:lift-product}. The cocycle measures how far $V$ is from being a homomorphism, namely by the phase $e^{im\zeta(\kappa,\kappa')/\hbar}$. This is most transparent for two commuting elements $\kappa\kappa'=\kappa'\kappa$ of $\mathcal{G}$. Applying the lemma to both orderings yields
	\begin{equation*}
		V_\kappa V_{\kappa'}=e^{\frac{im}{\hbar}\left(\zeta(\kappa,\kappa')-\zeta(\kappa',\kappa)\right)}\,V_{\kappa'}V_\kappa,
	\end{equation*}
	so that the lifts of commuting symmetries commute only up to a phase, and the antisymmetrized cocycle is the exact size of the failure. This phase cannot be removed by multiplying each $V_\kappa$ by a phase $e^{i\lambda(\kappa)}$. Such a rephasing changes the multiplier $\frac{m}{\hbar}\zeta(\kappa,\kappa')$ by $\lambda(\kappa)+\lambda(\kappa')-\lambda(\kappa\kappa')$, which is symmetric in $\kappa,\kappa'$ whenever $\kappa\kappa'=\kappa'\kappa$, and therefore cancels in the difference. For a boost and a translation the difference is $\zeta(\beta_{\mathbf{v}},\alpha_{\mathbf{a}})-\zeta(\alpha_{\mathbf{a}},\beta_{\mathbf{v}})=\mathbf{v}\cdot\mathbf{a}$, which produces the phase $e^{im\mathbf{v}\cdot\mathbf{a}/\hbar}$ of Relation~\eqref{eq:W2}. This is the instance, for the Bargmann cocycle, of Proposition~\ref{prop:commutator-pairing}, which shows in general that the antisymmetrized cocycle of two commuting elements depends only on the cohomology class. Since $\mathbf{v}\cdot\mathbf{a}\neq0$ for suitable $\mathbf{v}$ and $\mathbf{a}$, the Bargmann cocycle is not a coboundary, as made precise in Proposition~\ref{prop:bargmann-nontrivial}.
\end{remark}

We now record the commutation relations between the elementary transformations in $\widehat{\mathcal{G}}$. They are the group-level form of the Lie algebra relations of the Galilean generators, and, once represented by unitary operators, they take the place of the Heisenberg commutation relations of Equation~\eqref{eq:heisenberg-ccr}.

\begin{lemma}\label{lem:galilei-relations}
	Let $\widehat{\mathcal{G}}$ be as per Definition~\ref{def:bargmann-group}, with elementary transformations as per Definition~\ref{def:elementary-transformations}. For all $s,s'\in\mathbb{R}$, $\mathbf{a},\mathbf{a}',\mathbf{v},\mathbf{v}'\in\mathbb{R}^3$ and $u,u'\in SU(2)$,
	\begin{enumerate}
		\item $\hat\tau_s\hat\tau_{s'}=\hat\tau_{s+s'}$, $\hat\alpha_{\mathbf{a}}\hat\alpha_{\mathbf{a}'}=\hat\alpha_{\mathbf{a}+\mathbf{a}'}$, $\hat\beta_{\mathbf{v}}\hat\beta_{\mathbf{v}'}=\hat\beta_{\mathbf{v}+\mathbf{v}'}$ and $\hat\rho_u\hat\rho_{u'}=\hat\rho_{uu'}$,
		\item $\hat\tau_s\hat\alpha_{\mathbf{a}}=\hat\alpha_{\mathbf{a}}\hat\tau_s$,
		\item $\hat\beta_{\mathbf{v}}\hat\alpha_{\mathbf{a}}=z(\mathbf{v}\cdot\mathbf{a})\,\hat\alpha_{\mathbf{a}}\hat\beta_{\mathbf{v}}$,
		\item $\hat\tau_{-s}\hat\beta_{\mathbf{v}}\hat\tau_{s}=z\big(-\tfrac{1}{2}s|\mathbf{v}|^2\big)\,\hat\beta_{\mathbf{v}}\hat\alpha_{s\mathbf{v}}$,
		\item $\hat\rho_u\hat\tau_s\hat\rho_u^{-1}=\hat\tau_s$, $\hat\rho_u\hat\alpha_{\mathbf{a}}\hat\rho_u^{-1}=\hat\alpha_{R_u\mathbf{a}}$ and $\hat\rho_u\hat\beta_{\mathbf{v}}\hat\rho_u^{-1}=\hat\beta_{R_u\mathbf{v}}$,
		\item every element of $\widehat{\mathcal{G}}$ is a product $z(c)\,\hat\tau_s\hat\alpha_{\mathbf{a}}\hat\beta_{\mathbf{v}}\hat\rho_u$.
	\end{enumerate}
\end{lemma}

\begin{proof}
	\emph{Items 1, 2 and 5.} In every product involved the left factor has vanishing boost velocity, so $\zeta$ vanishes and the claim follows from Definition~\ref{def:galilei-group}.
	
	\emph{Item 3.} This is the identity in which the cocycle $\zeta$ and the central element $z$ meet, so we display every equality. By Lemma~\ref{lem:lift-product},
	\begin{align*}
		\hat\beta_{\mathbf{v}}\hat\alpha_{\mathbf{a}}
		&=\big(\zeta(\beta_{\mathbf{v}},\alpha_{\mathbf{a}}),\,\beta_{\mathbf{v}}\alpha_{\mathbf{a}}\big)
		=\big(\mathbf{v}\cdot\mathbf{a},\,(0,\mathbf{a},\mathbf{v},1)\big)\\
		&=z(\mathbf{v}\cdot\mathbf{a})\,\big(0,\alpha_{\mathbf{a}}\beta_{\mathbf{v}}\big)
		=z(\mathbf{v}\cdot\mathbf{a})\,\hat\alpha_{\mathbf{a}}\hat\beta_{\mathbf{v}}.
	\end{align*}
	The first equality is the product of Definition~\ref{def:bargmann-group}. In the second, $\zeta(\beta_{\mathbf{v}},\alpha_{\mathbf{a}})=\mathbf{v}\cdot\mathbf{a}+\tfrac{1}{2}|\mathbf{v}|^2\cdot 0=\mathbf{v}\cdot\mathbf{a}$, because the second factor has vanishing time component, and $\beta_{\mathbf{v}}\alpha_{\mathbf{a}}=(0,\mathbf{a},\mathbf{v},1)$ by Definition~\ref{def:galilei-group}. The third equality splits off the central coordinate through $(c,\kappa)=z(c)\,(0,\kappa)$, using that $\alpha_{\mathbf{a}}\beta_{\mathbf{v}}=(0,\mathbf{a},\mathbf{v},1)$ as well. The last one holds because $\zeta(\alpha_{\mathbf{a}},\beta_{\mathbf{v}})=0$, so that $(0,\alpha_{\mathbf{a}}\beta_{\mathbf{v}})=\hat\alpha_{\mathbf{a}}\hat\beta_{\mathbf{v}}$. The two orderings thus differ exactly by the central element $z$ whose coordinate is $\zeta(\beta_{\mathbf{v}},\alpha_{\mathbf{a}})-\zeta(\alpha_{\mathbf{a}},\beta_{\mathbf{v}})=\mathbf{v}\cdot\mathbf{a}$, in agreement with Remark~\ref{rem:cocycle-meaning}.
	
	\emph{Item 4.} $\tau_{-s}\beta_{\mathbf{v}}=(-s,\mathbf{0},\mathbf{v},1)$ with vanishing cocycle, and $\zeta\big((-s,\mathbf{0},\mathbf{v},1),\tau_s\big)=\tfrac{1}{2}s|\mathbf{v}|^2$, so that $\hat\tau_{-s}\hat\beta_{\mathbf{v}}\hat\tau_s=\big(\tfrac{1}{2}s|\mathbf{v}|^2,(0,s\mathbf{v},\mathbf{v},1)\big)$. On the other side $\zeta(\beta_{\mathbf{v}},\alpha_{s\mathbf{v}})=s|\mathbf{v}|^2$ and $\beta_{\mathbf{v}}\alpha_{s\mathbf{v}}=(0,s\mathbf{v},\mathbf{v},1)$, which yields the stated phase $\tfrac{1}{2}s|\mathbf{v}|^2-s|\mathbf{v}|^2$.
	
	\emph{Item 6.} $\hat\tau_s\hat\alpha_{\mathbf{a}}\hat\beta_{\mathbf{v}}\hat\rho_u=\hat{\kappa}$ with $\kappa=(s,\mathbf{a},\mathbf{v},u)$, because all the cocycles involved vanish, and $(c,\kappa)=z(c)\hat\kappa$.
\end{proof}

\noindent Items 3 and 4 are the only ones containing a phase, and they are the origin of the mass parameter and of the kinetic term respectively.

% =====================================================================
\section{Galilean Invariance of the Projective State Space}
\label{sec:invariance}

The Kähler structure of $\mathbb{P}\mathcal{H}$ constrains the transformations that a physical symmetry may induce on the pure states, since they must preserve the symplectic form, the metric and the complex structure. We require the Galilei group to act through such transformations, and this is the only assumption on the dynamics entering the chapter. A Wigner type argument turns the geometric action into a unitary ray representation, and a theorem of Bargmann turns the latter into a genuine unitary representation of the Bargmann group, at the price of a real parameter $m$. This parameter will be identified with the mass of the system in Section~\ref{sec:free-hamiltonian}, since it governs the commutation between boosts and translations. We state the invariance requirement first and then derive its consequences at the level of unitary operators.

\begin{definition}[Galilean action on the projective Hilbert space]\label{def:galilean-action}
	Let $\mathcal{H}$ be a separable Hilbert space as per Definition~\ref{def:hilbert-space} and let $\mathbb{P}\mathcal{H}$ be its projective space as per Definition~\ref{def:projective-hilbert-space}, equipped with the Kähler structure $(J,g,\omega)$ of Equation~\eqref{eq:physical-kahler}. A \emph{Galilean action} on $\mathbb{P}\mathcal{H}$ is a group homomorphism $\Phi$ from $\mathcal{G}$, as per Definition~\ref{def:galilei-group}, to the group of Kähler isomorphisms of $\mathbb{P}\mathcal{H}$, written $\kappa\mapsto\Phi_\kappa$, such that for every $\phi\in\mathbb{P}\mathcal{H}$ the orbit map $\kappa\mapsto\Phi_\kappa(\phi)$ is continuous with respect to the manifold topology of Proposition~\ref{prop:holomorphic-atlas}.
\end{definition}

\begin{remark}
	By Proposition~\ref{prop:topologies-coincide} the manifold topology of $\mathbb{P}\mathcal{H}$ is the weak$^*$ topology, hence continuity of the orbit maps means that expectation values of all observables vary continuously along the orbits. By Remark~\ref{rem:riemannian-distance} the distance $d_g$ induces the manifold topology, and by Theorem~\ref{thm:geometrical-born} the transition probability $P(\psi,\Phi_\kappa(\phi))$ is a strictly monotone function of $d_g(\psi,\Phi_\kappa(\phi))$ on its range. Therefore Definition~\ref{def:galilean-action} is equivalent to the continuity of $\kappa\mapsto P(\psi,\Phi_\kappa(\phi))$ for all $\psi,\phi\in\mathbb{P}\mathcal{H}$, which is the continuity hypothesis of the theory of ray representations.
\end{remark}

An elementary physical system is one whose pure states cannot be separated into two sectors by the symmetry group. This is expressed geometrically as follows.

\begin{definition}[Irreducible Galilean action]\label{def:irreducible-action}
	Let $\Phi$ be a Galilean action as per Definition~\ref{def:galilean-action}. The action is \emph{irreducible} if the only closed complex subspaces $\mathcal{K}\subseteq\mathcal{H}$ such that $\Phi_\kappa(\mathbb{P}\mathcal{K})\subseteq\mathbb{P}\mathcal{K}$ for all $\kappa\in\mathcal{G}$ are $\{0\}$ and $\mathcal{H}$.
\end{definition}

\subsection*{Multipliers, Cohomology and the Origin of the Mass}

The lifts $V_\kappa$ of a Galilean action are determined only up to a phase, so that $V_\kappa V_{\kappa'}$ and $V_{\kappa\kappa'}$ may differ by a phase which cannot be avoided. Whether this discrepancy can be removed by a suitable choice of the phases is a cohomological question in the subsection on cocycles and central extensions that follows Definition~\ref{def:strongly-continuous}. There it was shown that the obstruction is the class of the multiplier, and, in Proposition~\ref{prop:bargmann-correspondence}, that for a connected and simply connected group this class can be computed on the Lie algebra. In this subsection we carry out the computation for the Galilei group, in order to see where the real parameter $m$ comes from. We first show that the Bargmann cocycle is not a coboundary and that different multiples of it are not cohomologous. We then read Bargmann's theorem for the Galilei group in four steps, and only afterwards state it.

\begin{proposition}[The Bargmann cocycle is not trivial]\label{prop:bargmann-nontrivial}
	Let $\zeta$ be the Bargmann cocycle of Definition~\ref{def:bargmann-cocycle}.
	\begin{enumerate}
		\item For $m,m'\in\mathbb{R}$ the real cocycles $m\zeta$ and $m'\zeta$ are cohomologous if and only if $m=m'$. In particular $\zeta$ is not a coboundary.
		\item Let $X^\alpha_j$ and $X^\beta_j$ be the elements of the Lie algebra of $\mathcal{G}$ such that $\exp(tX^\alpha_j)=\alpha_{te_j}$ and $\exp(tX^\beta_j)=\beta_{te_j}$. Then $[X^\beta_i,X^\alpha_j]=0$ and the infinitesimal cocycle of Definition~\ref{def:infinitesimal-cocycle} satisfies $\psi_\zeta(X^\beta_i,X^\alpha_j)=d_{ij}$.
	\end{enumerate}
\end{proposition}

\begin{proof}
	Item 1. By Lemma~\ref{lem:galilei-relations}, $\beta_{\mathbf{v}}\alpha_{\mathbf{a}}=\alpha_{\mathbf{a}}\beta_{\mathbf{v}}=(0,\mathbf{a},\mathbf{v},1)$, so that $\beta_{\mathbf{v}}$ and $\alpha_{\mathbf{a}}$ commute in $\mathcal{G}$. The commutator pairing of Proposition~\ref{prop:commutator-pairing} gives $\Omega_{m\zeta}(\beta_{\mathbf{v}},\alpha_{\mathbf{a}})=m\big(\zeta(\beta_{\mathbf{v}},\alpha_{\mathbf{a}})-\zeta(\alpha_{\mathbf{a}},\beta_{\mathbf{v}})\big)=m\,\mathbf{v}\cdot\mathbf{a}$. If $m\zeta-m'\zeta=(m-m')\zeta$ were a coboundary, item 2 of that proposition would give $(m-m')\,\mathbf{v}\cdot\mathbf{a}=0$ for all $\mathbf{v},\mathbf{a}\in\mathbb{R}^3$, hence $m=m'$. The converse is trivial.
	
	Item 2. The one-parameter subgroups $t\mapsto\beta_{te_i}$ and $t\mapsto\alpha_{te_j}$ commute, so the Lie bracket of their generators vanishes. Moreover $\zeta(\beta_{se_i},\alpha_{te_j})-\zeta(\alpha_{te_j},\beta_{se_i})=st\,d_{ij}$, whose mixed derivative at the origin is $d_{ij}$.
\end{proof}

\noindent Item 2 gives a second proof that $\zeta$ is not a coboundary, at the level of the Lie algebra. Every coboundary $\psi_\lambda$ of Definition~\ref{def:lie-cohomology-group} vanishes on the pair $(X^\beta_i,X^\alpha_j)$, since $\lambda([X^\beta_i,X^\alpha_j])=0$, whereas $\psi_\zeta$ does not.

Proposition~\ref{prop:bargmann-nontrivial} shows that the numbers $m$ label distinct classes. It does not show that they exhaust the classes. This requires the computation of the whole second cohomology group of the Galilei algebra, which is long, and we recall its outcome.

\begin{proposition}[Second cohomology of the Galilei algebra]\label{prop:galilei-cohomology}
	Let $\mathfrak{g}$ be the Lie algebra of $\mathcal{G}$. Then $H^2(\mathfrak{g},\mathbb{R})\cong\mathbb{R}$, and it is generated by the class of $\psi_\zeta$.
\end{proposition}

\noindent See \cite[TODO-location]{TODO-Bargmann1954}. This is the statement that the Galilei group in three space dimensions has, up to equivalence, only the one-parameter family of cocycles $m\zeta$.

We can now read Bargmann's theorem for the Galilei group. Each step below is a result already established, and the only point that remains to be explained is how they combine.

\begin{remark}[Reading Bargmann's theorem for the Galilei group]\label{rem:bargmann-galilei}
	\emph{Step 1, the multiplier.} The lifts $V_\kappa$ of $\Phi$ satisfy $V_\kappa V_{\kappa'}=\omega(\kappa,\kappa')V_{\kappa\kappa'}$ with a multiplier $\omega$, which is a cocycle by Proposition~\ref{prop:multiplier-cocycle}.
	
	\emph{Step 2, the class.} Only the class of $\omega$ matters, by Proposition~\ref{prop:rephasing}. Boosts and translations commute in $\mathcal{G}$, but their lifts commute only up to a phase, which by Proposition~\ref{prop:commutator-pairing} does not depend on the choice of the phases. If this phase is nontrivial, the class is nontrivial and no rephasing yields a genuine representation of $\mathcal{G}$.
	
	\emph{Step 3, the Lie algebra.} Since $\mathcal{G}$ is connected and simply connected, the class is an element of $H^2(\mathfrak{g},\mathbb{R})$ by Proposition~\ref{prop:bargmann-correspondence}. Concretely, the generators satisfy $[K_i,P_j]=i\hbar m\,d_{ij}\mathbb{1}$ on the smooth vectors, as derived in the proof of Proposition~\ref{prop:galilean-poisson}, and shifting a generator by a multiple of $\mathbb{1}$ does not change this commutator.
	
	\emph{Step 4, the parameter.} By Proposition~\ref{prop:galilei-cohomology} the group $H^2(\mathfrak{g},\mathbb{R})$ is one dimensional, so the whole freedom is a single real number, and by Proposition~\ref{prop:bargmann-nontrivial} this number labels the class of $\frac{m}{\hbar}\zeta$ injectively. This number is the mass. The same argument for $SU(2)$ yields no parameter, because $H^2(\mathfrak{su}(2),\mathbb{R})=0$ by Example~\ref{ex:su2-cohomology}. The only obstruction for the rotations is the topological one of Example~\ref{ex:so3-su2}, which is removed by working with the universal covering.
\end{remark}

The passage from a Kähler action to a unitary representation relies on Bargmann's classification of the multipliers of the Galilei group, which we now state and derive from Proposition~\ref{prop:bargmann-correspondence}.

\begin{theorem}[Bargmann]\label{thm:bargmann}
	Let $\mathcal{H}$ be a separable Hilbert space and let $\kappa\mapsto V_\kappa$ be a continuous homomorphism from $\mathcal{G}$ to the group $U(\mathcal{H})/U(1)$ of unitary ray transformations of $\mathcal{H}$, endowed with the topology in which $V_n\to V$ if and only if $|(x,V_ny)|\to|(x,Vy)|$ for all $x,y\in\mathcal{H}$. Then there exist a unique $m\in\mathbb{R}$ and a strongly continuous unitary representation $\Pi:\widehat{\mathcal{G}}\to U(\mathcal{H})$, as per Definition~\ref{def:strongly-continuous}, such that $\Pi(z(c))=e^{imc/\hbar}\mathbb{1}$ for all $c\in\mathbb{R}$ and $\Pi(\hat\kappa)$ belongs to the ray $V_\kappa$ for all $\kappa\in\mathcal{G}$.
\end{theorem}

\begin{proof}
	Let $V$ be as in the statement. By items 1 and 2 of Proposition~\ref{prop:bargmann-correspondence}, the class of the multiplier of $V$ corresponds to an element of $H^2(\mathfrak{g},\mathbb{R})$, where $\mathfrak{g}$ is the Lie algebra of the connected and simply connected group $\mathcal{G}$. By Proposition~\ref{prop:galilei-cohomology} this element is a real multiple of $[\psi_\zeta]$, and since $\psi_{c\zeta}=c\,\psi_\zeta$ it can be written as $[\psi_{(m/\hbar)\zeta}]$ for a real number $m$. This number is unique by Proposition~\ref{prop:bargmann-nontrivial}, because $\frac{m}{\hbar}\zeta$ and $\frac{m'}{\hbar}\zeta$ are cohomologous only for $m=m'$. Item 3 of Proposition~\ref{prop:bargmann-correspondence}, applied to the continuous real cocycle $\frac{m}{\hbar}\zeta$, shows that the phases of $V$ can be chosen so that $V$ is strongly continuous with multiplier exactly $e^{im\zeta/\hbar}$, and that $\Pi(c,\kappa):=e^{imc/\hbar}V_\kappa$ is a strongly continuous unitary representation of $\widehat{\mathcal{G}}$. It satisfies $\Pi(z(c))=e^{imc/\hbar}\mathbb{1}$, and $\Pi(\hat\kappa)=V_\kappa$ belongs to the ray of the original operator.
\end{proof}

\noindent The inputs which are not proved in the text are Proposition~\ref{prop:bargmann-correspondence} and the computation of Proposition~\ref{prop:galilei-cohomology}.

The theorem is stated for unitary rays, whereas Definition~\ref{def:galilean-action} involves Kähler isomorphisms. The next proposition shows that the two coincide in the present setting and records how much freedom remains in the lift.

\begin{proposition}[Unitary lift of a Galilean action]\label{prop:lift}
	Let $\Phi$ be a Galilean action on $\mathbb{P}\mathcal{H}$. Then there exist a unique $m\in\mathbb{R}$ and a strongly continuous unitary representation $\Pi:\widehat{\mathcal{G}}\to U(\mathcal{H})$ such that $\Pi(z(c))=e^{imc/\hbar}\mathbb{1}$ and
	\begin{equation}
		\Phi_\kappa([x])=[\Pi(\hat\kappa)x], \qquad \kappa\in\mathcal{G},\ [x]\in\mathbb{P}\mathcal{H}.
		\label{eq:lift}
	\end{equation}
	If $\Pi'$ is another such representation with the same $m$, then $\Pi'(c,\kappa)=e^{i\epsilon s}\,\Pi(c,\kappa)$ for some $\epsilon\in\mathbb{R}$, where $s$ is the time translation of $\kappa=(s,\mathbf{a},\mathbf{v},u)$.
\end{proposition}

\begin{proof}
	By Proposition~\ref{prop:kahler-isomorphisms}, each $\Phi_\kappa$ is induced by a unitary operator $V_\kappa$, unique up to a phase. An antiunitary operator would reverse the sign of $\omega=2\hbar\,\mathrm{Im}(\cdot,\cdot)$ and is therefore excluded. The homomorphism property of $\Phi$ yields $V_\kappa V_{\kappa'}\in U(1)\,V_{\kappa\kappa'}$, so that $\kappa\mapsto[V_\kappa]$ is a homomorphism into $U(\mathcal{H})/U(1)$. It is continuous in the topology of Theorem~\ref{thm:bargmann}, since $|(x,V_\kappa y)|^2/(\|x\|^2\|y\|^2)$ is the transition probability between $[x]$ and $\Phi_\kappa[y]$, which is continuous in $\kappa$ by the previous remark. Theorem~\ref{thm:bargmann} yields $m$ and $\Pi$ with the required properties, and $m$ is unique because it labels the multiplier class of $V$.
	
	If $\Pi'$ is a second lift with the same $m$, then $\Pi'(\hat\kappa)=\chi(\kappa)\Pi(\hat\kappa)$ for a function $\chi:\mathcal{G}\to U(1)$, continuous by strong continuity. Both representations satisfy the same product law with the same phases, hence $\chi$ is a homomorphism. Since $\mathcal{G}$ is connected and simply connected, $\chi$ is determined by a character of the Lie algebra of $\mathcal{G}$, which vanishes on the derived algebra. By Lemma~\ref{lem:galilei-relations}, differentiating items 3, 4 and 5 shows that the derived algebra contains the generators of the translations, of the boosts and of the rotations, so that only the time translation survives. This yields $\chi(\kappa)=e^{i\epsilon s}$.
\end{proof}

\noindent The freedom $\Pi\mapsto e^{i\epsilon s}\Pi$ will correspond to an additive constant in the Hamiltonian, in agreement with Proposition~\ref{prop:kahler-isomorphisms}.

From now on $\Pi$ denotes a fixed lift of $\Phi$ as in Proposition~\ref{prop:lift}. Its restriction to the elementary transformations consists of strongly continuous unitary groups, and Stone's theorem identifies their generators.

\begin{definition}[Galilean generators]\label{def:generators}
	Let $\Pi$ be as in Proposition~\ref{prop:lift}. The unitary groups
	\begin{equation}
		S(s):=\Pi(\hat\tau_s), \qquad T(\mathbf{a}):=\Pi(\hat\alpha_{\mathbf{a}}), \qquad B(\mathbf{v}):=\Pi(\hat\beta_{\mathbf{v}})
	\end{equation}
	are strongly continuous. The selfadjoint operators $H$, $\mathbf{P}=(P_1,P_2,P_3)$ and $\mathbf{K}=(K_1,K_2,K_3)$ are defined, through Theorem~\ref{thm:stone} for $S$ and through Theorem~\ref{thm:snag} for $T$ and $B$, by
	\begin{equation}
		S(s)=e^{isH/\hbar}, \qquad T(\mathbf{a})=e^{-i\mathbf{a}\cdot\mathbf{P}/\hbar}, \qquad B(\mathbf{v})=e^{i\mathbf{v}\cdot\mathbf{K}/\hbar}.
	\end{equation}
	We denote by $E$ the projection valued measure, as per Definition~\ref{def:pvm}, which Theorem~\ref{thm:snag} associates with $\mathbf{t}\mapsto T(-\hbar\mathbf{t})$, so that $T(\mathbf{a})=\int e^{-i\mathbf{a}\cdot\mathbf{p}/\hbar}\,dE(\mathbf{p})$ and $P_j=\int p_j\,dE(\mathbf{p})$.
\end{definition}

\begin{remark}\label{rem:active-convention}
	The commutativity of the three groups of translations and boosts, stated in item 1 of Lemma~\ref{lem:galilei-relations}, is what allows $\mathbf{P}$ and $\mathbf{K}$ to have commuting components. The signs in Definition~\ref{def:generators} follow the active convention, in which $\tau_s$ shifts the events of a worldline towards the future. A state prepared on the hypersurface $t=0$ is mapped by $\Phi_{\tau_s}$ to the state which on $t=0$ looks like the original one on $t=-s$. Hence the Schrödinger flow of Proposition~\ref{prop:hamiltonian-schrodinger} is $t\mapsto\Phi_{\tau_{-t}}$, that is $[x]\mapsto[e^{-itH/\hbar}x]$, and $H$ is the energy operator of Section~\ref{sec:hamiltonian-flow}. The sign of the boost generator is fixed by requiring that $B(\mathbf{v})$ increases the momentum by $m\mathbf{v}$, as verified in Corollary~\ref{cor:schrodinger-realization}.
\end{remark}

The group relations of Lemma~\ref{lem:galilei-relations} become relations between unitary operators once $\Pi$ is applied, and the central elements are replaced by the phases $e^{imc/\hbar}$.

\begin{proposition}[Relations between the Galilean unitary groups]\label{prop:unitary-relations}
	Let $S$, $T$, $B$ be as in Definition~\ref{def:generators}. For all $s\in\mathbb{R}$, $\mathbf{a},\mathbf{v}\in\mathbb{R}^3$ and $u\in SU(2)$,
	\begin{align}
		S(s)\,T(\mathbf{a}) &= T(\mathbf{a})\,S(s), \tag{W1}\label{eq:W1}\\
		B(\mathbf{v})\,T(\mathbf{a}) &= e^{im\mathbf{v}\cdot\mathbf{a}/\hbar}\,T(\mathbf{a})\,B(\mathbf{v}), \tag{W2}\label{eq:W2}\\
		S(-s)\,B(\mathbf{v})\,S(s) &= e^{-ims|\mathbf{v}|^2/2\hbar}\,B(\mathbf{v})\,T(s\mathbf{v}), \tag{W3}\label{eq:W3}\\
		\Pi(\hat\rho_u)\,S(s)\,\Pi(\hat\rho_u)^{-1} &= S(s), \tag{W4a}\label{eq:W4a}\\
		\Pi(\hat\rho_u)\,T(\mathbf{a})\,\Pi(\hat\rho_u)^{-1} &= T(R_u\mathbf{a}), \tag{W4b}\label{eq:W4b}\\
		\Pi(\hat\rho_u)\,B(\mathbf{v})\,\Pi(\hat\rho_u)^{-1} &= B(R_u\mathbf{v}). \tag{W4c}\label{eq:W4c}
	\end{align}
	Moreover the operators $S(s)$, $T(\mathbf{a})$, $B(\mathbf{v})$, $\Pi(\hat\rho_u)$ and the scalars $e^{imc/\hbar}$ generate $\Pi(\widehat{\mathcal{G}})$.
\end{proposition}

\begin{proof}
	Apply $\Pi$ to items 2 to 5 of Lemma~\ref{lem:galilei-relations} and use $\Pi(z(c))=e^{imc/\hbar}\mathbb{1}$. The last statement is item 6 of the same lemma.
\end{proof}

\begin{remark}\label{rem:weyl}
	Relation~\eqref{eq:W2} is the Weyl form of the canonical commutation relations of Proposition~\ref{prop:weyl-relations}, and its content is best read in three steps.
	
	\emph{Matching the Weyl relations.} Setting $U(\boldsymbol{\alpha}):=B(\boldsymbol{\alpha}/m)=e^{i\boldsymbol{\alpha}\cdot\mathbf{K}/m\hbar}$ and $V(\boldsymbol{\beta}):=T(-\boldsymbol{\beta})=e^{i\boldsymbol{\beta}\cdot\mathbf{P}/\hbar}$, Relation~\eqref{eq:W2} becomes
	\begin{equation*}
		U(\boldsymbol{\alpha})V(\boldsymbol{\beta})=V(\boldsymbol{\beta})U(\boldsymbol{\alpha})\,e^{-i\boldsymbol{\alpha}\cdot\boldsymbol{\beta}/\hbar},
	\end{equation*}
	which is Equation~\eqref{eq:weyl-relations} up to the rescaling $(\boldsymbol{\alpha},\boldsymbol{\beta})\mapsto\hbar^{-1/2}(\boldsymbol{\alpha},\boldsymbol{\beta})$ that restores $\hbar$, set equal to one in Section~\ref{sec:weyl-algebra}. Formally, $\mathbf{Q}:=\mathbf{K}/m$ plays the role of the position operator and $\mathbf{P}$ that of the momentum. As in Section~\ref{sec:weyl-algebra}, neither of them is an element of a $C^*$-algebra, and they enter the theory only through the bounded unitary groups $B$ and $T$, which is consistent with Proposition~\ref{prop:heisenberg-obstruction}.
	
	\emph{Physical content of the phase.} Conjugation by a boost shifts the momentum. By Lemma~\ref{lem:boost-covariance}, $B(\mathbf{v})^{-1}\mathbf{P}B(\mathbf{v})=\mathbf{P}+m\mathbf{v}$, so that a boost of velocity $\mathbf{v}$ changes the momentum by $m\mathbf{v}$. The parameter $m$ is therefore the constant relating a change of velocity to a change of momentum, \textit{i.e.}, the mass, and it is not assumed but read off from the phase in Relation~\eqref{eq:W2}. By Proposition~\ref{prop:lift}, $m$ is determined by the action $\Phi$ and does not depend on the chosen lift.
	
	\emph{The case $m=0$.} For $m=0$, Relation~\eqref{eq:W2} states that boosts commute with translations, so that $\mathbf{K}$ and $\mathbf{P}$ commute and no pair of conjugate observables is available. Since the position operator is built from this pair, we exclude this case and assume $m\neq0$ from now on. The sign of $m$ is not restricted by the relations of this section, and it will be fixed by the positivity of the energy in Theorem~\ref{thm:free-hamiltonian}.
\end{remark}

% =====================================================================
\section{Hamiltonian Functions of the Galilean Generators}
\label{sec:galilean-functions}

The generators of Definition~\ref{def:generators} are selfadjoint but unbounded, whereas the mean value functions of Chapter~\ref{ch:projective} were introduced for bounded operators, for which they are smooth functions on the whole of $\mathbb{P}\mathcal{H}$. The Poisson bracket of Definition~\ref{def:poisson-riemann-brackets} requires such smoothness, and the proofs of Proposition~\ref{prop:poisson-commutator} and Proposition~\ref{prop:jordan-riemann} do not apply verbatim to unbounded operators. We observe however that those proofs are pointwise computations involving only the representatives $x$, $Ax$ and $Bx$, so that they extend to a dense domain of regular vectors. We make this extension precise, and we then use it to read the Galilean relations of Section~\ref{sec:invariance} as a Poisson algebra on $\mathbb{P}\mathcal{H}$. This translation is not needed for the proof of the main theorem, but it exhibits the geometric content of the result and connects it to the momentum map of Section~\ref{sec:momentum-map}.

\begin{definition}[Mean value function on a domain]\label{def:mean-value-domain}
	Let $\mathcal{H}$ be a Hilbert space and let $\mathcal{D}\subseteq\mathcal{H}$ be a dense linear subspace. Let $A$ be a symmetric operator, as per Definition~\ref{def:symmetric-operator}, with $\mathcal{D}\subseteq D(A)$. Denote by $\mathbb{P}\mathcal{D}\subseteq\mathbb{P}\mathcal{H}$ the set of rays $[x]$ with $x\in\mathcal{D}\setminus\{0\}$. The mean value function of $A$ on $\mathbb{P}\mathcal{D}$ is the map $\langle A\rangle:\mathbb{P}\mathcal{D}\to\mathbb{R}$ defined by
	\begin{equation}
		\langle A\rangle([x]) := \frac{(x,Ax)}{\|x\|^2}.
	\end{equation}
\end{definition}

\begin{remark}
	For $A\in B(\mathcal{H})$ and $\mathcal{D}=\mathcal{H}$ one recovers Definition~\ref{def:mean-value}. For unbounded $A$ the set $\mathbb{P}\mathcal{D}$ is dense in $\mathbb{P}\mathcal{H}$ but not open, so that $\langle A\rangle$ is not a smooth function on the Hilbertian manifold of Definition~\ref{def:hilbertian-manifold}. The brackets of Definition~\ref{def:poisson-riemann-brackets} cannot be defined through the musical isomorphisms, and we define them pointwise through the representatives of the Hamiltonian vector fields, following the proofs quoted above.
\end{remark}

\begin{lemma}[Pointwise Kähler formulas for unbounded operators]\label{lem:pointwise-formulas}
	Let $A$ and $B$ be selfadjoint operators on $\mathcal{H}$ and let $\mathcal{D}\subseteq D(A)\cap D(B)$ be a dense subspace with $A\mathcal{D}\subseteq\mathcal{D}$ and $B\mathcal{D}\subseteq\mathcal{D}$. Let $[x]\in\mathbb{P}\mathcal{D}$ with $\|x\|=1$, and write $f_A:=\langle A\rangle$, $f_B:=\langle B\rangle$ as per Definition~\ref{def:mean-value-domain}. Set
	\begin{equation}
		\eta_B := -\frac{i}{\hbar}\big(B-f_B([x])\big)x \in x^\perp\cap\mathcal{D}.
	\end{equation}
	Then the following hold.
	\begin{enumerate}
		\item The curve $t\mapsto[e^{-itB/\hbar}x]$ is differentiable at $t=0$ in $\mathbb{P}\mathcal{H}$, with velocity represented by $\eta_B$ under the identification $T_{[x]}\mathbb{P}\mathcal{H}\cong x^\perp$ of Proposition~\ref{prop:local-chart}.
		\item For every $\eta\in x^\perp\cap\mathcal{D}$ the function $s\mapsto f_A([x+s\eta])$ is differentiable at $s=0$, with derivative $2\,\mathrm{Re}(\eta,Ax)$.
	\end{enumerate}
	Define $\{f_A,f_B\}([x]):=2\,\mathrm{Re}(\eta_B,Ax)$, namely the derivative of $f_A$ along the velocity of the flow generated by $B$, and $((f_A,f_B))([x]):=2\hbar\,\mathrm{Re}(\eta_A,\eta_B)$. Then
	\begin{align}
		\{f_A,f_B\}([x]) &= \Big\langle\frac{1}{i\hbar}[A,B]\Big\rangle([x]), \label{eq:pointwise-poisson}\\
		((f_A,f_B))([x]) &= \frac{2}{\hbar}\Big(\langle A\circ B\rangle([x])-f_A([x])\,f_B([x])\Big), \label{eq:pointwise-riemann}
	\end{align}
	with $[A,B]=AB-BA$ and $A\circ B=\tfrac{1}{2}(AB+BA)$ defined on $\mathcal{D}$. For bounded $A$ and $B$ these are the values of the Poisson bracket and of the Riemann bracket of Definition~\ref{def:poisson-riemann-brackets} between $f_A$ and $f_B$.
\end{lemma}

\begin{proof}
	Item 1. Since $x\in D(B)$, the curve $t\mapsto e^{-itB/\hbar}x$ is strongly differentiable at $t=0$ with derivative $-\tfrac{i}{\hbar}Bx$. Composing with the chart $F_{[x]}$ of Definition~\ref{def:local-chart}, which sends $[\Phi]$ to $\Phi/(x,\Phi)-x$, the derivative of $\Phi_t/(x,\Phi_t)$ at $t=0$ is $-\tfrac{i}{\hbar}Bx+\tfrac{i}{\hbar}(x,Bx)x=\eta_B$, which lies in $x^\perp$.
	
	Item 2. By symmetry of $A$ on $\mathcal{D}$ and $(x,\eta)=0$,
	\begin{equation*}
		f_A([x+s\eta])=\frac{(x,Ax)+2s\,\mathrm{Re}(\eta,Ax)+s^2(\eta,A\eta)}{1+s^2\|\eta\|^2},
	\end{equation*}
	and the quotient rule at $s=0$ isolates the linear coefficient of the numerator.
	
	For Equation~\eqref{eq:pointwise-poisson}, $(\eta_B,Ax)=\tfrac{i}{\hbar}\big((Bx,Ax)-f_Af_B\big)$, hence $2\,\mathrm{Re}(\eta_B,Ax)=-\tfrac{2}{\hbar}\mathrm{Im}(Bx,Ax)$. Since $Bx\in D(A)$ and $Ax\in D(B)$, $(x,[A,B]x)=(Ax,Bx)-(Bx,Ax)=-2i\,\mathrm{Im}(Bx,Ax)$, and dividing by $i\hbar$ yields the claim. For Equation~\eqref{eq:pointwise-riemann}, $\eta_A,\eta_B\in x^\perp$ and Lemma~\ref{lem:metric-ambient} give $g(\eta_A,\eta_B)=2\hbar\,\mathrm{Re}(\eta_A,\eta_B)=\tfrac{2}{\hbar}\big(\mathrm{Re}(Ax,Bx)-f_Af_B\big)$, and $\mathrm{Re}(Ax,Bx)=\mathrm{Re}(x,ABx)=\langle A\circ B\rangle([x])$ by invariance of $\mathcal{D}$. The last statement follows by comparison with the proofs of Proposition~\ref{prop:poisson-commutator} and Proposition~\ref{prop:jordan-riemann}.
\end{proof}

\noindent Equation~\eqref{eq:pointwise-riemann} is the pointwise form of $\langle A\circ B\rangle=\tfrac{\hbar}{2}((f_A,f_B))+f_Af_B$, that is, of the product $\circ_{\hbar}$ of Definition~\ref{def:kahler-bracket}.

The Galilean relations of Proposition~\ref{prop:unitary-relations} are relations between unitary groups, and they must be differentiated to obtain commutators of the generators. This requires a common invariant domain, which exists by the following classical result on smooth vectors.

\begin{theorem}[Smooth vectors]\label{thm:smooth-vectors}
	Let $\Pi$ be a strongly continuous unitary representation of a Lie group $\mathcal{L}$ on a Hilbert space $\mathcal{H}$, and let $\mathcal{H}^\infty$ be the subspace of vectors $x$ such that $\kappa\mapsto\Pi(\kappa)x$ is smooth. Then $\mathcal{H}^\infty$ is dense, it is invariant under $\Pi(\kappa)$ for all $\kappa\in\mathcal{L}$, and it is invariant under $d\Pi(\xi)x:=\frac{d}{dt}\Pi(\exp t\xi)x\big|_{t=0}$ for all $\xi$ in the Lie algebra of $\mathcal{L}$. For each $\xi$, the operator $d\Pi(\xi)$ restricted to $\mathcal{H}^\infty$ is essentially skew-adjoint, and its closure is the generator of $t\mapsto\Pi(\exp t\xi)$ given by Theorem~\ref{thm:stone}. Moreover $d\Pi([\xi,\eta])=[d\Pi(\xi),d\Pi(\eta)]$ on $\mathcal{H}^\infty$.
\end{theorem}

\noindent These facts are due to G\aa rding and Nelson \cite[TODO-location]{TODO-Garding1947,TODO-Nelson1959}, and the essential skew-adjointness follows from the invariance of $\mathcal{H}^\infty$ under the unitary groups \cite[TODO-location]{TODO-RS03}.

Theorem~\ref{thm:smooth-vectors} defines $\mathcal{H}^\infty$ through a given representation, whereas Proposition~\ref{prop:lift} shows that the lift of a Galilean action is not unique. Since the functions on $\mathbb{P}\mathcal{H}^\infty$ considered below must be attached to the action $\Phi$ and not to a choice of lift, we show that the space of smooth vectors does not depend on this choice, nor on the Hilbert space realization of the action.

\begin{proposition}[Uniqueness of the space of smooth vectors]\label{prop:unique-smooth}
	Let $\Phi$ be a Galilean action on $\mathbb{P}\mathcal{H}$ and let $\Pi$, $\Pi'$ be two lifts as per Proposition~\ref{prop:lift}. Then
	\begin{enumerate}
		\item $\mathcal{H}^\infty(\Pi)=\mathcal{H}^\infty(\Pi')$, and we denote this space by $\mathcal{H}^\infty_\Phi$. The generators satisfy $\mathbf{P}'=\mathbf{P}$, $\mathbf{K}'=\mathbf{K}$ and $H'=H+\hbar\epsilon$ on $\mathcal{H}^\infty_\Phi$, where $\epsilon\in\mathbb{R}$ is as in Proposition~\ref{prop:lift}.
		\item If $W:\mathcal{H}\to\mathcal{H}'$ is a unitary operator and $\Phi'_\kappa=W\Phi_\kappa W^{-1}$ for all $\kappa\in\mathcal{G}$, then $W\mathcal{H}^\infty_\Phi=\mathcal{H}'^{\infty}_{\Phi'}$.
	\end{enumerate}
\end{proposition}

\begin{proof}
	Item 1. By Proposition~\ref{prop:lift}, $\Pi'(c,\kappa)=e^{i\epsilon s}\,\Pi(c,\kappa)$, where $s$ is the time component of $\kappa$. The orbit maps of $\Pi$ and $\Pi'$ therefore differ by a smooth scalar function of $\kappa$ with smooth inverse, so one is smooth if and only if the other is. For the generators, $S'(s)=e^{i\epsilon s}S(s)=e^{is(H+\hbar\epsilon)/\hbar}$, while $T'=T$ and $B'=B$ since $\alpha_{\mathbf{a}}$ and $\beta_{\mathbf{v}}$ have vanishing time component.
	
	Item 2. The representation $\Pi'':=W\Pi W^{-1}$ of $\widehat{\mathcal{G}}$ on $\mathcal{H}'$ is strongly continuous, it satisfies $\Pi''(z(c))=e^{imc/\hbar}\mathbb{1}$, and $[\Pi''(\hat\kappa)x']=W\Phi_\kappa W^{-1}[x']=\Phi'_\kappa[x']$. Hence $\Pi''$ is a lift of $\Phi'$, and by item 1, $\mathcal{H}'^{\infty}_{\Phi'}=\mathcal{H}^\infty(\Pi'')=W\mathcal{H}^\infty(\Pi)$.
\end{proof}

\noindent In words, $\mathcal{H}^\infty_\Phi$ is determined by the geometric action alone. It depends neither on the phases chosen to lift $\Phi$ to unitary operators nor on the realization of $\mathcal{H}$, and the only ambiguity that survives is the additive constant in $H$, which does not affect the flow on $\mathbb{P}\mathcal{H}$. In Proposition~\ref{prop:smooth-schwartz} this space will be identified with a classical function space.

We now apply this theorem to $\widehat{\mathcal{G}}$ and read the relations of Proposition~\ref{prop:unitary-relations} on $\mathcal{H}^\infty$. Throughout, $h$, $p_j$ and $k_j$ denote the mean value functions on $\mathbb{P}\mathcal{H}^\infty$ of $H$, $P_j$ and $K_j$ as per Definition~\ref{def:mean-value-domain}.

\begin{proposition}[Poisson algebra of the Galilean generators]\label{prop:galilean-poisson}
	Let $\Pi$ be as in Proposition~\ref{prop:lift} and let the brackets be those of Lemma~\ref{lem:pointwise-formulas} with $\mathcal{D}=\mathcal{H}^\infty$. For $i,j\in\{1,2,3\}$, on $\mathbb{P}\mathcal{H}^\infty$,
	\begin{equation}
		\{p_i,p_j\}=0,\quad \{k_i,k_j\}=0,\quad \{k_i,p_j\}=m\,d_{ij},\quad \{p_i,h\}=0,\quad \{k_i,h\}=p_i.
		\label{eq:galilean-poisson}
	\end{equation}
\end{proposition}

\begin{proof}
	By Theorem~\ref{thm:smooth-vectors}, $H$, $P_j$ and $K_j$ restricted to $\mathcal{H}^\infty$ are, up to constant factors, the operators $d\Pi$ of the corresponding Lie algebra elements, and they leave $\mathcal{H}^\infty$ invariant. By Equation~\eqref{eq:pointwise-poisson} it suffices to prove the commutation relations
	\begin{equation}
		[P_i,P_j]=0,\quad [K_i,K_j]=0,\quad [K_i,P_j]=i\hbar m\,d_{ij},\quad [P_i,H]=0,\quad [K_i,H]=i\hbar P_i
	\end{equation}
	on $\mathcal{H}^\infty$. The first two follow from the commutativity of the groups $T$ and $B$ and the fourth from Equation~\eqref{eq:W1}. For $x\in\mathcal{H}^\infty$ the maps $(\mathbf{v},\mathbf{a})\mapsto B(\mathbf{v})T(\mathbf{a})x$ and $(s,\mathbf{v})\mapsto S(-s)B(\mathbf{v})S(s)x$ are smooth. Applying $\partial_{v_i}\partial_{a_j}$ at the origin to Equation~\eqref{eq:W2} yields $\hbar^{-2}K_iP_jx=\tfrac{im}{\hbar}d_{ij}x+\hbar^{-2}P_jK_ix$, which is the third relation. Applying $\partial_s\partial_{v_i}$ at the origin to Equation~\eqref{eq:W3}, the phase gives no contribution since it is quadratic in $\mathbf{v}$, and one obtains $\hbar^{-2}[H,K_i]x=-\tfrac{i}{\hbar}P_ix$, which is the fifth relation.
\end{proof}

\begin{remark}\label{rem:momentum-map}
	By Proposition~\ref{prop:quantum-momentum-map} the momentum functions of the infinitesimal generators of $\Phi$ are $\langle\mu,\xi\rangle=\langle i\hbar\,d\Pi(\xi)\rangle$, which equal $-h$ for the time translation, $\mathbf{a}\cdot\mathbf{p}$ for the translation $\alpha_{\mathbf{a}}$ and $-\mathbf{v}\cdot\mathbf{k}$ for the boost $\beta_{\mathbf{v}}$. The overall signs depend on the active convention of Remark~\ref{rem:active-convention} and are immaterial. The relation $\{k_i,p_j\}=md_{ij}$ shows that the Poisson bracket between two momentum functions is not the momentum function of the Lie bracket, but differs from it by the constant $md_{ij}$. This constant is the value of the momentum function of the central element, in agreement with Example~\ref{ex:u1-phase}, and it is the cocycle that the Bargmann group removes.
\end{remark}

% =====================================================================
\section{The Hamiltonian of the Free Particle}
\label{sec:free-hamiltonian}

We now have all the ingredients to derive the Hamiltonian. The strategy is to compare the time translation group with the explicit one-parameter group $s\mapsto e^{is|\mathbf{P}|^2/2m\hbar}$ built from the momentum operators, and to show that the quotient of the two commutes with the whole representation. Relation~\eqref{eq:W3} between boosts and time translations has no counterpart in the Weyl relations of Section~\ref{sec:weyl-algebra}, and it is the identity that makes the comparison possible. For an irreducible action the quotient is a phase, and the Hamiltonian is determined up to the additive constant already identified in Proposition~\ref{prop:lift}. We first collect two lemmas, one on the covariance of the spectral measure of $\mathbf{P}$ under boosts and one on Schur's lemma in infinite dimensions, which replaces Lemma~\ref{lem:schur} when the representation space is not finite dimensional.

\begin{lemma}[Covariance of the momentum spectral measure]\label{lem:boost-covariance}
	Let $E$ be the projection valued measure of $\mathbf{P}$ as per Definition~\ref{def:generators}. Then for every Borel set $d\subseteq\mathbb{R}^3$, every $\mathbf{v}\in\mathbb{R}^3$ and every $u\in SU(2)$,
	\begin{equation}
		B(\mathbf{v})^{-1}E(d)B(\mathbf{v})=E(d-m\mathbf{v}),\qquad \Pi(\hat\rho_u)\,E(d)\,\Pi(\hat\rho_u)^{-1}=E(R_ud).
	\end{equation}
	Consequently, for every Borel function $f:\mathbb{R}^3\to\mathbb{C}$, as equalities between operators with their domains,
	\begin{equation*}
		B(\mathbf{v})^{-1}f(\mathbf{P})B(\mathbf{v})=f(\mathbf{P}+m\mathbf{v}),\qquad \Pi(\hat\rho_u)\,f(\mathbf{P})\,\Pi(\hat\rho_u)^{-1}=f(R_u^{-1}\mathbf{P}),
	\end{equation*}
	where $f(\mathbf{P}+m\mathbf{v})$ stands for $\int f(\mathbf{p}+m\mathbf{v})\,dE(\mathbf{p})$ and similarly for the second one. If $m\neq0$, then $E(U)\neq0$ for every nonempty open set $U\subseteq\mathbb{R}^3$, \textit{i.e.}, $\operatorname{supp}E=\mathbb{R}^3$.
\end{lemma}

\begin{proof}
	Set $U(\mathbf{t}):=T(-\hbar\mathbf{t})$, so that $E$ is the measure associated with $U$ by Theorem~\ref{thm:snag}. By Equation~\eqref{eq:W2}, $B(\mathbf{v})^{-1}U(\mathbf{t})B(\mathbf{v})=e^{im\mathbf{v}\cdot\mathbf{t}}\,U(\mathbf{t})$. This is the hypothesis of Proposition~\ref{prop:pvm-covariance} with $W=B(\mathbf{v})^{-1}$, $L=\mathbb{1}$ and $\mathbf{c}=m\mathbf{v}$. Here $\phi(\boldsymbol{\lambda})=\boldsymbol{\lambda}+m\mathbf{v}$ and $\phi^{-1}(d)=d-m\mathbf{v}$, which yields the first identity. By Equation~\eqref{eq:W4b}, $\Pi(\hat\rho_u)U(\mathbf{t})\Pi(\hat\rho_u)^{-1}=U(R_u\mathbf{t})$, which is the same hypothesis with $L=R_u$ and $\mathbf{c}=0$. Now $\phi(\boldsymbol{\lambda})=R_u^{T}\boldsymbol{\lambda}=R_u^{-1}\boldsymbol{\lambda}$ and $\phi^{-1}(d)=R_ud$, which yields the second identity. The statement on functions is the second part of Proposition~\ref{prop:pvm-covariance}.
	
	For the last claim, the first identity shows that $E(d)=0$ if and only if $E(d-m\mathbf{v})=0$. Since $m\neq0$, the vectors $m\mathbf{v}$ exhaust $\mathbb{R}^3$. If $E(U)=0$ for a nonempty open set $U$, then all its translates are $E$-null, and $\mathbb{R}^3$ is covered by countably many of them. This would give $E(\mathbb{R}^3)=\mathbb{1}=0$, a contradiction.
\end{proof}

\begin{lemma}[Schur's lemma for a selfadjoint operator]\label{lem:schur-infinite}
	Let $\mathcal{M}\subseteq B(\mathcal{H})$ be irreducible, \textit{i.e.}, the only closed subspaces of $\mathcal{H}$ invariant under every element of $\mathcal{M}$ are $\{0\}$ and $\mathcal{H}$. Let $A$ be a selfadjoint operator on $\mathcal{H}$ such that $e^{itA}X=Xe^{itA}$ for all $X\in\mathcal{M}$ and $t\in\mathbb{R}$. Then $A=a\,\mathbb{1}$ for some $a\in\mathbb{R}$.
\end{lemma}

\begin{proof}
	Let $E_A$ be the spectral measure of $A$ given by Theorem~\ref{thm:spectral}, so that $E_A$ is the measure associated with $t\mapsto e^{itA}$ by Theorem~\ref{thm:snag}. By Proposition~\ref{prop:pvm-commutation}, every $X\in\mathcal{M}$ commutes with every $E_A(d)$. Hence $E_A(d)\mathcal{H}$ is a closed subspace invariant under $\mathcal{M}$, and irreducibility gives $E_A(d)\in\{0,\mathbb{1}\}$ for every Borel set $d\subseteq\mathbb{R}$. Set
	\begin{equation*}
		a:=\sup\big\{\lambda\in\mathbb{R}\ \big|\ E_A((-\infty,\lambda))=0\big\}.
	\end{equation*}
	The set is nonempty and bounded above, because $E_A((-\infty,\lambda))$ converges strongly to $0$ for $\lambda\to-\infty$ and to $\mathbb{1}$ for $\lambda\to+\infty$ and takes only the values $0$ and $\mathbb{1}$. By monotonicity, $E_A((-\infty,\lambda))=0$ for $\lambda<a$ and $E_A((-\infty,\lambda))=\mathbb{1}$ for $\lambda>a$. By $\sigma$-additivity, $E_A(\{a\})=\lim_{\lambda\downarrow a}E_A((-\infty,\lambda))-\lim_{\lambda\uparrow a}E_A((-\infty,\lambda))=\mathbb{1}$, hence $A=\int\lambda\,dE_A(\lambda)=a\,\mathbb{1}$.
\end{proof}

\noindent This is the infinite dimensional version of Lemma~\ref{lem:schur}, and it is the only place where irreducibility of the action enters the proof of Theorem~\ref{thm:free-hamiltonian}.

\begin{theorem}[Hamiltonian of the free Galilean particle]\label{thm:free-hamiltonian}
	Let $\Phi$ be a Galilean action on $\mathbb{P}\mathcal{H}$ as per Definition~\ref{def:galilean-action}, with lift $\Pi$ and mass parameter $m\neq0$ as in Proposition~\ref{prop:lift}, and let $H$, $\mathbf{P}$ be as per Definition~\ref{def:generators}. Then
	\begin{enumerate}
		\item $S(s)=e^{is|\mathbf{P}|^2/2m\hbar}\,C(s)$, where $C(s)=e^{isH_{\mathrm{int}}/\hbar}$ for a selfadjoint operator $H_{\mathrm{int}}$ whose spectral projections commute with $\Pi(\widehat{\mathcal{G}})$. Hence $H$ is the closure of $\tfrac{1}{2m}|\mathbf{P}|^2+H_{\mathrm{int}}$ on $D(|\mathbf{P}|^2)\cap D(H_{\mathrm{int}})$.
		\item If $\Phi$ is irreducible as per Definition~\ref{def:irreducible-action}, then $H_{\mathrm{int}}=E_0\mathbb{1}$ for some $E_0\in\mathbb{R}$, and
		\begin{equation}
			H=\frac{|\mathbf{P}|^2}{2m}+E_0\,\mathbb{1}, \qquad D(H)=D(|\mathbf{P}|^2).
			\label{eq:free-hamiltonian}
		\end{equation}
		\item $\sigma(|\mathbf{P}|^2)=[0,\infty)$. Hence, in the irreducible case, $\sigma(H)=[E_0,\infty)$ if $m>0$ and $\sigma(H)=(-\infty,E_0]$ if $m<0$, so that $H$ is bounded below if and only if $m>0$.
	\end{enumerate}
\end{theorem}

\begin{proof}
	\emph{Idea.} The group $S$ and the explicit group $S_0(s):=e^{is|\mathbf{P}|^2/2m\hbar}$ built from the momentum operators obey the same commutation relation with the boosts, namely Relation~\eqref{eq:W3}. Their quotient $C(s):=S_0(s)^{-1}S(s)$ therefore commutes with the whole representation, and for an irreducible action it can only be a phase. The proof follows this plan in six steps. We write $E$ for the measure of Definition~\ref{def:generators}, so that $S_0(s)=\int e^{is|\mathbf{p}|^2/2m\hbar}\,dE(\mathbf{p})$ in the sense of Definition~\ref{def:pvm-integral}.
	
	\medskip
	\noindent\textbf{Step 1: the free group obeys Relation~\eqref{eq:W3}.} We show that
	\begin{equation*}
		S_0(-s)\,B(\mathbf{v})\,S_0(s)=e^{-ims|\mathbf{v}|^2/2\hbar}\,B(\mathbf{v})\,T(s\mathbf{v}).
	\end{equation*}
	Inserting $B(\mathbf{v})B(\mathbf{v})^{-1}$ and applying Lemma~\ref{lem:boost-covariance} to the bounded function $\mathbf{p}\mapsto e^{-is|\mathbf{p}|^2/2m\hbar}$,
	\begin{align*}
		S_0(-s)B(\mathbf{v})S_0(s)
		&=B(\mathbf{v})\,\big[B(\mathbf{v})^{-1}S_0(-s)B(\mathbf{v})\big]\,S_0(s)\\
		&=B(\mathbf{v})\,\exp\!\Big(-\frac{is}{2m\hbar}\big(|\mathbf{P}+m\mathbf{v}|^2-|\mathbf{P}|^2\big)\Big)\\
		&=B(\mathbf{v})\,\exp\!\Big(-\frac{is}{\hbar}\,\mathbf{v}\cdot\mathbf{P}\Big)\,e^{-ims|\mathbf{v}|^2/2\hbar}\\
		&=e^{-ims|\mathbf{v}|^2/2\hbar}\,B(\mathbf{v})\,T(s\mathbf{v}).
	\end{align*}
	The third line uses $|\mathbf{P}+m\mathbf{v}|^2-|\mathbf{P}|^2=2m\,\mathbf{v}\cdot\mathbf{P}+m^2|\mathbf{v}|^2$ and the multiplicativity of the functional calculus, and the last line uses $T(s\mathbf{v})=e^{-is\mathbf{v}\cdot\mathbf{P}/\hbar}$. With the coefficient $1/2m$, the shift of $\mathbf{P}$ by $m\mathbf{v}$ produces in the kinetic energy a term linear in $\mathbf{P}$, which is the translation $T(s\mathbf{v})$, and a constant term, which is the phase.
	
	\medskip
	\noindent\textbf{Step 2: $C(s)$ commutes with translations and boosts.} Both $S(s)$ and $S_0(s)$ commute with $T(\mathbf{a})$, the first by Relation~\eqref{eq:W1} and the second because it is a function of $\mathbf{P}$. Hence $C(s)$ commutes with $T(\mathbf{a})$. For the boosts, Step 1 with $s$ replaced by $-s$ gives $S_0(s)B(\mathbf{v})S_0(-s)=e^{ims|\mathbf{v}|^2/2\hbar}B(\mathbf{v})T(-s\mathbf{v})$, and then Relations~\eqref{eq:W3} and~\eqref{eq:W1} yield
	\begin{align*}
		C(s)^{-1}B(\mathbf{v})\,C(s)
		&=S(-s)\,\big[S_0(s)B(\mathbf{v})S_0(-s)\big]\,S(s)\\
		&=e^{ims|\mathbf{v}|^2/2\hbar}\,\big[S(-s)B(\mathbf{v})S(s)\big]\,T(-s\mathbf{v})\\
		&=e^{ims|\mathbf{v}|^2/2\hbar}\,e^{-ims|\mathbf{v}|^2/2\hbar}\,B(\mathbf{v})\,T(s\mathbf{v})\,T(-s\mathbf{v})
		=B(\mathbf{v}).
	\end{align*}
	The two phases cancel, and this is where it matters that $S_0$ and $S$ satisfy the same relation.
	
	\medskip
	\noindent\textbf{Step 3: $C$ is a one-parameter group and defines $H_{\mathrm{int}}$.} By Step 2 and Proposition~\ref{prop:pvm-commutation}, the operator $S(s')$ commutes with every $E(d)$, hence with $S_0(s)$. Therefore
	\begin{equation*}
		C(s+s')=S_0(-s)S_0(-s')S(s)S(s')=S_0(-s)S(s)\,S_0(-s')S(s')=C(s)\,C(s'),
	\end{equation*}
	and $C$ is strongly continuous. Theorem~\ref{thm:stone} yields a selfadjoint operator $H_{\mathrm{int}}$ with $C(s)=e^{isH_{\mathrm{int}}/\hbar}$. Write $H_0:=\frac{1}{2m}|\mathbf{P}|^2$, so that $S_0(s)=e^{isH_0/\hbar}$. By Step 2 and Proposition~\ref{prop:pvm-commutation}, $C(s)$ commutes with every $E(d)$, hence with $S_0(s')$. Applying Proposition~\ref{prop:pvm-commutation} twice, first to $X=S_0(s')$ and the group $C$, then to $X=E_{H_{\mathrm{int}}}(d)$ and the group $S_0$, shows that $H_0$ and $H_{\mathrm{int}}$ commute strongly. Theorem~\ref{thm:joint-spectral} then implies that $S(s)=S_0(s)C(s)$ is generated by the closure of $H_0+H_{\mathrm{int}}$ on $D(H_0)\cap D(H_{\mathrm{int}})$, and the uniqueness in Theorem~\ref{thm:stone} identifies this closure with $H$.
	
	\medskip
	\noindent\textbf{Step 4: $C(s)$ commutes with the whole representation.} By Lemma~\ref{lem:boost-covariance}, $\Pi(\hat\rho_u)\,|\mathbf{P}|^2\,\Pi(\hat\rho_u)^{-1}=|R_u^{-1}\mathbf{P}|^2=|\mathbf{P}|^2$, so $\Pi(\hat\rho_u)$ commutes with $S_0(s)$. It commutes with $S(s)$ by Relation~\eqref{eq:W4a}, hence with $C(s)$. Together with Steps 2 and 3 and the trivial commutation with scalars, $C(s)$ commutes with all the generators listed in Proposition~\ref{prop:unitary-relations}, and therefore with $\Pi(\widehat{\mathcal{G}})$. By Proposition~\ref{prop:pvm-commutation} applied to the group $C$, the spectral projections of $H_{\mathrm{int}}$ commute with $\Pi(\widehat{\mathcal{G}})$ as well. This proves item 1.
	
	\medskip
	\noindent\textbf{Step 5: the irreducible case.} The set $\Pi(\widehat{\mathcal{G}})$ is irreducible. Indeed, if $\mathcal{K}\subseteq\mathcal{H}$ is a closed subspace invariant under $\Pi(\widehat{\mathcal{G}})$, then Equation~\eqref{eq:lift} gives $\Phi_\kappa(\mathbb{P}\mathcal{K})\subseteq\mathbb{P}\mathcal{K}$ for all $\kappa\in\mathcal{G}$, and Definition~\ref{def:irreducible-action} forces $\mathcal{K}\in\{\{0\},\mathcal{H}\}$. Since $e^{itH_{\mathrm{int}}/\hbar}=C(t)$ commutes with $\Pi(\widehat{\mathcal{G}})$ by Step 4, Lemma~\ref{lem:schur-infinite} applied to $A=H_{\mathrm{int}}/\hbar$ yields $H_{\mathrm{int}}=E_0\mathbb{1}$ for some $E_0\in\mathbb{R}$. Then
	\begin{equation*}
		S(s)=S_0(s)\,e^{isE_0/\hbar}=\exp\!\Big(\frac{is}{\hbar}\Big(\frac{|\mathbf{P}|^2}{2m}+E_0\Big)\Big),
	\end{equation*}
	and the uniqueness in Theorem~\ref{thm:stone} yields Equation~\eqref{eq:free-hamiltonian}, with $D(H)=D(|\mathbf{P}|^2)$ because the addition of a constant does not change the domain. This proves item 2.
	
	\medskip
	\noindent\textbf{Step 6: the spectrum.} By Lemma~\ref{lem:boost-covariance} the support of $E$ is $\mathbb{R}^3$. Item 5 of Proposition~\ref{prop:pvm-calculus} applied to the continuous function $\mathbf{p}\mapsto|\mathbf{p}|^2$ gives $\sigma(|\mathbf{P}|^2)=\overline{\{|\mathbf{p}|^2\mid\mathbf{p}\in\mathbb{R}^3\}}=[0,\infty)$. The statement on $\sigma(H)$ follows from Equation~\eqref{eq:free-hamiltonian}, since $m>0$ gives $H\ge E_0$ with spectrum $[E_0,\infty)$ and $m<0$ gives $H\le E_0$ with spectrum $(-\infty,E_0]$. This proves item 3.
\end{proof}

\noindent The constant $E_0$ is an internal energy which is not determined by Galilean invariance, and by Proposition~\ref{prop:lift} it can be shifted to any value by changing the lift, so that it has no effect on the flow on $\mathbb{P}\mathcal{H}$.

\begin{remark}
	Each hypothesis of Theorem~\ref{thm:free-hamiltonian} has a specific role. Continuity of the action is needed to apply Theorem~\ref{thm:bargmann}. The condition $m\neq0$ is used in Lemma~\ref{lem:boost-covariance} and it is the point at which the kinetic term acquires the factor $1/2m$. Irreducibility enters only in Step 5, and without it item 1 still holds, with $H_{\mathrm{int}}$ describing internal degrees of freedom that commute with the whole Galilean representation. The additive constant is the same ambiguity as in Proposition~\ref{prop:kahler-isomorphisms}.
\end{remark}

The result can be rewritten in the language of Chapter~\ref{ch:projective}. The kinetic term of Equation~\eqref{eq:free-hamiltonian} becomes the sum of the classical expression in the momentum functions and a term built from the Riemann bracket, which accounts for the quantum dispersion.

\begin{corollary}[Geometric form of the free Hamiltonian]\label{cor:geometric-hamiltonian}
	In the hypotheses of Theorem~\ref{thm:free-hamiltonian}, item 2, for every $[x]\in\mathbb{P}\mathcal{H}^\infty$
	\begin{equation}
		h([x])=\frac{1}{2m}\sum_{j=1}^3 p_j\circ_{\hbar} p_j+E_0=\frac{|\mathbf{p}|^2}{2m}+\frac{\hbar}{4m}\sum_{j=1}^3((p_j,p_j))+E_0,
		\label{eq:geometric-hamiltonian}
	\end{equation}
	where $\circ_{\hbar}$ is the pointwise product $f\circ_{\hbar} l=\tfrac{\hbar}{2}((f,l))+fl$ of Definition~\ref{def:kahler-bracket}. Equivalently, $h=|\mathbf{p}|^2/2m+\tfrac{1}{2m}\sum_j(d p_j)^2+E_0$, with the dispersions of Definition~\ref{def:dispersion}.
\end{corollary}

\begin{proof}
	Since $\mathcal{H}^\infty\subseteq D(|\mathbf{P}|^2)$ and $P_j\mathcal{H}^\infty\subseteq\mathcal{H}^\infty$, Equation~\eqref{eq:free-hamiltonian} yields $h=\tfrac{1}{2m}\sum_j\langle P_j^2\rangle+E_0$ on $\mathbb{P}\mathcal{H}^\infty$. Equation~\eqref{eq:pointwise-riemann} with $A=B=P_j$ gives $\langle P_j^2\rangle=\tfrac{\hbar}{2}((p_j,p_j))+p_j^2$, and the dispersion form follows from Equation~\eqref{eq:dispersion-relation}.
\end{proof}

\noindent The first term of Equation~\eqref{eq:geometric-hamiltonian} is the kinetic energy of a classical particle of mass $m$ and momentum $\mathbf{p}$, as in Chapter~\ref{ch:classical}. The second term has no classical counterpart, and it is the subject of the following remark and proposition.

\begin{remark}[Meaning of the dispersion term]\label{rem:dispersion-term}
	The function $p_j$ is the mean momentum of the state, while $(d p_j)^2$ is the variance of the momentum in the same state. Equation~\eqref{eq:geometric-hamiltonian} is therefore the identity $\langle|\mathbf{P}|^2\rangle=|\langle\mathbf{P}\rangle|^2+\sum_j(d p_j)^2$, divided by $2m$ and shifted by $E_0$. The energy of a free particle splits into the kinetic energy of the motion of the mean momentum and the energy stored in the spread of the momentum around it. The second part is a genuinely quantum contribution, since in a classical state the momentum is a point and the spread vanishes. In geometric terms $(d p_j)^2=\tfrac{\hbar}{2}((p_j,p_j))$ is proportional to the squared length of the Hamiltonian vector field of $p_j$, by Equation~\eqref{eq:pointwise-riemann}. The next proposition shows that this length never vanishes and that it is bounded below by the uncertainty relation, so that localizing a free particle always costs kinetic energy. Example~\ref{ex:gaussian-packet} makes both statements explicit on a Gaussian packet.
\end{remark}

\begin{proposition}[Positivity of the dispersion]\label{prop:dispersion-positive}
	In the hypotheses of Theorem~\ref{thm:free-hamiltonian}, item 2, let $[x]\in\mathbb{P}\mathcal{H}^\infty$ and set $d p_j([x]):=\sqrt{\hbar/2}\,\sqrt{((p_j,p_j))([x])}$ as in Definition~\ref{def:dispersion}, and in the same way $d k_j([x])$. Then, for $j\in\{1,2,3\}$,
	\begin{enumerate}
		\item $d p_j([x])>0$,
		\item $d k_j([x])\,d p_j([x])\ge\frac{\hbar}{2}\,|m|$. In particular, for $q_j:=k_j/m$ and $d q_j:=d k_j/|m|$ one has $d q_j\,d p_j\ge\hbar/2$.
	\end{enumerate}
\end{proposition}

\begin{proof}
	Let $x$ be a unit representative of $[x]$. Equation~\eqref{eq:pointwise-riemann} with $A=B=P_j$ gives $(d p_j)^2=\langle P_j^2\rangle-p_j^2=\|(P_j-p_j([x]))x\|^2$, and the same holds for $K_j$.
	
	Item 1. If $d p_j([x])=0$, then $P_jx=p_j([x])x$, so that $x$ is an eigenvector of $P_j$. By Lemma~\ref{lem:boost-covariance}, with $\mathbf{v}=ve_j$, one has $B(ve_j)^{-1}P_jB(ve_j)=P_j+mv$ as operators with equal domains. Hence $B(ve_j)x\in D(P_j)$ and $P_jB(ve_j)x=(p_j([x])+mv)\,B(ve_j)x$ for every $v\in\mathbb{R}$. Since $m\neq0$, these are eigenvectors of the selfadjoint operator $P_j$ with pairwise distinct eigenvalues, hence pairwise orthogonal unit vectors. There are uncountably many of them, which contradicts the separability of $\mathcal{H}$.
	
	Item 2. Take $A=K_j$ and $B=P_j$, and let $\eta_A,\eta_B$ be the vectors of Lemma~\ref{lem:pointwise-formulas}. Since $\eta_B\perp x$ and $(A-f_A([x]))x=i\hbar\,\eta_A$, one has $(\eta_B,Ax)=(\eta_B,(A-f_A([x]))x)=i\hbar\,(\eta_B,\eta_A)$, hence $\{f_A,f_B\}([x])=2\hbar\,\mathrm{Im}(\eta_A,\eta_B)$. The Cauchy--Schwarz inequality and $\|\eta_A\|=d a/\hbar$, with $d a:=\|(A-f_A([x]))x\|$ and similarly for $B$, give
	\begin{equation*}
		|\{f_A,f_B\}([x])|\le2\hbar\,\|\eta_A\|\,\|\eta_B\|=\frac{2}{\hbar}\,d a\,d b.
	\end{equation*}
	Proposition~\ref{prop:galilean-poisson} gives $\{k_j,p_j\}=m$, which yields the claim.
\end{proof}

\noindent By item 2, for $m>0$ the dispersion term of Equation~\eqref{eq:geometric-hamiltonian} is bounded below by $\frac{\hbar^2}{8m}\sum_j(d q_j)^{-2}$. A state localized in a region of size $d q_j$ therefore carries a kinetic energy of order $\hbar^2/m(d q_j)^2$, which grows as the localization increases.

We now derive the equations of motion on $\mathbb{P}\mathcal{H}$. They reproduce uniform rectilinear motion, and they exhibit the explicitly time dependent integral of motion associated with the boost.

\begin{corollary}[Uniform motion]\label{cor:uniform-motion}
	In the hypotheses of Theorem~\ref{thm:free-hamiltonian}, item 2, let $[x]\in\mathbb{P}\mathcal{H}^\infty$ and let $\Phi^{\mathrm{S}}_t[x]=[e^{-itH/\hbar}x]$ be the Schrödinger flow. Then $\Phi^{\mathrm{S}}_t[x]\in\mathbb{P}\mathcal{H}^\infty$ and, writing $q_j:=k_j/m$,
	\begin{equation}
		p_j(\Phi^{\mathrm{S}}_t[x])=p_j([x]),\qquad q_j(\Phi^{\mathrm{S}}_t[x])=q_j([x])+\frac{t}{m}\,p_j([x]).
	\end{equation}
	In particular the function $k_j-t\,p_j$ is conserved along the flow, although $k_j$ is not.
\end{corollary}

\begin{proof}
	The subspace $\mathcal{H}^\infty$ is invariant under $S(t)$ by Theorem~\ref{thm:smooth-vectors}. For $A\in\{P_j,K_j\}$ and $x_t=e^{-itH/\hbar}x$, differentiating $(x_t,Ax_t)$ yields $\tfrac{d}{dt}\langle A\rangle([x_t])=\langle\tfrac{i}{\hbar}[H,A]\rangle([x_t])=\{\langle A\rangle,h\}([x_t])$ by Equation~\eqref{eq:pointwise-poisson}. Proposition~\ref{prop:galilean-poisson} gives $\dot p_j=0$ and $\dot k_j=p_j$, and integration yields the claim.
\end{proof}

The geometric statements above do not use the structure of $\mathcal{H}$ beyond Theorem~\ref{thm:free-hamiltonian}. To connect the result with the Schrödinger representation of Section~\ref{sec:weyl-algebra}, and to identify the role of spin, we combine the theorem with the uniqueness theorem for the Weyl relations.

\begin{corollary}[Schrödinger realization]\label{cor:schrodinger-realization}
	In the hypotheses of Theorem~\ref{thm:free-hamiltonian}, item 2, there exist a $j\in\{0,\tfrac{1}{2},1,\dots\}$ and a unitary operator $\mathcal{H}\to L^2(\mathbb{R}^3)\otimes\mathbb{C}^{2j+1}$ under which
	\begin{equation}
		(T(\mathbf{a})\psi)(\mathbf{x})=\psi(\mathbf{x}-\mathbf{a}),\qquad (B(\mathbf{v})\psi)(\mathbf{x})=e^{im\mathbf{v}\cdot\mathbf{x}/\hbar}\psi(\mathbf{x}),\qquad \Pi(\hat\rho_u)=\Pi_{\mathrm{orb}}(u)\otimes d_j(u),
	\end{equation}
	where $\Pi_{\mathrm{orb}}(u)\psi=\psi(R_u^{-1}\,\cdot)$ and $d_j$ is the irreducible representation of $SU(2)$ of dimension $2j+1$ acting on $\mathbb{C}^{2j+1}$, with $\psi\in L^2(\mathbb{R}^3,\mathbb{C}^{2j+1})$. In this realization
	\begin{equation}
		H=-\frac{\hbar^2}{2m}\,d\otimes\mathbb{1}+E_0\,\mathbb{1}, \qquad \mathbf{P}=-i\hbar\nabla\otimes\mathbb{1}, \qquad \mathbf{K}=m\mathbf{x}\otimes\mathbb{1}.
	\end{equation}
\end{corollary}

\begin{proof}
	\emph{Strategy.} Relation~\eqref{eq:W2} makes $(B,T)$ a representation of the Weyl relations, which fixes the action of $T$ and $B$ up to a multiplicity space $\mathcal{K}$. The rotations and the time translations then act on $\mathcal{K}$ in a controlled way, and irreducibility forces $\mathcal{K}$ to carry a single irreducible representation of $SU(2)$.
	
	\medskip
	\noindent\textbf{Step 1: the tensor structure.} By Remark~\ref{rem:weyl}, the pair $U(\boldsymbol{\alpha})=B(\boldsymbol{\alpha}/m)$, $V(\boldsymbol{\beta})=T(-\boldsymbol{\beta})$ is a regular representation of the Weyl relations in three degrees of freedom, as per Definition~\ref{def:regular-representation}, on the separable space $\mathcal{H}$. The von Neumann uniqueness theorem, Theorem~\ref{thm:von-neumann}, concerns irreducible representations, and an arbitrary regular representation on a separable space is a direct sum of irreducible ones \cite[TODO-location]{TODO-RS03}. Hence there is a unitary operator $\mathcal{H}\to L^2(\mathbb{R}^3)\otimes\mathcal{K}$, with $\mathcal{K}$ a separable Hilbert space, under which $U$ and $V$ act as the Schrödinger representation of Definition~\ref{def:schrodinger-representation} on the first factor and trivially on $\mathcal{K}$. This gives the stated form of $T(\mathbf{a})$ and $B(\mathbf{v})$ and, by Theorem~\ref{thm:snag}, the stated $\mathbf{P}$ and $\mathbf{K}$.
	
	\medskip
	\noindent\textbf{Step 2: the rotations.} The operator $\Pi_{\mathrm{orb}}(u)\otimes\mathbb{1}$ transforms $T(\mathbf{a})$ into $T(R_u\mathbf{a})$ and $B(\mathbf{v})$ into $B(R_u\mathbf{v})$, as one checks on the explicit formulas, and so does $\Pi(\hat\rho_u)$ by Relations~\eqref{eq:W4b} and~\eqref{eq:W4c}. Hence $(\Pi_{\mathrm{orb}}(u)\otimes\mathbb{1})^{-1}\Pi(\hat\rho_u)$ commutes with all $T(\mathbf{a})$ and $B(\mathbf{v})$. By irreducibility of the Schrödinger representation these operators generate $B(L^2(\mathbb{R}^3))\otimes\mathbb{1}$ as a von Neumann algebra, whose commutant is $\mathbb{1}\otimes B(\mathcal{K})$. Therefore $\Pi(\hat\rho_u)=\Pi_{\mathrm{orb}}(u)\otimes d(u)$ for a unitary operator $d(u)$ on $\mathcal{K}$. Both $\Pi$ and $\Pi_{\mathrm{orb}}$ are representations of $SU(2)$, and $\mathbb{1}\otimes d(u)$ commutes with $\Pi_{\mathrm{orb}}(u')\otimes\mathbb{1}$, so that $d(uu')=d(u)d(u')$. The map $d$ is strongly continuous, hence a unitary representation of $SU(2)$ on $\mathcal{K}$.
	
	\medskip
	\noindent\textbf{Step 3: irreducibility of $d$.} By Theorem~\ref{thm:free-hamiltonian}, $S(s)=e^{is|\mathbf{P}|^2/2m\hbar}e^{isE_0/\hbar}$, and $|\mathbf{P}|^2=-\hbar^2d\otimes\mathbb{1}$ acts on the first factor only. Let $\mathcal{K}_1\subseteq\mathcal{K}$ be a closed subspace invariant under $d$. Then $L^2(\mathbb{R}^3)\otimes\mathcal{K}_1$ is invariant under $T$, $B$, $S$, $\Pi(\hat\rho_u)$ and the scalars, hence under $\Pi(\widehat{\mathcal{G}})$ by Proposition~\ref{prop:unitary-relations}. Irreducibility of $\Phi$, in the form obtained in Step 5 of the proof of Theorem~\ref{thm:free-hamiltonian}, gives $\mathcal{K}_1\in\{\{0\},\mathcal{K}\}$. Thus $d$ is an irreducible strongly continuous unitary representation of the compact group $SU(2)$. Such a representation is finite dimensional and isomorphic to the representation $d_j$ of dimension $2j+1$ for a unique $j\in\{0,\tfrac{1}{2},1,\dots\}$ \cite[TODO-location]{TODO-Hal03}, so that $\mathcal{K}\cong\mathbb{C}^{2j+1}$.
	
	\medskip
	\noindent\textbf{Step 4: the formulas.} The expressions for $\mathbf{P}$ and $\mathbf{K}$ follow from Definition~\ref{def:generators} and the form of $T$ and $B$ in Step 1, and the expression for $H$ follows from Equation~\eqref{eq:free-hamiltonian}. In particular $B(\mathbf{v})^{-1}\mathbf{P}B(\mathbf{v})=\mathbf{P}+m\mathbf{v}$, so that boosts increase the momentum by $m\mathbf{v}$, as stated in Remark~\ref{rem:active-convention}.
\end{proof}

\noindent It descends that the free Schrödinger equation $i\hbar\,\partial_t\psi=-\tfrac{\hbar^2}{2m}d\psi$ is recovered with no reference to a classical Hamiltonian, and that spin enters only through the representation $d_j$, without affecting the form of $H$. An elementary Galilean system is thus labelled by the pair $(m,j)$ and by the irrelevant constant $E_0$.

The identification just obtained also pins down the space of smooth vectors of Section~\ref{sec:galilean-functions}, which was defined abstractly and shown to be unique in Proposition~\ref{prop:unique-smooth}. In the realization above it is a classical function space, which makes the domain of the Galilean mean value functions completely explicit.

\begin{proposition}[Smooth vectors of the free particle]\label{prop:smooth-schwartz}
	In the realization of Corollary~\ref{cor:schrodinger-realization},
	\begin{equation*}
		\mathcal{H}^\infty_\Phi=\mathcal{S}\big(\mathbb{R}^3;\mathbb{C}^{2j+1}\big)=\mathcal{S}(\mathbb{R}^3)\otimes\mathbb{C}^{2j+1},
	\end{equation*}
	where $\mathcal{S}(\mathbb{R}^3)$ is the Schwartz space of Definition~\ref{def:schwartz}. By Proposition~\ref{prop:unique-smooth}, this identification does not depend on the lift of $\Phi$.
\end{proposition}

\begin{proof}
	Write $N=2j+1$.
	
	\emph{Inclusion $\subseteq$.} Let $\psi\in\mathcal{H}^\infty_\Phi$. By Theorem~\ref{thm:smooth-vectors}, $\mathcal{H}^\infty_\Phi$ is invariant under $d\Pi(\xi)$ for every $\xi$ in the Lie algebra of $\widehat{\mathcal{G}}$, and by Definition~\ref{def:generators} these operators restrict to multiples of $P_j=-i\hbar\partial_j$ and $K_j=mx_j$. By induction, $\psi$ belongs to the domain of every finite product of these operators, taken in any order, and the result is again in $L^2$. In particular $K^\beta P^\alpha\psi\in L^2$ for all multi-indices, that is, $\mathbf{x}^\beta\partial^\alpha\psi\in L^2$ for the weak derivatives of each component of $\psi$. Item 4 of Proposition~\ref{prop:schwartz-properties} gives $\psi\in\mathcal{S}(\mathbb{R}^3;\mathbb{C}^N)$.
	
	\emph{Inclusion $\supseteq$.} Let $\psi\in\mathcal{S}(\mathbb{R}^3;\mathbb{C}^N)$. By item 6 of Lemma~\ref{lem:galilei-relations}, $\Pi(c,\kappa)\psi=e^{imc/\hbar}\,S(s)T(\mathbf{a})B(\mathbf{v})\Pi(\hat\rho_u)\psi$ for $\kappa=(s,\mathbf{a},\mathbf{v},u)$, and $(c,s,\mathbf{a},\mathbf{v},u)$ are global coordinates on $\widehat{\mathcal{G}}$. It suffices to show that each of the four maps $s\mapsto S(s)\phi$, $\mathbf{a}\mapsto T(\mathbf{a})\phi$, $\mathbf{v}\mapsto B(\mathbf{v})\phi$ and $u\mapsto\Pi(\hat\rho_u)\phi$ is a smooth action on the Fr\'echet space $\mathcal{S}(\mathbb{R}^3;\mathbb{C}^N)$. Indeed, their composition is then smooth with values in $\mathcal{S}$, hence in $L^2$ by the continuity of the inclusion. Now $T(\mathbf{a})$ and $B(\mathbf{v})$ are translations and multiplications by $e^{im\mathbf{v}\cdot\mathbf{x}/\hbar}$, with derivatives $\partial_{a_j}T(\mathbf{a})\phi=T(\mathbf{a})(-\partial_j\phi)$ and $\partial_{v_j}B(\mathbf{v})\phi=\tfrac{im}{\hbar}\,x_jB(\mathbf{v})\phi$. The rotations act through $\phi\mapsto\phi\circ R_u^{-1}$ and the matrices $d_j(u)$. The group $S(s)$ acts in Fourier variables as the multiplication by $e^{is(|\mathbf{p}|^2/2m+E_0)/\hbar}$, whose derivatives are polynomially bounded. By items 2 and 3 of Proposition~\ref{prop:schwartz-properties} each of these maps is smooth into $\mathcal{S}$, and each derivative is a continuous operator on $\mathcal{S}$. Hence $\psi\in\mathcal{H}^\infty_\Phi$.
\end{proof}

\noindent In particular, the Galilean mean value functions $h$, $p_j$ and $k_j$ are defined on the rays of Schwartz functions, on which $Q_j$, $P_j$ and $d$ act as the familiar differential operators.

The following example shows the content of Equation~\eqref{eq:geometric-hamiltonian} and of Proposition~\ref{prop:dispersion-positive} on the simplest family of states with nonzero momentum dispersion.

\begin{example}[Gaussian wave packet]\label{ex:gaussian-packet}
	In the realization of Corollary~\ref{cor:schrodinger-realization}, fix $\sigma>0$, $\mathbf{p}_0\in\mathbb{R}^3$ and a unit vector $\chi\in\mathbb{C}^{2j+1}$, and consider
	\begin{equation*}
		\psi(\mathbf{x})=(2\pi\sigma^2)^{-3/4}\,e^{i\mathbf{p}_0\cdot\mathbf{x}/\hbar}\,e^{-|\mathbf{x}|^2/4\sigma^2}\,\chi.
	\end{equation*}
	It is normalized and it belongs to $\mathcal{S}(\mathbb{R}^3;\mathbb{C}^{2j+1})$ by item 1 of Proposition~\ref{prop:schwartz-properties}, hence $[\psi]\in\mathbb{P}\mathcal{H}^\infty_\Phi$ by Proposition~\ref{prop:smooth-schwartz}. The position density is $(2\pi\sigma^2)^{-3/2}e^{-|\mathbf{x}|^2/2\sigma^2}$, so that $\langle x_j\rangle=0$ and $\langle x_j^2\rangle=\sigma^2$. For the momentum, $P_j\psi=\big(p_{0,j}+\tfrac{i\hbar}{2\sigma^2}x_j\big)\psi$, hence
	\begin{equation*}
		p_j([\psi])=p_{0,j},\qquad (P_j-p_{0,j})\psi=\frac{i\hbar}{2\sigma^2}\,x_j\psi,\qquad d p_j([\psi])^2=\frac{\hbar^2}{4\sigma^4}\,\sigma^2=\frac{\hbar^2}{4\sigma^2}.
	\end{equation*}
	Thus $d p_j=\hbar/2\sigma$, $d q_j=\sigma$ and $d q_jd p_j=\hbar/2$, so that the bound of Proposition~\ref{prop:dispersion-positive} is saturated. The Riemann bracket is $((p_j,p_j))([\psi])=2(d p_j)^2/\hbar=\hbar/2\sigma^2$, and Equation~\eqref{eq:geometric-hamiltonian} yields
	\begin{equation*}
		h([\psi])=\frac{|\mathbf{p}_0|^2}{2m}+\frac{3\hbar^2}{8m\sigma^2}+E_0,
	\end{equation*}
	which agrees with the direct evaluation of $(\psi,H\psi)$. The energy exceeds the classical value $|\mathbf{p}_0|^2/2m+E_0$ by $3\hbar^2/8m\sigma^2$, and the excess grows as the packet is squeezed. It vanishes only in the limit $\sigma\to\infty$, in which $d p_j\to0$ but the packet ceases to be normalizable and leaves $\mathbb{P}\mathcal{H}$. By Corollary~\ref{cor:uniform-motion}, the mean position evolves as $q_j(t)=t\,p_{0,j}/m$.
\end{example}

\begin{remark}[Poisson reading of the proof]\label{rem:poisson-shadow}
	Theorem~\ref{thm:free-hamiltonian} is the operator version of a statement on $\mathbb{P}\mathcal{H}^\infty$ that follows from Proposition~\ref{prop:galilean-poisson}. Setting $q_i=k_i/m$ one has $\{q_i,p_j\}=d_{ij}$, and $\{k_i,\langle P_j^2\rangle\}=\langle\tfrac{1}{i\hbar}[K_i,P_j^2]\rangle=2m\,d_{ij}\,p_i$. Hence the function $h_{\mathrm{int}}:=h-\tfrac{1}{2m}\sum_j\langle P_j^2\rangle$ satisfies $\{k_i,h_{\mathrm{int}}\}=\{p_i,h_{\mathrm{int}}\}=0$, so it is invariant under the translations and the boosts. For an irreducible action the only such functions are constants, which is the content of Step 5.
\end{remark}

\begin{remark}[Outlook]\label{rem:outlook}
	The mechanism used in this chapter is that group relations between boosts and time translations fix the form of the energy. In the Galilean case the mass is a parameter of a central extension, and the phase in Relation~\eqref{eq:W2} is what makes $\mathbf{K}/m$ conjugate to $\mathbf{P}$. For the universal covering of the Poincaré group the Lie algebra cohomology with real coefficients is trivial, so every projective unitary representation lifts to a genuine one without any extension, and the mass appears instead as the eigenvalue of a Casimir operator. This difference will have to be taken into account when the strategy of this chapter is applied to the relativistic case.
\end{remark}

% =====================================================================
% CROSS-REFERENCE MAP (labels used above -> thesis numbers; adapt to real labels)
% ch:projective            Chapter 3
% ch:classical             Chapter 1
% def:lie-group            Definition 1.2.8
% def:strongly-continuous  Definition 1.2.28
% lem:schur                Lemma 1.2.26
% def:hilbert-space        Definition 2.2.6
% def:symmetric-operator   Definition 2.2.28
% thm:stone                Theorem 2.10
% thm:von-neumann          Theorem 2.11
% def:regular-representation Definition 2.4.16
% prop:weyl-relations      Proposition 2.4.8
% eq:weyl-relations        Equations (2.44)-(2.45)
% eq:heisenberg-ccr        Equation (2.37)
% sec:weyl-algebra         Section 2.4
% def:hilbertian-manifold  Definition 3.1.1
% def:poisson-riemann-brackets Definition 3.1.13
% def:kahler-bracket       Definition 3.1.16 (circ_nu product)
% def:projective-hilbert-space Definition 3.2.2
% def:local-chart          Definition 3.2.5
% prop:local-chart         Proposition 3.2.6
% prop:holomorphic-atlas   Proposition 3.2.7
% prop:topologies-coincide Proposition 3.2.9
% eq:physical-kahler       Equation (3.24)
% def:mean-value           Definition 3.3.1
% prop:hamiltonian-schrodinger Proposition 3.3.2
% sec:hamiltonian-flow     Section 3.3
% prop:poisson-commutator  Proposition 3.3.4
% prop:kahler-isomorphisms Proposition 3.3.6
% lem:metric-ambient       Lemma 3.4.1
% prop:jordan-riemann      Proposition 3.4.3
% def:dispersion           Definition 3.4.4
% eq:dispersion-relation   Equation (3.29)
% rem:riemannian-distance  Remark 3.4.11
% thm:geometrical-born     Theorem 3.8
% sec:momentum-map         Section 3.6
% prop:quantum-momentum-map Proposition 3.6.1
% ex:u1-phase              Example 3.6.2
% NEW LABELS FROM THE CHAPTER 1 INSERT (end of Section 1.2), file chapter1_insert_cohomology.tex:
% def:cocycle, prop:multiplier-cocycle, prop:rephasing, def:cohomology-group, def:central-extension,
% prop:cocycle-extension, prop:commutator-pairing, def:lie-cohomology-group,
% def:infinitesimal-cocycle, prop:bargmann-correspondence, ex:su2-cohomology, ex:so3-su2
% ex:su-n  Example 1.2.5 (su(n)), used in the Chapter 1 insert
% prop:heisenberg-obstruction Proposition 2.4.3
% def:schrodinger-representation Definition 2.4.19
% NEW LABELS FROM THE CHAPTER 2 INSERT (end of Section 2.2), file chapter2_insert_pvm_schwartz.tex:
% def:pvm, def:pvm-integral, def:pvm-support, prop:pvm-calculus, thm:spectral,
% def:strongly-commuting, thm:joint-spectral, thm:snag, prop:pvm-commutation,
% prop:pvm-covariance, def:schwartz, prop:schwartz-properties
%
% BIBLIOGRAPHY TO ADD OR MATCH (verify every location before use)
% TODO-Bargmann1954   V. Bargmann, On unitary ray representations of continuous groups, Ann. Math. 59 (1954) 1-46
% TODO-LevyLeblond1971 J.-M. Levy-Leblond, Galilei group and Galilean invariance, in Group Theory and its Applications II (1971)
% TODO-Garding1947    L. Garding, Note on continuous representations of Lie groups, PNAS 33 (1947)
% TODO-Nelson1959     E. Nelson, Analytic vectors, Ann. Math. 70 (1959) 572-615
% TODO-RS03           Reed-Simon I (already in the thesis bibliography as [RS03])
% TODO-SNAG           Stone-Naimark-Ambrose-Godement theorem (Reed-Simon or Moretti)
% TODO-Mor18, TODO-FourierUniqueness, TODO-Sobolev  (used in the Chapter 2 insert)
% TODO-Hal03          Hall (already in the thesis bibliography as [Hal03])
\end{document}